% =====================================================================
%  CAPITULO 11 - ELETROMAGNETISMO
% =====================================================================
\section{Eletromagnetismo}

\begin{conceito}[Antes de começar]
\textbf{Campo de um fio reto longo:}
$B = \dfrac{\mu_0\,I}{2\pi\,r}$.\par
\textbf{Campo no centro de uma espira circular:}
$B = \dfrac{\mu_0\,I}{2\,R}$.\par
\textbf{Campo no interior de um solenoide:}
$B = \mu_0\,n\,I$, onde $n = N/\ell$ é o número de espiras por metro.\par
\textbf{Fluxo magnético:} $\Phi = B\,A\cos\theta$, medido em webers (Wb).\par
\textbf{Força sobre um condutor:} $F = B\,I\,\ell\sin\theta$.\par
\textbf{Força sobre uma carga em movimento:} $F = q\,v\,B\sin\theta$.\par
\textbf{Lei de Faraday--Lenz:} $\varepsilon = -N\,\dfrac{\mathrm{d}\Phi}{\mathrm{d}t}$;
o sinal negativo (Lenz) indica que a corrente induzida se opõe à variação de fluxo
que a criou. Constante: $\mu_0 = 4\pi\times10^{-7}\,\si{\tesla\meter\per\ampere}$.
\end{conceito}

\legendaniveis

% =====================================================================
\subsection{Conceitos iniciais e materiais magnéticos}
% =====================================================================

\blocofixacao

\begin{questao}[1]{}
O princípio da inseparabilidade dos polos magnéticos garante que:
\begin{alternativas}
\item existem monopolos magnéticos isolados.
\item ao dividir um ímã, obtêm-se dois novos ímãs, cada um com dois polos.
\item o polo norte pode ser isolado do polo sul.
\item as linhas de campo magnético têm início e fim.
\item um ímã pode ter apenas um polo.
\end{alternativas}
\resposta{(b)}
\resolucao{Diferentemente das cargas elétricas, não existem monopolos magnéticos:
ao partir um ímã, cada pedaço volta a ter polos norte e sul. Por isso as linhas de
campo magnético são sempre \emph{fechadas} (sem começo nem fim), o que elimina
(d).}
\end{questao}

\begin{questao}[1]{}
Quanto à permeabilidade magnética relativa $\mu_r$, um material
\textbf{ferromagnético} caracteriza-se por:
\begin{alternativas}
\item $\mu_r \approx 1$.
\item $\mu_r$ ligeiramente menor que 1.
\item $\mu_r$ muito maior que 1 (centenas a milhares).
\item $\mu_r = 0$.
\item $\mu_r$ negativa.
\end{alternativas}
\resposta{(c)}
\resolucao{Ferromagnéticos (ferro, níquel, cobalto) têm $\mu_r$ da ordem de
centenas ou milhares --- por isso concentram intensamente as linhas de campo e são
usados em núcleos de transformadores e motores. Paramagnéticos têm
$\mu_r\gtrsim1$; diamagnéticos, $\mu_r\lesssim1$.}
\end{questao}

\begin{questao}[1]{}
Dois materiais são submetidos ao mesmo campo externo $H$. Um apresenta
$B_1 = \SI{0.002}{\tesla}$ e o outro, $B_2 = \SI{1.8}{\tesla}$. Qual é o
ferromagnético?
\resposta{O de $B_2 = \SI{1.8}{\tesla}$}
\resolucao{Como $B = \mu\,H$, para o mesmo $H$ o material com maior $B$ tem maior
permeabilidade. O valor $B_2 = \SI{1.8}{\tesla}$ é muito superior, indicando
$\mu$ elevado: material \textbf{ferromagnético}. O de $B_1 = \SI{0.002}{\tesla}$ é
paramagnético ou diamagnético (resposta muito fraca ao campo).}
\end{questao}

\begin{questao}[1]{Apostila original revisada}
Dois materiais apresentam permeabilidades relativas
$\mu_{rA}=1{,}0005$ e $\mu_{rB}=0{,}99998$. Classifique cada material como
paramagnético ou diamagnético e descreva, qualitativamente, como reage a um campo
magnético externo.
\resposta{$A$: paramagnético; $B$: diamagnético}
\resolucao{Materiais com $\mu_r$ ligeiramente maior que 1 são paramagnéticos e são
fracamente atraídos pelo campo. Materiais com $\mu_r$ ligeiramente menor que 1 são
diamagnéticos e são fracamente repelidos.}
\end{questao}

\blocoaplicacao

\begin{questao}[2]{}
Explique por que os materiais ferromagnéticos são preferidos na construção de
núcleos de máquinas elétricas (transformadores, motores, geradores).
\resposta{Ver comentário}
\resolucao{Por terem $\mu_r$ muito elevado, os ferromagnéticos \emph{concentram} e
intensificam o campo magnético para uma mesma corrente de excitação, o que aumenta
o fluxo e, portanto, a eficiência do acoplamento magnético (mais tensão induzida,
mais torque). Além disso, oferecem um "caminho" de baixa relutância para o fluxo,
canalizando-o pelo núcleo em vez de dispersá-lo pelo ar. A desvantagem são as
perdas por histerese e correntes parasitas, minimizadas usando núcleos laminados
e ligas específicas (aço-silício).}
\end{questao}

\blocoaprofundamento

\begin{questao}[3]{}
Um mesmo enrolamento é usado com dois núcleos diferentes: um de ar
($\mu_r \approx 1$) e outro de ferro-silício ($\mu_r = 4000$). Mantendo a mesma
corrente:
\begin{itensq}
\item Qual é a razão entre as induções magnéticas $B$ obtidas nos dois casos?
\item Se o núcleo de ferro \emph{saturar} a partir de certo ponto, o que acontece
      com essa razão ao se aumentar muito a corrente?
\item Por que a saturação impõe um limite de projeto às máquinas elétricas?
\end{itensq}
\resposta{(a) $B_{\text{Fe}}/B_{\text{ar}} = 4000$;
(b) a razão \emph{cai}; (c) ver comentário}
\resolucao{(a) Como $B = \mu H = \mu_r\mu_0 H$, com $H$ fixo (mesma corrente e
geometria):
\[
\frac{B_{\text{Fe}}}{B_{\text{ar}}}=\frac{\mu_{r,\text{Fe}}}{\mu_{r,\text{ar}}}
=4000.
\]
O núcleo ferromagnético multiplica o campo por quatro mil --- é o que torna
viáveis transformadores e motores compactos.\\
(b) Na \textbf{saturação}, praticamente todos os domínios magnéticos já estão
alinhados; o material não consegue mais contribuir. A partir daí, aumentos de
corrente produzem acréscimos de $B$ apenas na proporção do ar
($\mu_r \to 1$ incremental). A razão $B_{\text{Fe}}/B_{\text{ar}}$, portanto,
\emph{diminui} progressivamente --- a curva $B\times H$ "achata".\\
(c) Porque o torque de um motor e a tensão de um transformador dependem do fluxo,
e o fluxo não pode crescer indefinidamente. Ultrapassado o joelho da curva de
saturação:
\begin{itemize}[leftmargin=1.6em,itemsep=1pt,topsep=2pt]
\item a corrente de magnetização aumenta acentuadamente (aquecendo o enrolamento) sem ganho
      proporcional de fluxo;
\item a forma de onda distorce, gerando harmônicos;
\item as perdas aumentam desproporcionalmente.
\end{itemize}
Por isso o projetista dimensiona a \emph{seção do núcleo} para operar
\textbf{abaixo} do joelho --- tipicamente em torno de \SI{1.5}{\tesla} para
aço-silício. \emph{Saturação é o limite físico que define o tamanho mínimo de
uma máquina elétrica.}}
\end{questao}

\begin{questao}[3]{}
Um material ferromagnético submetido a um campo alternado percorre um
\textbf{ciclo de histerese}: a curva $B \times H$ não é uma reta e não retorna
pelo mesmo caminho.

\circuitoqq{%
\begin{tikzpicture}
 \begin{axis}[width=8.2cm,height=5.6cm,xlabel={$H$ (A/m)},ylabel={$B$ (T)},
   xmin=-3,xmax=3,ymin=-1.6,ymax=1.6,xtick=\empty,ytick=\empty,
   axis lines=middle,axis line style={-Stealth}]
   \addplot[Icol,very thick,smooth] coordinates
     {(-2.6,-1.35)(-1.6,-1.25)(-0.6,-1.0)(0,-0.75)(0.5,-0.2)(1.0,0.75)(1.7,1.2)(2.6,1.35)};
   \addplot[Vcol,very thick,smooth] coordinates
     {(2.6,1.35)(1.6,1.25)(0.6,1.0)(0,0.75)(-0.5,0.2)(-1.0,-0.75)(-1.7,-1.2)(-2.6,-1.35)};
   \node[font=\footnotesize,Vcol] at (axis cs:0.42,0.95) {$B_r$};
   \node[font=\footnotesize,Icol] at (axis cs:-1.0,-0.35) {$H_c$};
 \end{axis}
\end{tikzpicture}}{Ciclo de histerese de um material ferromagnético.}

\begin{itensq}
\item Identifique o que representam $B_r$ (remanência) e $H_c$ (coercividade).
\item A \emph{área} interna do ciclo tem significado físico. Qual?
\item Compare os requisitos de um material para (i) núcleo de transformador e
      (ii) ímã permanente. Qual deve ter ciclo "largo" e qual "estreito"?
\end{itensq}
\resposta{Ver análise comentada}
\resolucao{\textbf{(a)} $B_r$ é a \textbf{remanência}: a indução que \emph{permanece}
no material quando o campo externo é removido ($H = 0$). É o que faz um prego
continuar magnetizado depois de encostar num ímã.\\
$H_c$ é a \textbf{coercividade}: o campo externo, de sentido oposto, necessário
para \emph{anular} a magnetização ($B = 0$). Quantifica a resistência do material à desmagnetização.

\textbf{(b)} A área interna do ciclo representa a \textbf{energia dissipada por
unidade de volume e por ciclo} --- as chamadas \emph{perdas por histerese}. Cada
inversão do campo exige reorientar os domínios magnéticos, e esse processo é
irreversível, gerando calor. A potência perdida é
\[
P_h = k\,f\,V\,(\text{área do ciclo}),
\]
proporcional à \emph{frequência} $f$: por isso máquinas de alta frequência exigem
materiais especiais.

\textbf{(c) Requisitos opostos.}
\begin{center}\small\sffamily
\begin{tabular}{@{}l l l@{}}
\toprule
 & \textbf{Núcleo de transformador} & \textbf{Ímã permanente}\\
\midrule
Ciclo desejado & \textbf{estreito} (fino) & \textbf{largo} (amplo)\\
$H_c$ & baixa & alta\\
$B_r$ & baixa & alta\\
Objetivo & minimizar perdas & reter magnetização\\
Material típico & ferro-silício, ferrite & alnico, neodímio (NdFeB)\\
Denominação & magneticamente \emph{mole} & magneticamente \emph{duro}\\
\bottomrule
\end{tabular}
\end{center}
\textbf{Raciocínio central:} o transformador inverte o campo 50 ou 60 vezes por
segundo --- cada inversão custa a área do ciclo em calor, logo o ciclo deve ser o
mais estreito possível. O ímã permanente, ao contrário, precisa \emph{resistir} à
desmagnetização: quanto maior a área do ciclo, maior é sua resistência à desmagnetização. \emph{O mesmo
gráfico, lido com objetivos opostos, leva a escolhas opostas de material.}}
\end{questao}

% BEGIN BANCO COMPLEMENTAR Conceitos iniciais e materiais magnéticos
% =====================================================================
%  Banco complementar --- Conceitos Iniciais e Materiais Magneticos
%  REVISADO. Substitui a versao gerada por template (10 clones em N2,
%  9 em N3 e 8 questoes conceituais de uma linha marcadas como N4).
%  O cap11 ja cobre: inseparabilidade dos polos (MC), definicao de
%  ferromagnetico por mu_r (MC), comparacao de dois materiais por B sob
%  mesmo H, classificacao para/diamagnetica a partir de mu_r, por que
%  ferromagneticos sao usados em nucleos, comparacao ar x ferro-silicio
%  no mesmo enrolamento e a descricao qualitativa do ciclo de histerese.
%  A LAMINACAO ja foi tratada quantitativamente no banco de Lenz, e o
%  ENTREFERRO no banco de fluxo -- por isso as N4 antigas sobre esses
%  dois temas foram removidas daqui.
%  Este banco explora: H, B e M como grandezas distintas, curva de
%  magnetizacao inicial e mu diferencial, temperatura de Curie, energia
%  do ciclo, materiais moles x duros com numeros, ponto de operacao em
%  nucleo nao linear, circuito magnetico com dois materiais e ensaio
%  experimental do ciclo.
%  Distribuicao mantida: 4 N1 + 11 N2 + 10 N3 + 8 N4 = 33
% =====================================================================

\subsubsection*{Banco complementar --- Conceitos iniciais e materiais magnéticos}

\begin{atencao}[Organização do banco]
Três grandezas convivem e são frequentemente confundidas. $\vec H$
($\si{\ampere\per\meter}$) é o campo \emph{aplicado}, determinado apenas pelas
correntes de condução; $\vec M$ é a \emph{magnetização} do material; e
$\vec B=\mu_0(\vec H+\vec M)$ é o que efetivamente age sobre cargas e produz
fluxo. Em materiais lineares, $\vec B=\mu_r\mu_0\vec H$ --- mas ferromagnéticos
\textbf{não} são lineares, e neles $\mu_r$ é uma função, não uma constante. O N4
trata de energia do ciclo, ponto de operação e ensaio.
\end{atencao}

\blocofixacao

\begin{questao}[1]{}
Um material linear tem $\mu_r=500$ e é submetido a
$H=\SI{100}{\ampere\per\meter}$. Determine $B$.
\dados{$\mu_0=4\pi\times10^{-7}\,\si{\tesla\meter\per\ampere}$}
\resposta{$B=\SI{62.8}{\milli\tesla}$}
\resolucao{
\[
B=\mu_r\mu_0H=(500)(4\pi\times10^{-7})(100)=\pot{6.28}{-2}\,\si{\tesla}.
\]
\textbf{Sem o material} ($\mu_r=1$), o mesmo $H$ produziria apenas
$\SI{126}{\micro\tesla}$. O material multiplicou $B$ por $500$ --- e é
exatamente esse o papel de um núcleo magnético.\\
\textbf{Atenção:} $H$ é fixado pelas correntes; $B$ é o que o material
\emph{responde}. Trocar o núcleo muda $B$, mas não muda $H$.}
\end{questao}

\begin{questao}[1]{}
Determine a magnetização $M$ do material da questão anterior.
\resposta{$M=\SI{49.9}{\kilo\ampere\per\meter}$}
\resolucao{De $B=\mu_0(H+M)$:
\[
M=\frac{B}{\mu_0}-H=\mu_rH-H=(\mu_r-1)H
=(499)(100)=\pot{4.99}{4}\,\si{\ampere\per\meter}.
\]
\textbf{Interpretação:} a magnetização do material equivale a uma corrente
superficial de $\SI{49.9}{\kilo\ampere\per\meter}$ --- quinhentas vezes maior
que a corrente de condução que a provocou.\\
\textbf{É por isso que o núcleo "amplifica":} ele não cria energia, apenas
alinha os momentos magnéticos já existentes no material, cujo efeito somado
supera em muito o do enrolamento.}
\end{questao}

\begin{questao}[1]{}
Define-se a suscetibilidade magnética como $\chi=\mu_r-1$. Classifique os
materiais com $\chi=+\pot{5}{-4}$ e $\chi=-\pot{2}{-5}$.
\resposta{Paramagnético e diamagnético, respectivamente}
\resolucao{\textbf{Paramagnético} ($\chi>0$, pequeno): os momentos atômicos
alinham-se \emph{fracamente} com o campo, reforçando-o de forma quase
imperceptível. Exemplos: alumínio, platina, oxigênio líquido.\\
\textbf{Diamagnético} ($\chi<0$): o campo induz momentos \emph{opostos}, e o
material é levemente repelido. Exemplos: cobre, água, bismuto, grafite. Todo
material tem componente diamagnética; ela só é observável quando não há
efeito maior a mascará-la.\\
\textbf{Ordens de grandeza:} $|\chi|\sim10^{-5}$ a $10^{-4}$ nesses dois casos,
contra $\chi\sim10^{3}$ em ferromagnéticos --- diferença de sete ordens de
grandeza.}
\end{questao}

\begin{questao}[1]{}
Um enrolamento de $500$ espiras percorrido por $\SI{0.4}{\ampere}$ envolve um
núcleo de caminho médio $\SI{25}{\centi\meter}$. Determine $H$.
\resposta{$H=\SI{800}{\ampere\per\meter}$}
\resolucao{Da lei de Ampère aplicada ao caminho médio,
\[
H=\frac{NI}{\ell}=\frac{500(0{,}4)}{0{,}25}=\SI{800}{\ampere\per\meter}.
\]
\textbf{Observação central:} $H$ \textbf{não depende do material}. Ele é
determinado exclusivamente pelas correntes e pela geometria.\\
Trocar o núcleo de ar por ferro não altera $H$ --- altera $B$, e portanto o
fluxo. É essa separação que torna útil distinguir as duas grandezas.}
\end{questao}

\blocoaplicacao

\begin{questao}[2]{}
Um núcleo de $\SI{10}{\centi\meter}$ de caminho médio, seção
$\SI{1}{\centi\meter\squared}$ e $\mu_r=1000$ recebe $\SI{200}{}$
ampère-espiras.
\begin{itensq}
\item Determine a relutância.
\item Determine o fluxo e a densidade de fluxo.
\item Avalie o resultado.
\end{itensq}
\resposta{$\mathcal{R}=\pot{7.96}{5}\,\si{\per\henry}$;
$\Phi=\SI{0.251}{\milli\weber}$; $B=\SI{2.51}{\tesla}$ --- \textbf{saturado}}
\resolucao{(a)
\[
\mathcal{R}=\frac{\ell}{\mu_r\mu_0A}
=\frac{0{,}10}{(1000)(4\pi\times10^{-7})(\pot{1}{-4})}
=\pot{7.96}{5}\,\si{\per\henry}.
\]
(b)
\[
\Phi=\frac{NI}{\mathcal{R}}=\frac{200}{\pot{7.96}{5}}
=\pot{2.51}{-4}\,\si{\weber};
\qquad
B=\frac{\Phi}{A}=\SI{2.51}{\tesla}.
\]
(c) \textbf{Resultado impossível.} Nenhum aço comum sustenta
$\SI{2.51}{\tesla}$ --- todos saturam entre $1{,}5$ e $\SI{2}{\tesla}$.\\
\textbf{O que realmente ocorreria:} conforme $B$ se aproxima da saturação,
$\mu_r$ desaba de $1000$ para valores próximos de $1$. A relutância dispara e o
fluxo estaciona em torno de $B_{sat}A=\SI{0.15}{\milli\weber}$.\\
\textbf{Regra de método:} calcular com $\mu_r$ constante e \emph{depois}
verificar $B$ contra $B_{sat}$. Se o valor exceder, o modelo linear não se
aplica e é preciso recorrer à curva $B\times H$ real.}
\end{questao}

\begin{questao}[2]{}
Um mesmo enrolamento produz $H=\SI{500}{\ampere\per\meter}$. Determine $B$ para
núcleo de ar, de alumínio ($\mu_r=1{,}00002$) e de ferro-silício
($\mu_r=4000$).
\resposta{$\SI{628}{\micro\tesla}$; $\SI{628}{\micro\tesla}$;
$\SI{2.51}{\tesla}$}
\resolucao{
\[
\text{ar}:\ B=\mu_0(500)=\SI{628}{\micro\tesla};
\]
\[
\text{alumínio}:\ B=(1{,}00002)\mu_0(500)=\SI{628}{\micro\tesla};
\]
\[
\text{ferro-silício}:\ B=(4000)\mu_0(500)=\SI{2.51}{\tesla}
\ \text{(saturaria antes)}.
\]
\textbf{Duas leituras.}\\
\emph{(i)} O alumínio é \emph{indistinguível} do ar para fins práticos --- sua
suscetibilidade de $\pot{2}{-5}$ altera $B$ na quinta casa decimal. É por isso
que se pode enrolar bobinas em carretéis de alumínio sem consequência magnética.\\
\emph{(ii)} O ferro-silício produziria $\SI{2.51}{\tesla}$ \emph{se fosse
linear} --- mas ele satura antes. Com $H=\SI{500}{\ampere\per\meter}$, um
ferro-silício real já está bem dentro do joelho da curva, com $B$ real em torno
de $1{,}5$ a $\SI{1.7}{\tesla}$.}
\end{questao}

\begin{questao}[2]{}
Um material apresenta $B=\SI{1.2}{\tesla}$ para $H=\SI{400}{\ampere\per\meter}$
e $B=\SI{1.5}{\tesla}$ para $H=\SI{1200}{\ampere\per\meter}$.
\begin{itensq}
\item Determine $\mu_r$ em cada ponto.
\item Interprete a variação.
\end{itensq}
\resposta{$\mu_r=2390$ e $995$}
\resolucao{(a)
\[
\mu_r=\frac{B}{\mu_0H}:\quad
\frac{1{,}2}{(4\pi\times10^{-7})(400)}=2390;
\qquad
\frac{1{,}5}{(4\pi\times10^{-7})(1200)}=995.
\]
(b) \textbf{A permeabilidade caiu para menos da metade} enquanto $H$
triplicou.\\
\textbf{Interpretação:} o material está entrando na região de \emph{joelho} da
curva de magnetização. Os domínios já estão em boa parte alinhados, e cada
acréscimo de $H$ produz cada vez menos acréscimo de $B$.\\
\textbf{Consequência prática:} triplicar a corrente aumentou o fluxo em apenas
$25\%$. Operar nessa região é ineficiente --- o cobre gasta muito para
render pouco.\\
\textbf{Lição de fundo:} $\mu_r$ de catálogo é sempre um valor \emph{para uma
condição específica}. Usá-lo fora dessa condição é um erro comum e silencioso.}
\end{questao}

\begin{questao}[2]{}
Explique a diferença entre a permeabilidade \textbf{normal}
($\mu=B/H$) e a permeabilidade \textbf{diferencial}
($\mu_d=\mathrm{d}B/\mathrm{d}H$), e indique quando cada uma se aplica.
\resposta{Normal para o ponto de operação; diferencial para pequenos sinais
sobrepostos}
\resolucao{\textbf{Permeabilidade normal:} razão entre $B$ e $H$ no ponto de
operação. Ela governa o \emph{fluxo total} e, portanto, a relutância e o
dimensionamento do núcleo.\\
\textbf{Permeabilidade diferencial:} inclinação da curva naquele ponto. Ela
governa a resposta a \emph{pequenas variações} sobrepostas --- e portanto a
indutância incremental.\\
\textbf{Onde diferem drasticamente:} perto da saturação, $B/H$ ainda é grande
(centenas), mas $\mathrm{d}B/\mathrm{d}H$ já caiu para quase $\mu_0$. Um
indutor com corrente contínua elevada pode ter indutância incremental dezenas
de vezes menor que a nominal.\\
\textbf{Consequência de projeto:} em filtros com componente contínua --- muito
comuns em fontes chaveadas --- é a permeabilidade \emph{diferencial} que
determina o desempenho, e é ela que precisa ser consultada no catálogo.}
\end{questao}

\begin{questao}[2]{}
Compare a temperatura de Curie do ferro ($\SI{770}{\celsius}$), do níquel
($\SI{358}{\celsius}$) e de uma ferrite típica
($\approx\SI{250}{\celsius}$), e explique o que ocorre acima dela.
\resposta{Acima de $T_C$ o material torna-se paramagnético e perde
praticamente toda a permeabilidade}
\resolucao{\textbf{O que ocorre acima de $T_C$.} A agitação térmica supera a
interação de troca que mantém os momentos atômicos alinhados dentro dos
domínios. A estrutura de domínios se desfaz, e o material passa de
\textbf{ferromagnético} a \textbf{paramagnético}:
\[
\mu_r:\ 10^{3}\text{--}10^{4}\ \longrightarrow\ \approx1.
\]
\textbf{A transição é abrupta} --- trata-se de uma transição de fase de segunda
ordem, e $\mu_r$ despenca em poucos graus em torno de $T_C$.\\
\textbf{Consequências práticas.}
\begin{itemize}
\item \textbf{Ferrites são as mais vulneráveis.} Com $T_C$ próxima de
$\SI{250}{\celsius}$, e temperatura de operação já em torno de
$\SI{100}{\celsius}$, a margem é pequena. Um transformador de fonte chaveada
que superaqueça perde indutância, a corrente dispara e ele se destrói --- fuga
térmica.
\item \textbf{Aplicação deliberada:} panelas de indução com fundo cuja $T_C$ é
próxima da temperatura desejada param de aquecer sozinhas ao atingi-la. É um
termostato sem sensor.
\item \textbf{Desmagnetização por aquecimento:} aquecer um ímã acima de
$T_C$ e resfriá-lo sem campo o deixa completamente desmagnetizado.
\end{itemize}}
\end{questao}

\begin{questao}[2]{}
Um ímã permanente de ferrite tem coeficiente de temperatura de
$\SI{-0.2}{\percent\per\celsius}$ para a remanência. Determine a variação de
$B_r$ entre $\SI{20}{\celsius}$ e $\SI{120}{\celsius}$.
\resposta{Redução de $20\%$}
\resolucao{
\[
\frac{\Delta B_r}{B_r}=(-0{,}002)(100)=-0{,}20=-20\%.
\]
\textbf{Reversível ou não?} Até certo limite, a perda é \textbf{reversível} --- o
ímã recupera a magnetização ao esfriar. Acima de uma temperatura crítica
(bem abaixo de $T_C$), parte da perda torna-se \textbf{irreversível}, porque
domínios se reorientam permanentemente.\\
\textbf{Comparação entre materiais:}
\begin{center}
\begin{tabular}{lcc}
\hline
Material & Coef. de $B_r$ & $T_{\max}$ típica\\
\hline
Ferrite & $\SI{-0.2}{\percent\per\celsius}$ & $\SI{250}{\celsius}$\\
Neodímio (NdFeB) & $\SI{-0.12}{\percent\per\celsius}$ & $\SI{80}{}$ a
 $\SI{200}{\celsius}$\\
Samário-cobalto & $\SI{-0.03}{\percent\per\celsius}$ & $\SI{300}{\celsius}$\\
\hline
\end{tabular}
\end{center}
\textbf{Compromisso revelador:} o neodímio tem a maior densidade de energia mas
péssima estabilidade térmica; o samário-cobalto é bem mais estável e mais caro.
Em motores que aquecem, a escolha raramente é o material mais forte.}
\end{questao}

\begin{questao}[2]{}
Um núcleo de aço-silício dissipa perdas por histerese proporcionais à área do
ciclo. Com área de $\SI{200}{\joule\per\meter\cubed}$ por ciclo e volume de
$\SI{1000}{\centi\meter\cubed}$, determine a potência a $\SI{60}{\hertz}$ e a
$\SI{400}{\hertz}$.
\resposta{$\SI{12}{\watt}$ e $\SI{80}{\watt}$}
\resolucao{
\[
P_h=f\times(\text{área})\times V:
\quad
60(200)(\pot{1}{-3})=\SI{12}{\watt};
\qquad
400(200)(\pot{1}{-3})=\SI{80}{\watt}.
\]
\textbf{Proporcionalidade linear com $f$} --- cada ciclo custa a mesma energia,
e mais ciclos por segundo significam mais potência.\\
\textbf{Contraste com as correntes parasitas}, cujas perdas vão com
$f^{2}$. Em baixa frequência a histerese domina; em alta, as parasitas.\\
\textbf{Consequência:} o mesmo núcleo que funciona bem em $\SI{60}{\hertz}$
pode ser inviável em $\SI{400}{\hertz}$ --- e é por isso que equipamentos
aeronáuticos, que operam nessa frequência, usam ligas especiais de lâminas
finíssimas.}
\end{questao}

\begin{questao}[2]{}
Um material magnético \textbf{mole} tem $H_c=\SI{50}{\ampere\per\meter}$ e um
material \textbf{duro} tem $H_c=\SI{500}{\kilo\ampere\per\meter}$. Compare-os
e indique aplicações.
\resposta{Razão de $10^{4}$ na coercividade; moles para núcleos, duros para ímãs
permanentes}
\resolucao{\textbf{Coercividade $H_c$} é o campo necessário para \emph{anular} a
magnetização remanente. Ela mede o quanto o material "resiste" a mudar de
estado.
\begin{center}
\begin{tabular}{lccl}
\hline
Tipo & $H_c$ & Área do ciclo & Aplicação\\
\hline
Mole & baixa & pequena & núcleos de transformador e motor\\
Duro & alta & grande & ímãs permanentes\\
\hline
\end{tabular}
\end{center}
\textbf{Materiais moles} (ferro-silício, permalloy, ferrites moles):
magnetizam e desmagnetizam facilmente, com pouca perda por ciclo. É o que se
quer em algo que inverte a magnetização $120$ vezes por segundo.\\
\textbf{Materiais duros} (alnico, ferrite dura, NdFeB): resistem à
desmagnetização, mantendo o campo sem corrente alguma.\\
\textbf{O critério é a área do ciclo:} pequena para quem vai percorrê-lo muitas
vezes, grande para quem precisa permanecer magnetizado.}
\end{questao}

\begin{questao}[2]{}
Explique por que a curva de magnetização \textbf{inicial} de um material
virgem difere do ciclo de histerese subsequente.
\resposta{A curva inicial parte da origem; o ciclo nunca mais passa por ela}
\resolucao{\textbf{Material virgem:} os domínios estão orientados
aleatoriamente, e a magnetização líquida é nula. Aplicando $H$ crescente, a
curva parte da \textbf{origem} e sobe até a saturação --- é a \emph{curva de
magnetização inicial}.\\
\textbf{Ao reduzir $H$ a zero}, a magnetização \emph{não} retorna a zero: resta
a \textbf{remanência} $B_r$. Os domínios que se reorientaram não voltam
espontaneamente.\\
\textbf{A partir daí}, ciclando $H$ entre $\pm H_{\max}$, o material percorre
um laço fechado que \textbf{nunca mais passa pela origem}.\\
\textbf{Como voltar ao estado virgem:} aplicar campo alternado de amplitude
decrescente até zero (desmagnetização por ciclos decrescentes), ou aquecer acima
de $T_C$ e resfriar sem campo.\\
\textbf{Importância prática:} ao levantar experimentalmente a curva
$B\times H$, é preciso desmagnetizar o material antes --- caso contrário, os
primeiros pontos refletem a história anterior, não a propriedade do material.}
\end{questao}

\begin{questao}[2]{}
Um transformador é desligado no pico da corrente de magnetização. Explique o
que ocorre na religação e por que a corrente de \emph{inrush} pode ser elevada.
\resposta{O fluxo remanente soma-se ao fluxo imposto, podendo levar o núcleo a
saturação profunda}
\resolucao{\textbf{Ao desligar}, o núcleo retém fluxo remanente
$\Phi_r$ --- consequência direta da histerese.\\
\textbf{Ao religar}, o fluxo evolui a partir desse valor. No pior caso, a tensão
é aplicada no instante em que exigiria fluxo de mesma polaridade, e o fluxo
total pode chegar a
\[
\Phi\approx\Phi_r+2\Phi_{\max},
\]
muito além da saturação.\\
\textbf{Consequência:} saturado, o núcleo perde permeabilidade, a indutância
desaba e a corrente é limitada apenas pela resistência do enrolamento. A
corrente de \emph{inrush} pode atingir dez a vinte vezes a nominal, por alguns
ciclos.\\
\textbf{Soluções empregadas:}
\begin{itemize}
\item \textbf{Chaveamento sincronizado:} energizar no instante correto da onda
de tensão.
\item \textbf{Resistor ou NTC de pré-carga}, curto-circuitado após alguns
ciclos.
\item \textbf{Proteção com curva temporizada}, que tolera o transitório sem
atuar.
\end{itemize}
\textbf{Note a conexão:} este fenômeno só existe \emph{porque} há histerese ---
um núcleo hipotético sem remanência não teria \emph{inrush} por essa causa.}
\end{questao}

\begin{questao}[2]{}
Um circuito magnético tem dois trechos em série: ferro
($\ell_1=\SI{30}{\centi\meter}$, $\mu_r=2000$) e ar
($\ell_2=\SI{2}{\milli\meter}$), ambos com seção
$\SI{5}{\centi\meter\squared}$. Determine a relutância total e a fração de cada
trecho.
\resposta{$\mathcal{R}_{fe}=\pot{2.39}{5}$, $\mathcal{R}_{ar}=\pot{3.18}{6}$;
o ar responde por $93\%$}
\resolucao{
\[
\mathcal{R}_{fe}=\frac{0{,}30}{(2000)\mu_0(\pot{5}{-4})}
=\pot{2.39}{5}\,\si{\per\henry};
\]
\[
\mathcal{R}_{ar}=\frac{\pot{2}{-3}}{\mu_0(\pot{5}{-4})}
=\pot{3.18}{6}\,\si{\per\henry}.
\]
\[
\mathcal{R}_{tot}=\pot{3.42}{6};
\qquad
\frac{\mathcal{R}_{ar}}{\mathcal{R}_{tot}}=93\%.
\]
\textbf{Dois milímetros de ar dominam trinta centímetros de ferro} --- porque a
razão de permeabilidades ($2000$) supera de longe a razão de comprimentos
($150$).\\
\textbf{As relutâncias somam-se em série}, exatamente como resistências. E o
fluxo é o mesmo nos dois trechos, o que garante $B$ igual --- já que a seção
também é a mesma.\\
\textbf{Consequência:} o comportamento do circuito é governado quase
inteiramente pelo entreferro, e a não linearidade do ferro fica mascarada. É
esse efeito que \emph{lineariza} indutores com entreferro.}
\end{questao}

\blocoaprofundamento

\begin{questao}[3]{}
Distinga rigorosamente $\vec H$, $\vec B$ e $\vec M$.
\begin{itensq}
\item Escreva a relação entre elas e as unidades.
\item Determine qual delas muda ao inserir um núcleo em um solenoide, mantida a
      corrente.
\item Explique qual é "mais fundamental" e por quê.
\end{itensq}
\resposta{$\vec B=\mu_0(\vec H+\vec M)$; ao inserir o núcleo, $H$ não muda e
$B$ aumenta}
\resolucao{(a)
\[
\vec B=\mu_0\left(\vec H+\vec M\right),
\]
com $[H]=[M]=\si{\ampere\per\meter}$ e $[B]=\si{\tesla}$.\\
\emph{($H$ é às vezes chamado "intensidade de campo magnético" e $B$
"densidade de fluxo" --- nomenclatura infeliz, porque sugere que $H$ é o campo
principal.)}\\
(b) \textbf{Ao inserir o núcleo com corrente constante:}
\begin{itemize}
\item $H=NI/\ell$ \textbf{não muda} --- ele depende só das correntes de
condução.
\item $M$ \textbf{passa de zero a $(\mu_r-1)H$} --- o material se magnetiza.
\item $B$ \textbf{aumenta por} $\mu_r$ --- consequência das duas anteriores.
\end{itemize}
(c) \textbf{Qual é mais fundamental: $\vec B$.} Argumentos:
\begin{itemize}
\item É $\vec B$ que aparece na força de Lorentz
($\vec F=q\vec v\times\vec B$) --- é ele que age sobre cargas.
\item É $\vec B$ que aparece na lei de Faraday --- é seu fluxo que induz f.e.m.
\item $\nabla\cdot\vec B=0$ é lei universal; $\nabla\cdot\vec H=-\nabla\cdot
\vec M$ não é nula em geral.
\end{itemize}
\textbf{Para que serve $\vec H$, então:} ele é conveniente porque sua circulação
depende \emph{apenas} das correntes de condução ($\oint\vec H\cdot
\mathrm{d}\vec\ell=NI$), sem precisar contabilizar as correntes de
magnetização do material. É uma grandeza \emph{auxiliar} de enorme utilidade
prática --- mas auxiliar.}
\end{questao}

\begin{questao}[3]{}
A curva $B\times H$ de um material apresenta três regiões distintas.
\begin{itensq}
\item Descreva cada uma em termos de domínios.
\item Identifique onde $\mu_r$ é máxima.
\item Indique a região de operação recomendada.
\end{itensq}
\resposta{Deslocamento reversível, deslocamento irreversível e rotação;
$\mu_r$ máxima no trecho intermediário}
\resolucao{(a) \textbf{Três regiões.}
\begin{enumerate}
\item \textbf{Campos baixos --- deslocamento reversível de paredes.} As paredes
entre domínios deslocam-se ligeiramente, favorecendo os domínios alinhados com
$H$. Removido o campo, elas retornam. A curva é pouco inclinada.
\item \textbf{Campos médios --- deslocamento irreversível (saltos de
Barkhausen).} As paredes superam obstáculos (impurezas, contornos de grão) em
saltos abruptos. É aqui que a curva é mais \textbf{íngreme} e onde nasce a
histerese, pois os saltos não se desfazem ao reduzir $H$.
\item \textbf{Campos altos --- rotação dos domínios.} Já não há paredes a
mover; os momentos giram progressivamente até se alinharem com $H$. A curva
achata-se: é o \emph{joelho} e depois a saturação.
\end{enumerate}
(b) \textbf{$\mu_r$ é máxima na região intermediária}, onde
$\mathrm{d}B/\mathrm{d}H$ é maior. Nem no início (paredes presas) nem no fim
(domínios já alinhados).\\
\textbf{Consequência contraintuitiva:} a permeabilidade em campos muito baixos é
\emph{menor} que a máxima --- daí catálogos especificarem $\mu_i$ (inicial) e
$\mu_{\max}$ separadamente.\\
(c) \textbf{Região recomendada:} abaixo do joelho, tipicamente com
$B_{\max}$ entre $60$ e $80\%$ de $B_{sat}$. Isso aproveita a alta
permeabilidade sem risco de saturação em transitórios.}
\end{questao}

\begin{questao}[3]{}
Um núcleo tem curva $B\times H$ não linear, e é preciso determinar o ponto de
operação com entreferro.
\begin{itensq}
\item Formule o problema.
\item Descreva o método gráfico da reta de entreferro.
\item Explique por que o entreferro estabiliza o ponto de operação.
\end{itensq}
\resposta{Interseção da curva $B\times H$ do material com a reta imposta pelo
entreferro}
\resolucao{(a) \textbf{Formulação.} A lei de Ampère no circuito magnético dá
\[
NI=H_{fe}\,\ell_{fe}+H_g\,g,
\]
e no entreferro $H_g=B/\mu_0$ (com $B$ igual ao do ferro, por conservação do
fluxo e seção constante). Então
\[
NI=H_{fe}\ell_{fe}+\frac{B}{\mu_0}g.
\]
Duas incógnitas ($B$ e $H_{fe}$) e uma equação --- a segunda relação é a
\emph{curva do material}, que não tem forma algébrica simples.\\
(b) \textbf{Método gráfico.} Reescrevendo,
\[
B=\frac{\mu_0\,NI}{g}-\frac{\mu_0\ell_{fe}}{g}\,H_{fe},
\]
que é uma \textbf{reta} no plano $B\times H$ --- a \emph{reta de entreferro} ou
\emph{reta de carga magnética}.\\
Traçando-a sobre a curva $B\times H$ do material, o \textbf{ponto de operação} é
a interseção.\\
\emph{(É exatamente o mesmo procedimento da reta de carga usada com diodos: uma
característica não linear do dispositivo e uma reta imposta pelo circuito.)}\\
(c) \textbf{Por que o entreferro estabiliza.} A inclinação da reta é
$-\mu_0\ell_{fe}/g$: quanto \emph{maior} o entreferro, mais \textbf{deitada} a
reta.\\
Uma reta deitada cruza a curva do material em uma região de inclinação
comparável --- e ali a interseção é \emph{pouco sensível} a variações de
$\mu_r$ do material. Trocar o lote de ferro, ou aquecê-lo, quase não desloca o
ponto.\\
\textbf{Sem entreferro}, a reta é quase vertical e a interseção depende
inteiramente da curva do material --- que varia com temperatura, lote e
história magnética. \textbf{O entreferro compra previsibilidade ao custo de mais
ampère-espiras.}}
\end{questao}

\begin{questao}[3]{}
Compare aço-silício e ferrite para aplicação em transformadores.
\begin{itensq}
\item Compare $B_{sat}$, resistividade e faixa de frequência.
\item Determine em qual faixa cada um é preferível.
\item Justifique fisicamente.
\end{itensq}
\resposta{Aço-silício: alto $B_{sat}$, baixa resistividade, até
$\si{\kilo\hertz}$; ferrite: baixo $B_{sat}$, altíssima resistividade,
até $\si{\mega\hertz}$}
\resolucao{(a)
\begin{center}
\begin{tabular}{lcc}
\hline
Propriedade & Aço-silício & Ferrite (MnZn)\\
\hline
$B_{sat}$ & $1{,}5$--$\SI{2}{\tesla}$ & $0{,}3$--$\SI{0.5}{\tesla}$\\
Resistividade & $\approx\pot{5}{-7}\,\si{\ohm\meter}$ &
 $10^{-1}$ a $\SI{10}{\ohm\meter}$\\
$\mu_r$ típica & $2000$--$8000$ & $1000$--$5000$\\
Faixa útil & até $\approx\SI{1}{\kilo\hertz}$ & $\SI{20}{\kilo\hertz}$ a
 $\si{\mega\hertz}$\\
\hline
\end{tabular}
\end{center}
\textbf{Note a resistividade:} a ferrite é um \emph{cerâmico}, com
resistividade até $10^{7}$ vezes maior que a do aço.\\
(b) \textbf{Baixa frequência} ($\SI{50}{}$--$\SI{400}{\hertz}$): aço-silício.
Seu $B_{sat}$ elevado permite núcleos menores para a mesma potência, e as perdas
parasitas são controladas por laminação.\\
\textbf{Alta frequência} ($>\SI{20}{\kilo\hertz}$): ferrite. Nenhuma laminação
seria fina o bastante para conter as correntes parasitas no aço.\\
(c) \textbf{Justificativa física.} As perdas por correntes parasitas vão com
\[
P_F\propto\frac{B^{2}f^{2}t^{2}}{\rho}.
\]
Em alta frequência, o termo $f^{2}$ explode. Há duas formas de compensar:
reduzir $t$ (laminação) ou aumentar $\rho$.\\
A ferrite ataca $\rho$ --- e ganha por muitas ordens de grandeza, o que nenhuma
laminação alcança.\\
\textbf{O preço:} $B_{sat}$ baixo obriga a usar seção maior ou aceitar menos
fluxo. \textbf{Mas em alta frequência isso não importa}, porque
$V=4{,}44Nf\Phi_{\max}$ --- com $f$ mil vezes maior, o fluxo necessário é mil
vezes menor. \textbf{É a alta frequência que torna a ferrite viável, e a ferrite
que torna a alta frequência viável.}}
\end{questao}

\begin{questao}[3]{}
Um ímã permanente é retirado de seu circuito magnético e deixado com os polos
livres ao ar.
\begin{itensq}
\item Explique o que ocorre com seu ponto de operação.
\item Determine o efeito de aproximar uma peça de ferro.
\item Justifique a prática de armazenar ímãs com "quebra-fluxo".
\end{itensq}
\resposta{Ao ar, o ponto de operação desce na curva de desmagnetização; o ferro
o eleva}
\resolucao{(a) \textbf{Ao ar livre.} O ímã precisa "empurrar" o fluxo por um
caminho de alta relutância --- o próprio ar em torno dele. Isso equivale a um
entreferro enorme, e a reta de carga fica muito deitada.\\
O ponto de operação desce na \emph{curva de desmagnetização} (segundo quadrante
do laço), afastando-se de $B_r$. O ímã fica submetido a campo desmagnetizante
produzido por ele mesmo.\\
(b) \textbf{Aproximando ferro.} O ferro oferece caminho de baixa relutância, e a
reta de carga se levanta. O ponto de operação sobe, aproximando-se de
$B_r$ --- o ímã fica "mais forte".\\
É por isso que um ímã em ferradura fechado por uma armadura de ferro sustenta
muito mais peso que aberto.\\
(c) \textbf{Quebra-fluxo (\emph{keeper}).} É uma peça de ferro doce colocada
entre os polos durante o armazenamento. Ela:
\begin{itemize}
\item mantém o ponto de operação alto, evitando desmagnetização progressiva;
\item protege contra campos externos adversos;
\item reduz a atração de limalha e detritos.
\end{itemize}
\textbf{Ressalva histórica:} a prática era essencial com alnico, cuja curva de
desmagnetização é frágil. Com neodímio moderno --- cuja curva é praticamente
reta no segundo quadrante --- o efeito é bem menor, embora o quebra-fluxo
continue útil por outras razões.}
\end{questao}

\begin{questao}[3]{}
Um material tem $B_r=\SI{1.2}{\tesla}$ e $H_c=\SI{900}{\kilo\ampere\per\meter}$.
\begin{itensq}
\item Estime o produto de energia máximo.
\item Compare com ferrite ($B_r=\SI{0.4}{\tesla}$,
      $H_c=\SI{250}{\kilo\ampere\per\meter}$).
\item Explique o significado do produto de energia.
\end{itensq}
\resposta{$(BH)_{\max}\approx\SI{270}{\kilo\joule\per\meter\cubed}$ contra
$\SI{25}{\kilo\joule\per\meter\cubed}$}
\resolucao{(a) Para curva de desmagnetização aproximadamente \emph{reta} (caso
do neodímio), o produto máximo ocorre no ponto médio:
\[
(BH)_{\max}\approx\frac{B_rH_c}{4}
=\frac{(1{,}2)(\pot{9}{5})}{4}
=\SI{270}{\kilo\joule\per\meter\cubed}.
\]
(b) Para a ferrite:
\[
(BH)_{\max}\approx\frac{(0{,}4)(\pot{2.5}{5})}{4}
=\SI{25}{\kilo\joule\per\meter\cubed},
\]
cerca de \textbf{onze vezes menor}.\\
(c) \textbf{Significado.} O produto de energia mede a energia magnética
disponível \emph{por unidade de volume} do ímã. Ele é a figura de mérito que
determina quão pequeno o ímã pode ser para produzir dado campo em dado
entreferro.\\
\textbf{Consequência prática:} para o mesmo desempenho, um ímã de neodímio ocupa
cerca de um décimo do volume de um de ferrite. Isso é o que viabilizou
servomotores compactos, fones de ouvido pequenos e discos rígidos.\\
\textbf{Contrapartida:} o neodímio é caro, oxida (precisa de revestimento),
desmagnetiza acima de $\SI{80}{}$--$\SI{200}{\celsius}$ e depende de
terras-raras. A ferrite é barata, quimicamente estável e térmicamente mais
tolerante --- e por isso continua dominando em volume de produção.}
\end{questao}

\begin{questao}[3]{}
Explique a diferença entre desmagnetização por \textbf{aquecimento} e por
\textbf{campo alternado decrescente}.
\begin{itensq}
\item Descreva cada processo.
\item Compare eficácia e aplicabilidade.
\end{itensq}
\resposta{Aquecimento acima de $T_C$ desordena os domínios; campo alternado os
distribui simetricamente}
\resolucao{(a) \textbf{Aquecimento acima de $T_C$.} A agitação térmica destrói o
alinhamento dentro dos domínios. Resfriando \emph{sem campo aplicado}, os
domínios se reformam em orientações aleatórias, e a magnetização líquida é
nula.\\
\textbf{Campo alternado decrescente.} Aplica-se campo alternado de amplitude
inicialmente alta (o bastante para saturar) e reduz-se lentamente a zero. O
material percorre laços de histerese cada vez menores, espiralando até a origem.
Ao final, tantos domínios apontam em um sentido quanto no outro.\\
(b) \textbf{Comparação.}
\begin{center}
\begin{tabular}{lcc}
\hline
& Aquecimento & Campo alternado\\
\hline
Eficácia & completa & muito boa\\
Danifica a peça? & possivelmente & não\\
Aplicável a peça montada? & raramente & sim\\
Rapidez & lento & segundos\\
\hline
\end{tabular}
\end{center}
\textbf{Na prática industrial, usa-se campo alternado} --- em desmagnetizadores
de ferramentas, de peças usinadas e de fitas magnéticas. O aquecimento é
inviável para peças que não toleram $\SI{770}{\celsius}$.\\
\textbf{Cuidado essencial:} a redução do campo deve ser \emph{lenta} em relação
ao período. Desligar bruscamente o desmagnetizador pode deixar a peça
magnetizada no valor correspondente à última meia-onda --- pior que antes.\\
\textbf{Alternativa equivalente:} afastar lentamente a peça de um campo
alternado fixo, o que produz o mesmo decaimento gradual.}
\end{questao}

\begin{questao}[3]{}
Um núcleo opera com $B_{\max}=\SI{1.2}{\tesla}$ a $\SI{60}{\hertz}$. Analise o
efeito de elevar $B_{\max}$ para $\SI{1.5}{\tesla}$.
\begin{itensq}
\item Determine o efeito sobre o número de espiras.
\item Determine o efeito sobre as perdas.
\item Avalie o compromisso.
\end{itensq}
\resposta{$N$ cai $20\%$; perdas por histerese sobem cerca de $60\%$ e
parasitas $56\%$}
\resolucao{(a) De $V=4{,}44Nf\Phi_{\max}=4{,}44NfB_{\max}A$, com $V$, $f$ e
$A$ fixos:
\[
N\propto\frac{1}{B_{\max}}
\;\Longrightarrow\;
\frac{N'}{N}=\frac{1{,}2}{1{,}5}=0{,}80.
\]
\textbf{Vinte por cento menos espiras} --- menos cobre, menos resistência, menos
perdas no enrolamento.\\
(b) \textbf{Perdas no núcleo aumentam.}
\begin{itemize}
\item \textbf{Parasitas:} $P_F\propto B_{\max}^{2}$, logo
$(1{,}5/1{,}2)^{2}=1{,}56$ --- aumento de $56\%$.
\item \textbf{Histerese:} $P_h\propto B_{\max}^{n}$ com $n\approx1{,}6$ a
$2$ (expoente de Steinmetz), dando aumento de $60$ a $56\%$.
\end{itemize}
(c) \textbf{O compromisso.}
\begin{center}
\begin{tabular}{lc}
\hline
Elevar $B_{\max}$ & Efeito\\
\hline
Espiras & $-20\%$\\
Perdas no cobre & diminuem\\
Perdas no núcleo & $+\approx58\%$\\
Volume e custo & diminuem\\
Margem para saturação & \textbf{diminui}\\
\hline
\end{tabular}
\end{center}
\textbf{Existe um ótimo econômico}, em que a soma das perdas é mínima --- e ele
depende do custo relativo do cobre e do aço e do regime de carga.\\
\textbf{Mas há um limite duro:} $\SI{1.5}{\tesla}$ já está próximo do joelho. A
margem para sobretensões e para o transitório de energização encolhe, e o risco
de \emph{inrush} severo cresce. \textbf{Projetar perto da saturação economiza
material e compra problemas operacionais.}}
\end{questao}

\begin{questao}[3]{}
Explique por que a permeabilidade de um material ferromagnético não é uma
constante, e quais parâmetros a afetam.
\resposta{$\mu$ depende de $H$, da temperatura, da frequência e da história
magnética}
\resolucao{\textbf{Quatro dependências, todas relevantes.}
\begin{enumerate}
\item \textbf{Da intensidade $H$.} É a mais óbvia: $\mu_r$ parte de um valor
inicial, cresce até um máximo e despenca na saturação. Variação de uma ordem de
grandeza ou mais.
\item \textbf{Da temperatura.} $\mu_r$ tende a \emph{crescer} ao se aproximar de
$T_C$ (efeito Hopkinson) e depois colapsa abruptamente. Isso torna o
comportamento térmico não monotônico e difícil de prever.
\item \textbf{Da frequência.} Em alta frequência, as paredes de domínio não
conseguem acompanhar o campo --- há \emph{relaxação}. A permeabilidade efetiva
cai, e aparece uma componente de perda. Catálogos de ferrite trazem
$\mu$ complexo, com parte real e imaginária.
\item \textbf{Da história magnética.} Por causa da histerese, o mesmo $H$ pode
corresponder a diferentes $B$, conforme o material esteja subindo ou descendo o
laço.
\end{enumerate}
\textbf{Consequência para o projetista.} Um valor único de $\mu_r$ em catálogo é
sempre uma \emph{simplificação} referida a condições específicas --- geralmente
campo baixo, temperatura ambiente e baixa frequência.\\
\textbf{Prática correta:} usar $\mu_r$ para estimativas iniciais, e recorrer às
\emph{curvas} do fabricante (de $B\times H$, de perdas por frequência, de
$\mu$ por temperatura) para o dimensionamento final. Projetos de fontes
chaveadas que ignoram isso costumam funcionar na bancada e falhar em campo.}
\end{questao}

\begin{questao}[3]{}
Um sensor de proximidade indutivo detecta metais pela variação de indutância.
\begin{itensq}
\item Explique o princípio para metais ferromagnéticos.
\item Explique para metais não ferromagnéticos (alumínio, cobre).
\item Determine por que o alcance difere entre eles.
\end{itensq}
\resposta{Ferromagnéticos aumentam $L$ por permeabilidade; não ferromagnéticos
reduzem $L$ e o fator de qualidade por correntes parasitas}
\resolucao{(a) \textbf{Ferromagnéticos.} A aproximação do material oferece
caminho de menor relutância, aumentando a indutância da bobina sensora. A
variação de $L$ é detectada pelo circuito.\\
Simultaneamente, correntes parasitas no aço amortecem o oscilador.\\
(b) \textbf{Não ferromagnéticos.} Alumínio e cobre têm $\mu_r\approx1$ --- não
alteram a relutância. Mas são \emph{bons condutores}, e o campo alternado induz
neles correntes parasitas intensas.\\
Pela lei de Lenz, essas correntes criam campo oposto, o que \textbf{reduz} a
indutância efetiva e, sobretudo, \textbf{drena energia} do oscilador, reduzindo
seu fator de qualidade.\\
(c) \textbf{Por que o alcance difere.} Sensores indutivos comuns são
especificados para aço, e apresentam \emph{fatores de correção}:
\begin{center}
\begin{tabular}{lc}
\hline
Material & Fator de alcance\\
\hline
Aço & $1{,}00$\\
Aço inoxidável & $\approx0{,}70$\\
Latão & $\approx0{,}50$\\
Alumínio & $\approx0{,}40$\\
Cobre & $\approx0{,}30$\\
\hline
\end{tabular}
\end{center}
\textbf{Razão:} no aço, os dois efeitos (permeabilidade e parasitas) somam-se.
Nos não ferromagnéticos, resta apenas o segundo --- e ele é tanto mais fraco
quanto \emph{melhor} o condutor, porque as correntes se concentram numa
película superficial cada vez mais fina.\\
\textbf{Consequência de aplicação:} um sensor ajustado para detectar peças de
aço a $\SI{10}{\milli\meter}$ só detectará cobre a $\SI{3}{\milli\meter}$ ---
erro de projeto frequente em linhas de montagem com peças mistas.}
\end{questao}

\blocodesafio

\begin{questao}[4]{}
\textbf{(Energia do ciclo de histerese.)} Demonstre que a área do laço
$B\times H$ representa energia dissipada por unidade de volume e por ciclo.
\begin{itensq}
\item Deduza a expressão.
\item Determine a potência dissipada em função da frequência.
\item Explique a origem microscópica da perda.
\end{itensq}
\resposta{$w=\oint H\,\mathrm{d}B$ por ciclo; $P=fVw$}
\resolucao{(a) \textbf{Dedução.} A potência instantânea fornecida ao núcleo pela
fonte é $p=\varepsilon i$. Com $\varepsilon=NA\,\mathrm{d}B/\mathrm{d}t$ e
$i=H\ell/N$:
\[
p=NA\frac{\mathrm{d}B}{\mathrm{d}t}\cdot\frac{H\ell}{N}
=(A\ell)\,H\,\frac{\mathrm{d}B}{\mathrm{d}t}
=V\,H\,\frac{\mathrm{d}B}{\mathrm{d}t}.
\]
A energia por ciclo é
\[
W=\int p\,\mathrm{d}t=V\oint H\,\mathrm{d}B,
\]
e $\oint H\,\mathrm{d}B$ é precisamente a \textbf{área do laço} no plano
$B\times H$, com unidade $\si{\joule\per\meter\cubed}$.\\
\textbf{Se o material fosse sem histerese}, o laço colapsaria em uma curva
única, a integral de ida cancelaria a de volta e a energia por ciclo seria
\textbf{zero} --- toda a energia magnética seria devolvida.\\
(b)
\[
P_h=f\,V\oint H\,\mathrm{d}B,
\]
\textbf{linear em $f$}. (As perdas por correntes parasitas, distintas destas,
vão com $f^{2}$.)\\
(c) \textbf{Origem microscópica.} A perda vem dos \textbf{saltos de Barkhausen}:
as paredes de domínio ficam presas em impurezas, defeitos e contornos de grão, e
se libertam abruptamente quando o campo cresce o suficiente.\\
Cada salto é um processo \textbf{irreversível} --- a energia elástica acumulada
converte-se em vibração da rede, ou seja, em \textbf{calor}.\\
\textbf{Consequência de engenharia:} materiais mais puros e com grãos maiores e
melhor orientados têm menos obstáculos, laços mais estreitos e menores perdas.
É exatamente isso que o processo de \emph{grão orientado} do aço-silício
persegue.}
\end{questao}

\begin{questao}[4]{}
\textbf{(Circuito magnético com dois materiais.)} Um circuito tem trecho de
ferro ($\ell_1$, $\mu_1$) em série com trecho de ferrite ($\ell_2$,
$\mu_2$), mesma seção.
\begin{itensq}
\item Formule as equações.
\item Determine $B$ e $H$ em cada trecho.
\item Explique qual grandeza é contínua na interface e por quê.
\end{itensq}
\resposta{$B$ é contínuo; $H$ é descontínuo, na razão inversa das
permeabilidades}
\resolucao{(a) \textbf{Duas equações.}\\
\emph{Conservação do fluxo} (equivalente à LKC):
\[
\Phi_1=\Phi_2
\;\Longrightarrow\;
B_1A=B_2A
\;\Longrightarrow\;
B_1=B_2=B.
\]
\emph{Lei de Ampère} (equivalente à LKT):
\[
NI=H_1\ell_1+H_2\ell_2.
\]
(b) Com $H_i=B/\mu_i$:
\[
NI=B\left(\frac{\ell_1}{\mu_1}+\frac{\ell_2}{\mu_2}\right)
\;\Longrightarrow\;
B=\frac{NI}{\dfrac{\ell_1}{\mu_1}+\dfrac{\ell_2}{\mu_2}}.
\]
\[
H_1=\frac{B}{\mu_1};
\qquad
H_2=\frac{B}{\mu_2};
\qquad
\frac{H_1}{H_2}=\frac{\mu_2}{\mu_1}.
\]
(c) \textbf{Qual grandeza é contínua --- e por quê.}
\begin{itemize}
\item \textbf{$B$ é contínuo} através da interface (na componente normal),
porque $\nabla\cdot\vec B=0$: o fluxo que chega precisa continuar. Não há
"acúmulo de fluxo" possível.
\item \textbf{$H$ é descontínuo}, na razão inversa das permeabilidades. O
material de menor $\mu$ exige \emph{muito mais} $H$ para sustentar o mesmo
$B$.
\end{itemize}
\textbf{Exemplo numérico revelador.} Com $\mu_{r1}=2000$ (ferro) e
$\mu_{r2}=2$ (um material quase não magnético), $H_2$ seria \emph{mil vezes}
$H_1$ --- e praticamente toda a força magnetomotriz seria consumida naquele
trecho, ainda que ele fosse curtíssimo.\\
\textbf{É exatamente o que acontece no entreferro} ($\mu_r=1$), e por isso ele
domina o circuito. Um entreferro não é um caso especial: é apenas o caso extremo
deste problema de dois materiais.}
\end{questao}

\begin{questao}[4]{}
\textbf{(Ensaio do ciclo de histerese.)} Projete um experimento para levantar o
laço $B\times H$ de um material em osciloscópio.
\begin{itensq}
\item Descreva o circuito e o que cada canal mede.
\item Determine as constantes de conversão.
\item Aponte os cuidados experimentais.
\end{itensq}
\resposta{Canal horizontal $\propto H$ (corrente por shunt); vertical
$\propto B$ (f.e.m. integrada)}
\resolucao{(a) \textbf{Circuito.} Um núcleo toroidal com dois enrolamentos:
primário $N_1$ excitado por fonte alternada e secundário $N_2$ para medição.
\begin{itemize}
\item \textbf{Canal horizontal ($X$):} tensão sobre um resistor \emph{shunt}
$R_s$ em série com o primário. Como $H=N_1i_1/\ell$, essa tensão é proporcional
a $H$.
\item \textbf{Canal vertical ($Y$):} a f.e.m. do secundário passa por um
\textbf{integrador} $RC$. Como
$\varepsilon_2=N_2A\,\mathrm{d}B/\mathrm{d}t$, sua integral é proporcional a
$B$.
\end{itemize}
Em modo $XY$, o osciloscópio traça diretamente o laço.\\
(b) \textbf{Constantes de conversão.}
\[
H=\frac{N_1}{\ell R_s}\,v_X;
\qquad
B=\frac{R_iC_i}{N_2A}\,v_Y,
\]
com $R_iC_i$ os componentes do integrador.\\
\textbf{Condição do integrador:} $R_iC_i\gg1/f$, para que ele integre de fato em
vez de apenas atenuar. Com $f=\SI{60}{\hertz}$, adote
$R_iC_i\ge\SI{0.1}{\second}$.\\
(c) \textbf{Cuidados experimentais.}
\begin{itemize}
\item \textbf{Desmagnetizar antes} --- caso contrário os primeiros ciclos
refletem a história anterior. Comece com amplitude alta e reduza.
\item \textbf{$R_s$ pequeno} (fração de ohm), para não distorcer a corrente nem
dissipar demais. Use quatro fios se possível.
\item \textbf{Núcleo toroidal}, que garante $H$ uniforme e caminho médio bem
definido --- em núcleos em E o caminho é ambíguo.
\item \textbf{Deriva do integrador:} qualquer \emph{offset} faz o laço "subir"
lentamente na tela. Use acoplamento CA ou integrador com realimentação CC.
\item \textbf{Varrer a amplitude} para obter a família de laços e localizar a
curva de magnetização normal, que é o lugar geométrico dos vértices.
\end{itemize}
\textbf{O que se extrai:} $B_r$ (interseção com o eixo vertical), $H_c$
(interseção com o horizontal), $B_{sat}$ e a área --- e daí as perdas por
ciclo.}
\end{questao}

\begin{questao}[4]{}
\textbf{(Domínios magnéticos.)} Explique a estrutura de domínios e sua relação
com as propriedades macroscópicas.
\begin{itensq}
\item Explique por que os domínios existem.
\item Relacione a estrutura com $B_r$, $H_c$ e $\mu_r$.
\item Explique a diferença entre materiais moles e duros nesse nível.
\end{itensq}
\resposta{Domínios minimizam a energia magnetostática; a mobilidade das paredes
determina $\mu_r$ e $H_c$}
\resolucao{(a) \textbf{Por que existem.} Um cristal ferromagnético
uniformemente magnetizado criaria campo externo intenso, com alta energia
magnetostática armazenada no espaço.\\
Dividir-se em \textbf{domínios} de orientações opostas faz esse campo externo
quase desaparecer --- as linhas se fecham internamente. O custo é a energia das
\emph{paredes de domínio} (paredes de Bloch), onde os momentos giram
gradualmente.\\
\textbf{O equilíbrio} entre esses dois termos fixa o tamanho típico dos
domínios, tipicamente de micrômetros a dezenas de micrômetros.\\
(b) \textbf{Relação com as propriedades macroscópicas.}
\begin{itemize}
\item \textbf{$\mu_r$ alta} exige paredes \emph{móveis} --- material puro, sem
tensões internas, com grãos grandes e bem orientados.
\item \textbf{$H_c$ baixa} tem a mesma origem: paredes que deslizam facilmente
não resistem à inversão.
\item \textbf{$B_r$ alta} exige que os domínios permaneçam alinhados após a
retirada do campo --- o que requer \emph{anisotropia} cristalina forte,
"ancorando" a magnetização em direções preferenciais.
\end{itemize}
(c) \textbf{Moles contra duros, no nível microscópico.}
\begin{center}
\begin{tabular}{lll}
\hline
& Moles & Duros\\
\hline
Paredes & muito móveis & fixadas ou inexistentes\\
Impurezas & minimizadas & introduzidas de propósito\\
Anisotropia & baixa & alta\\
Grãos & grandes, orientados & finos, monodomínio\\
\hline
\end{tabular}
\end{center}
\textbf{O contraste é elegante:} ímãs modernos de neodímio são feitos de
partículas tão pequenas que cada uma é um \textbf{monodomínio} --- não há
paredes para mover, e inverter a magnetização exige girar todos os momentos
contra a anisotropia. Daí a coercividade enorme.\\
\textbf{Já o aço-silício de grão orientado} persegue o oposto: grãos grandes,
alinhados por laminação a frio, com paredes livres para deslizar.\\
\textbf{Mesma física, objetivos opostos --- e processos de fabricação
inteiramente diferentes.}}
\end{questao}

\begin{questao}[4]{}
\textbf{(Projeto sob restrição de perdas.)} Um transformador de
$\SI{60}{\hertz}$ deve ter perdas no núcleo inferiores a $\SI{2}{\watt}$, com
material cuja perda específica é
$\SI{1.5}{\watt\per\kilogram}$ a $\SI{1.5}{\tesla}$ e
$\SI{0.9}{\watt\per\kilogram}$ a $\SI{1.2}{\tesla}$.
\begin{itensq}
\item Determine a massa admissível em cada condição.
\item Determine o efeito sobre o número de espiras.
\item Estabeleça o critério de escolha.
\end{itensq}
\resposta{$\SI{1.33}{\kilogram}$ a $\SI{1.5}{\tesla}$;
$\SI{2.22}{\kilogram}$ a $\SI{1.2}{\tesla}$}
\resolucao{(a)
\[
m\le\frac{2}{1{,}5}=\SI{1.33}{\kilogram}\ \text{a } \SI{1.5}{\tesla};
\qquad
m\le\frac{2}{0{,}9}=\SI{2.22}{\kilogram}\ \text{a } \SI{1.2}{\tesla}.
\]
(b) De $V=4{,}44NfB_{\max}A$, com tensão e frequência fixas,
\[
N\,B_{\max}A=\text{const}.
\]
Reduzir $B_{\max}$ de $1{,}5$ para $\SI{1.2}{\tesla}$ exige aumentar
$NA$ em $25\%$ --- mais espiras, ou seção maior, ou ambos.\\
(c) \textbf{Critério de escolha.} As duas opções levam a projetos diferentes:
\begin{center}
\begin{tabular}{lcc}
\hline
& $B=\SI{1.5}{\tesla}$ & $B=\SI{1.2}{\tesla}$\\
\hline
Massa de núcleo & menor & maior ($+67\%$)\\
Espiras & menos & mais ($+25\%$)\\
Perdas no cobre & menores & maiores\\
Margem para saturação & pequena & confortável\\
Corrente de \emph{inrush} & severa & moderada\\
\hline
\end{tabular}
\end{center}
\textbf{Note o conflito:} operar em $B$ alto economiza \emph{ferro} mas gasta
mais \emph{cobre} por unidade de perda; operar em $B$ baixo faz o inverso.\\
\textbf{Decisão:}
\begin{itemize}
\item \textbf{Se o custo de material dominar} (transformadores de baixo custo,
uso intermitente), escolha $B$ alto.
\item \textbf{Se o custo de energia dominar} (operação contínua por décadas),
escolha $B$ baixo --- as perdas em vazio existem $\SI{8760}{\hour}$ por ano.
\item \textbf{Se houver risco de sobretensão}, escolha $B$ baixo por segurança:
um transformador operando a $\SI{1.5}{\tesla}$ satura com apenas $10\%$ de
sobretensão.
\end{itemize}}
\end{questao}

\begin{questao}[4]{}
\textbf{(Anisotropia e grão orientado.)} Explique o que é o aço-silício de grão
orientado e por que ele é usado em transformadores mas não em motores.
\begin{itensq}
\item Explique o processo e o efeito.
\item Compare as perdas nas duas direções.
\item Justifique a diferença de aplicação.
\end{itensq}
\resposta{Grãos alinhados numa direção preferencial; excelente ao longo dela,
ruim transversalmente}
\resolucao{(a) \textbf{Processo e efeito.} O aço-silício é laminado a frio e
recozido de modo que os grãos cristalinos se alinhem com o eixo de fácil
magnetização do ferro (a direção $\langle100\rangle$ da célula cúbica) paralelo
à direção de laminação.\\
\textbf{Resultado:} nessa direção, magnetizar é muito mais fácil --- as paredes
de domínio se movem livremente, $\mu_r$ é altíssima e as perdas caem
drasticamente.\\
(b) \textbf{Comparação por direção.}
\begin{center}
\begin{tabular}{lcc}
\hline
& Direção de laminação & Transversal\\
\hline
Perdas específicas & $\approx\SI{1}{\watt\per\kilogram}$ &
 $2$ a $3\times$ maiores\\
$\mu_r$ & muito alta & bem menor\\
$B$ para dado $H$ & alto & reduzido\\
\hline
\end{tabular}
\end{center}
\textbf{A anisotropia é grande} --- o material é excelente em uma direção e
medíocre nas outras.\\
(c) \textbf{Por que serve a transformadores e não a motores.}
\begin{itemize}
\item \textbf{Transformador:} o fluxo percorre um caminho \emph{fixo e
conhecido} --- as colunas e as culatras do núcleo. É possível cortar as lâminas
de modo que o fluxo siga sempre a direção favorável. \emph{(Nos cantos, onde o
fluxo precisa dobrar, usam-se juntas em $\ang{45}$ para minimizar a
transição.)}
\item \textbf{Motor:} o fluxo no estator \textbf{gira}, percorrendo todas as
direções ao longo de cada revolução. Não existe direção preferencial a
privilegiar --- metade do tempo o fluxo estaria na direção ruim.
\end{itemize}
\textbf{Por isso motores usam aço de grão \emph{não} orientado}, deliberadamente
\textbf{isotrópico}: pior que o grão orientado em sua melhor direção, mas
uniforme em todas --- e portanto melhor \emph{na média} sobre uma rotação
completa.\\
\textbf{Princípio geral:} otimizar para o caso médio quando a solicitação é
variável, e para o caso específico quando ela é fixa.}
\end{questao}

\begin{questao}[4]{}
\textbf{(Blindagem magnética com material de alta permeabilidade.)} Analise
quantitativamente a blindagem de campo estático por invólucro de mu-metal.
\begin{itensq}
\item Explique o mecanismo por relutância.
\item Estime o fator de blindagem para invólucro cilíndrico.
\item Aponte as limitações práticas.
\end{itensq}
\resposta{Fator $\approx\mu_r t/D$; o mecanismo é desvio de fluxo, não
cancelamento}
\resolucao{(a) \textbf{Mecanismo.} Ao contrário da blindagem elétrica --- que
\emph{cancela} o campo por redistribuição de cargas --- a blindagem magnética
\textbf{desvia} o fluxo.\\
O invólucro de alta $\mu_r$ oferece caminho de relutância muito menor que o ar
interno. O fluxo "prefere" circular pela parede e contorna a região protegida.\\
\textbf{Analogia:} é um divisor de fluxo, em que o material é o ramo de baixa
relutância e a cavidade, o de alta.\\
(b) \textbf{Fator de blindagem.} Para um cilindro de diâmetro $D$ e parede
$t$, com $t\ll D$, a estimativa clássica é
\[
S\approx\frac{\mu_r\,t}{D}.
\]
Com $\mu_r=20000$ (mu-metal recozido), $t=\SI{1}{\milli\meter}$ e
$D=\SI{100}{\milli\meter}$:
\[
S\approx\frac{20000(1)}{100}=200,
\]
ou seja, atenuação de cerca de $\SI{46}{\decibel}$.\\
\textbf{Compare com a blindagem elétrica}, que atinge muitas ordens de grandeza
com uma folha fina de qualquer metal. A magnética é \emph{muito} mais modesta.\\
(c) \textbf{Limitações práticas.}
\begin{itemize}
\item \textbf{Saturação.} Se o campo externo for intenso, a parede satura,
$\mu_r$ desaba e a blindagem falha justamente quando mais se precisa dela.
\emph{Solução:} camadas múltiplas, com a externa de material de alto
$B_{sat}$ (aço comum) e a interna de mu-metal.
\item \textbf{Sensibilidade mecânica.} O mu-metal perde permeabilidade se for
dobrado, cortado ou submetido a choque. Exige \textbf{recozimento após a
conformação}, em atmosfera de hidrogênio --- processo caro.
\item \textbf{Aberturas.} Furos e fendas quebram o caminho de baixa relutância e
degradam a blindagem muito mais que no caso elétrico.
\item \textbf{Custo e peso}, ambos elevados.
\end{itemize}
\textbf{Conclusão:} blindar campo magnético estático é caro, frágil e de eficácia
limitada. Sempre que possível, é preferível \emph{afastar} a fonte, reduzir a
área das espiras ou compensar ativamente o campo com bobinas.}
\end{questao}

\begin{questao}[4]{}
\textbf{(Ponto de operação de ímã permanente.)} Determine o ponto de operação
de um ímã em um circuito magnético com entreferro.
\begin{itensq}
\item Formule a equação da reta de carga.
\item Determine a condição de máximo aproveitamento.
\item Explique o efeito de aumentar o entreferro.
\end{itensq}
\resposta{A reta de carga tem inclinação
$-\dfrac{\ell_m A_g}{g A_m}$; o aproveitamento é máximo em
$(BH)_{\max}$}
\resolucao{(a) \textbf{Reta de carga.} Duas condições:
\emph{conservação do fluxo} ($B_mA_m=B_gA_g$) e \emph{lei de Ampère}
($H_m\ell_m+H_gg=0$, pois não há corrente).\\
Da segunda, $H_m\ell_m=-H_gg=-\dfrac{B_g}{\mu_0}g$. Combinando com a primeira:
\[
B_m=-\mu_0\,\frac{\ell_m A_g}{g\,A_m}\,H_m.
\]
É uma reta pela \textbf{origem} no segundo quadrante ($H_m<0$, $B_m>0$) ---
a chamada \emph{linha de carga} do ímã.\\
\textbf{Note:} o ímã opera com $H$ \emph{negativo} --- ele está sujeito ao
próprio campo desmagnetizante. Sem entreferro ($g\to0$), a reta seria vertical e
o ponto de operação estaria em $B_r$.\\
(b) \textbf{Máximo aproveitamento.} A energia disponível no entreferro é
proporcional ao produto $B_mH_m$. O ponto ótimo é aquele em que
$|B_mH_m|$ é \textbf{máximo} --- o já mencionado $(BH)_{\max}$.\\
Dimensiona-se o circuito \emph{para que a reta de carga passe por esse ponto},
escolhendo a razão $\dfrac{\ell_m A_g}{g A_m}$ adequada. Isso minimiza o volume
de ímã necessário.\\
(c) \textbf{Efeito de aumentar o entreferro.} A inclinação
$-\mu_0\ell_mA_g/(gA_m)$ diminui em módulo --- a reta \textbf{deita}. O ponto de
operação desce na curva de desmagnetização:
\begin{itemize}
\item $B_m$ diminui (menos fluxo);
\item $|H_m|$ aumenta (mais campo desmagnetizante sobre o próprio ímã);
\item se o entreferro for grande demais, o ponto pode ultrapassar o
\textbf{joelho} da curva de desmagnetização --- e aí ocorre perda
\textbf{irreversível} de magnetização.
\end{itemize}
\textbf{Consequência prática decisiva:} desmontar um circuito magnético com ímã
de alnico pode danificá-lo permanentemente, pois o ponto de operação despenca.
Ímãs de neodímio, cuja curva de desmagnetização é praticamente reta até o
segundo quadrante inteiro, toleram isso muito melhor --- e é uma das razões de
sua popularidade, além do produto de energia.}
\end{questao}

% END BANCO COMPLEMENTAR

\gabarito[12.1 Conceitos iniciais e materiais magnéticos]

% =====================================================================
\subsection{Campo magnético}
% =====================================================================

\blocofixacao

\begin{questao}[1]{}
O campo magnético a uma distância $r$ de um fio retilíneo longo percorrido por
corrente $I$:
\begin{alternativas}
\item é proporcional a $r$.
\item é inversamente proporcional a $r$.
\item é inversamente proporcional a $r^2$.
\item independe de $r$.
\item é proporcional a $r^2$.
\end{alternativas}
\resposta{(b)}
\resolucao{Para um fio reto longo, $B = \dfrac{\mu_0 I}{2\pi r}$: o campo é
\emph{inversamente} proporcional à distância $r$ (não ao quadrado, como no campo
elétrico de uma carga puntiforme).}
\end{questao}

\begin{questao}[1]{}
Um fio reto longo conduz $I = \SI{10}{\ampere}$. Determine o módulo do campo
magnético a $\SI{5}{\centi\meter}$ do fio.
\dados{$\mu_0 = 4\pi\times10^{-7}\,\si{\tesla\meter\per\ampere}$.}
\resposta{$B = \pot{4}{-5}\,\si{\tesla} = \SI{40}{\micro\tesla}$}
\resolucao{$B = \dfrac{\mu_0 I}{2\pi r}
= \dfrac{(\pot{4\pi}{-7})(10)}{2\pi(0{,}05)}
= \dfrac{\pot{2}{-6}}{0{,}05} = \pot{4}{-5}\,\si{\tesla}$.\\
Um atalho útil: $\dfrac{\mu_0}{2\pi} = \pot{2}{-7}$, então
$B = \pot{2}{-7}\cdot\dfrac{I}{r}$.}
\end{questao}

\begin{questao}[1]{}
Um fio retilíneo é percorrido por corrente para cima (saindo do chão em direção ao
teto). Usando a regra da mão direita, as linhas de campo magnético ao redor do fio:
\begin{alternativas}
\item apontam radialmente para fora do fio.
\item apontam radialmente para dentro do fio.
\item formam círculos concêntricos em torno do fio.
\item são paralelas ao fio.
\item não existem (fio reto não gera campo).
\end{alternativas}
\resposta{(c)}
\resolucao{O campo de um fio reto forma \emph{círculos concêntricos} em planos
perpendiculares ao fio. Pela regra da mão direita (polegar no sentido da corrente,
dedos envolvendo o fio), determina-se o sentido dessas linhas. O campo nunca é
radial (isso é característico do campo \emph{elétrico} de uma carga).}
\end{questao}

\blocoaplicacao

\begin{questao}[2]{}
Um solenoide de $\SI{1000}{}$ espiras por metro é percorrido por
$I = \SI{2}{\ampere}$. Determine o campo magnético no seu interior.
\dados{$\mu_0 = 4\pi\times10^{-7}\,\si{\tesla\meter\per\ampere}$.}
\resposta{$B \approx \SI{2.51}{\milli\tesla}$}
\resolucao{$B = \mu_0\,n\,I = (\pot{4\pi}{-7})(1000)(2)
= \pot{8\pi}{-4} \approx\SI{2.51}{\milli\tesla}$. O campo no interior de um
solenoide é \emph{uniforme} e independe do raio.}
\end{questao}

\begin{questao}[2]{}
Uma espira circular de raio $R = \SI{10}{\centi\meter}$ conduz
$I = \SI{5}{\ampere}$. Determine o campo magnético no seu centro.
\dados{$\mu_0 = 4\pi\times10^{-7}\,\si{\tesla\meter\per\ampere}$.}
\resposta{$B \approx \SI{31.4}{\micro\tesla}$}
\resolucao{$B = \dfrac{\mu_0 I}{2R}
= \dfrac{(\pot{4\pi}{-7})(5)}{2(0{,}1)}
= \dfrac{\pot{2\pi}{-6}}{0{,}2}=\pot{\pi}{-5}
\approx\SI{31.4}{\micro\tesla}$.}
\end{questao}

\begin{questao}[2]{}
Um fio reto conduz $I = \SI{4}{\ampere}$. Determine o campo magnético a
$\SI{2}{\centi\meter}$ do fio.
\dados{$\mu_0 = 4\pi\times10^{-7}\,\si{\tesla\meter\per\ampere}$.}
\resposta{$B = \SI{40}{\micro\tesla}$}
\resolucao{$B = \dfrac{\mu_0 I}{2\pi r} = \pot{2}{-7}\cdot\dfrac{4}{0{,}02}
= \pot{4}{-5}\,\si{\tesla} = \SI{40}{\micro\tesla}$
(usando o atalho $\mu_0/2\pi = \pot{2}{-7}$).}
\end{questao}

\blocoaprofundamento

\begin{questao}[3]{}
Dois fios paralelos, distantes $\SI{10}{\centi\meter}$, conduzem correntes de
$\SI{3}{\ampere}$ e $\SI{6}{\ampere}$ no \textbf{mesmo sentido}. Determine o campo
magnético resultante no ponto médio entre eles e diga se os fios se atraem ou se
repelem.
\dados{$\mu_0 = 4\pi\times10^{-7}\,\si{\tesla\meter\per\ampere}$.}
\resposta{$B \approx \SI{12}{\micro\tesla}$; os fios se atraem}
\resolucao{No ponto médio ($r = \SI{5}{\centi\meter}$ de cada fio), os campos dos
dois fios têm \emph{sentidos opostos} (correntes no mesmo sentido), então os
módulos se subtraem:
\[
B = \frac{\mu_0}{2\pi r}(I_2 - I_1)
= \pot{2}{-7}\cdot\frac{6-3}{0{,}05}
= \pot{2}{-7}\cdot60 = \pot{1.2}{-5}\,\si{\tesla}\approx\SI{12}{\micro\tesla}.
\]
Quanto à força \emph{entre} os fios: correntes de mesmo sentido se
\textbf{atraem} (regra fundamental do eletromagnetismo, base da definição do
ampère). Note a distinção: o \emph{campo} no ponto médio se subtrai, mas a
\emph{força} entre os fios é atrativa --- são perguntas diferentes.}
\end{questao}

\begin{questao}[3]{}
Dois fios paralelos e longos, separados por $\SI{5}{\centi\meter}$, conduzem
$\SI{10}{\ampere}$ cada, em sentidos \textbf{opostos}. Determine a força por
unidade de comprimento entre eles e diga se é de atração ou repulsão.
\dados{$\mu_0 = 4\pi\times10^{-7}\,\si{\tesla\meter\per\ampere}$.}
\resposta{$F/\ell = \SI{4e-4}{\newton\per\meter}$, repulsão}
\resolucao{A força por comprimento entre dois fios paralelos é
\[
\frac{F}{\ell}=\frac{\mu_0 I_1 I_2}{2\pi d}
=\pot{2}{-7}\cdot\frac{10\cdot10}{0{,}05}
=\pot{4}{-4}\,\si{\newton\per\meter}.
\]
Correntes em sentidos \emph{opostos} se \textbf{repelem} (e as de mesmo sentido se
atraem). Esta força é a base da definição do ampère no SI.}
\end{questao}

\begin{questao}[3]{}
Um \textbf{toroide} de raio médio $r = \SI{10}{\centi\meter}$ possui $N = 200$
espiras percorridas por $I = \SI{3}{\ampere}$.
\begin{itensq}
\item Aplicando a lei de Ampère a uma circunferência de raio $r$ interna ao
      toroide, deduza a expressão de $B$ e calcule seu valor.
\item Mostre que o campo \emph{fora} do toroide é nulo.
\item Compare com um solenoide reto de mesmo $N$ e comprimento igual ao
      perímetro médio do toroide. O que se conclui?
\end{itensq}
\dados{$\mu_0 = 4\pi\times10^{-7}\,\si{\tesla\meter\per\ampere}$.}
\resposta{(a) $B = \dfrac{\mu_0 N I}{2\pi r} = \SI{1.2}{\milli\tesla}$;
(b) e (c) ver comentário}
\resolucao{\textbf{(a)} A lei de Ampère estabelece
$\oint \vec B\cdot\mathrm{d}\vec\ell = \mu_0 I_{\text{env}}$. Escolhendo uma
circunferência de raio $r$ concêntrica ao toroide, por simetria $B$ é constante e
tangente ao longo dela:
\[
B\,(2\pi r) = \mu_0\,N I
\;\Longrightarrow\;
B=\frac{\mu_0 N I}{2\pi r}
=\frac{(\pot{4\pi}{-7})(200)(3)}{2\pi(0{,}1)}
=\pot{1.2}{-3}\,\si{\tesla}=\SI{1.2}{\milli\tesla}.
\]
\textbf{(b) Campo externo nulo.} Para um percurso \emph{fora} do toroide há dois
casos: (i) circunferência de raio maior que o externo --- a corrente entra e sai
do enrolamento, de modo que a corrente líquida envolvida é \emph{zero};
(ii) circunferência de raio menor que o interno --- não envolve corrente alguma.
Em ambos, $\oint \vec B\cdot\mathrm{d}\vec\ell = 0$ e, pela simetria,
$B = 0$.\\
\textbf{Consequência prática:} o toroide \textbf{confina} o campo em seu interior,
praticamente sem dispersão. Por isso transformadores toroidais produzem muito
menos interferência eletromagnética que os de núcleo E-I, e são preferidos em
áudio e instrumentação de precisão.

\textbf{(c) Comparação com o solenoide.} O perímetro médio é
$\ell = 2\pi r = \SI{0.628}{\meter}$. Um solenoide reto com o mesmo $N$ e esse
comprimento teria
\[
B = \mu_0\,\frac{N}{\ell}\,I = \frac{\mu_0 N I}{2\pi r},
\]
\emph{exatamente a mesma expressão}. Ou seja: o toroide é um solenoide
"fechado sobre si mesmo". A diferença essencial é que o solenoide reto tem
\textbf{extremidades}, onde as linhas de campo escapam (efeito de borda) e o
campo é menos uniforme; o toroide não tem bordas, sendo o dispositivo ideal para
estudar o comportamento magnético puro de um material --- razão pela qual as
curvas de histerese são levantadas em amostras toroidais.}
\end{questao}

\begin{questao}[3]{Apostila original revisada}
Duas espiras circulares, concêntricas e coplanares, têm raios
$R_1=10\pi\,\si{\centi\metre}$ e $R_2=2{,}5\pi\,\si{\centi\metre}$. A primeira é
percorrida por $I_1=\SI{2}{\ampere}$. Determine o valor e o sentido de $I_2$ para
que o campo magnético resultante no centro seja nulo.
\resposta{$I_2=\SI{0.5}{\ampere}$, em sentido oposto a $I_1$}
\resolucao{No centro de uma espira, $B=\mu_0 I/(2R)$. Para cancelamento,
$I_1/R_1=I_2/R_2$. Assim,
$I_2=I_1(R_2/R_1)=2(2{,}5/10)=\SI{0.5}{\ampere}$. Os sentidos devem ser opostos.}
\end{questao}

% BEGIN BANCO COMPLEMENTAR Campo magnético
% =====================================================================
%  Banco complementar --- Campo Magnetico  (REVISADO)
%  Substitui a versao gerada por template (5 clones em N1, 9 em N2,
%  8 em N3 e 8 questoes conceituais marcadas como N4).
%  O cap11 ja cobre: B proporcional a 1/r (MC), fio de 10 A a 5 cm,
%  regra da mao direita, solenoide de 1000 esp/m, ESPIRA CIRCULAR de
%  R=10 cm no centro, fio de 4 A a 2 cm, DOIS FIOS PARALELOS (mesmo
%  sentido e sentidos opostos), TOROIDE e duas espiras concentricas.
%  Este banco evita esses casos e explora: fio de secao finita por
%  lei de Ampere, cabo coaxial, campo no eixo de espira, dipolo
%  magnetico a grande distancia, projeto de solenoide, mapeamento com
%  sensor Hall, blindagem magnetica e os limites praticos de eletroimas.
%  Distribuicao mantida: 5 N1 + 9 N2 + 8 N3 + 8 N4 = 30
% =====================================================================

\subsubsection*{Banco complementar --- Campo magnético}

\begin{atencao}[Organização do banco]
Três resultados cobrem quase tudo: fio retilíneo longo
($B=\mu_0I/2\pi r$), centro de espira ($B=\mu_0I/2R$) e solenoide longo
($B=\mu_0nI$, \emph{independente} do raio). Note que apenas o primeiro decai com
a distância --- dentro do solenoide o campo é uniforme. O N3 trata de condutores
de seção finita, cabo coaxial e campo fora do eixo; o N4 exige dedução por
Biot--Savart, projeto, blindagem e limites físicos.
\end{atencao}

\blocofixacao

\begin{questao}[1]{}
Um fio retilíneo longo conduz $\SI{2}{\ampere}$. Determine o campo magnético a
$\SI{1}{\centi\meter}$ dele.
\dados{$\mu_0=4\pi\times10^{-7}\,\si{\tesla\meter\per\ampere}$}
\resposta{$B=\SI{40}{\micro\tesla}$}
\resolucao{
\[
B=\frac{\mu_0 I}{2\pi r}=\frac{(4\pi\times10^{-7})(2)}{2\pi(0{,}01)}
=\frac{\pot{2}{-7}\times2}{0{,}01}=\pot{4}{-5}\,\si{\tesla}.
\]
\textbf{Atalho útil:} $\dfrac{\mu_0}{2\pi}=\pot{2}{-7}$ exatamente, de modo que
$B=\pot{2}{-7}\dfrac{I}{r}$ --- evita carregar $\pi$ pela conta toda.\\
\textbf{Ordem de grandeza:} $\SI{40}{\micro\tesla}$ é comparável ao campo
magnético da Terra ($\approx\SI{50}{\micro\tesla}$). Uma bússola perto desse fio
já seria perturbada.}
\end{questao}

\begin{questao}[1]{}
Um solenoide tem $\SI{2000}{}$ espiras por metro e conduz
$\SI{3}{\ampere}$. Determine o campo em seu interior.
\resposta{$B=\SI{7.54}{\milli\tesla}$}
\resolucao{
\[
B=\mu_0 nI=(4\pi\times10^{-7})(2000)(3)=\pot{7.54}{-3}\,\si{\tesla}.
\]
\textbf{Note o que \emph{não} aparece na fórmula:} o raio do solenoide e o
número total de espiras. Só importa a \textbf{densidade linear}
$n=N/L$.\\
Dois solenoides, um fino e um grosso, com a mesma densidade de espiras e a mesma
corrente, produzem o \emph{mesmo} campo interno.}
\end{questao}

\begin{questao}[1]{}
Determine a corrente necessária para que um fio retilíneo produza
$B=\SI{100}{\micro\tesla}$ a $\SI{2}{\centi\meter}$.
\resposta{$I=\SI{10}{\ampere}$}
\resolucao{
\[
I=\frac{2\pi rB}{\mu_0}=\frac{rB}{\pot{2}{-7}}
=\frac{(0{,}02)(\pot{1}{-4})}{\pot{2}{-7}}=\SI{10}{\ampere}.
\]
\textbf{Leitura prática:} são necessários $\SI{10}{\ampere}$ para produzir, a
dois centímetros, apenas o dobro do campo terrestre. Campos magnéticos
intensos exigem correntes altas --- ou muitas espiras.}
\end{questao}

\begin{questao}[1]{}
Sobre o campo magnético de um solenoide longo, é correto afirmar:
\begin{alternativas}[2]
\item é uniforme no interior e praticamente nulo no exterior
\item é maior no exterior
\item aumenta com a distância ao eixo, no interior
\item é nulo no interior
\end{alternativas}
\resposta{(a)}
\resolucao{Aplicando a lei de Ampère a um retângulo com um lado dentro e outro
fora, obtém-se $B_{int}=\mu_0nI$ e $B_{ext}\approx0$ para um solenoide
\emph{ideal} (infinitamente longo).\\
\textbf{Por que o campo externo quase se anula:} as linhas que saem de uma
extremidade retornam pela outra, espalhando-se por um volume enorme --- a
densidade de linhas fora é desprezível diante da interna.\\
\textbf{Em solenoides reais}, o campo externo não é exatamente nulo, e nas
extremidades o campo interno cai para cerca de \emph{metade} do valor central.}
\end{questao}

\begin{questao}[1]{}
O campo magnético da Terra vale aproximadamente $\SI{50}{\micro\tesla}$.
Expresse-o em gauss.
\dados{$\SI{1}{\tesla}=10^{4}\,\si{\gauss}$}
\resposta{$\SI{0.5}{\gauss}$}
\resolucao{
\[
B=\pot{5}{-5}\,\si{\tesla}\times10^{4}=\SI{0.5}{\gauss}.
\]
\textbf{Escala de referência para o tópico:}
\begin{center}
\begin{tabular}{lcc}
\hline
Situação & Tesla & Gauss\\
\hline
Campo terrestre & $\pot{5}{-5}$ & $0{,}5$\\
Ímã de geladeira & $\pot{5}{-3}$ & $50$\\
Ímã de neodímio (superfície) & $0{,}5$ & $5000$\\
Ressonância magnética & $1{,}5$ a $3$ & $\pot{1.5}{4}$ a $\pot{3}{4}$\\
\hline
\end{tabular}
\end{center}
O tesla é uma unidade \emph{grande}: campos de laboratório raramente passam de
alguns tesla, e o gauss continua em uso justamente por ser mais próximo dos
valores cotidianos.}
\end{questao}

\blocoaplicacao

\begin{questao}[2]{}
Um solenoide de $100$ espiras e $\SI{10}{\centi\meter}$ de comprimento conduz
$\SI{1}{\ampere}$, com núcleo de ar.
\begin{itensq}
\item Determine a densidade de espiras.
\item Determine o campo interno.
\end{itensq}
\resposta{$n=\SI{1000}{\per\meter}$; $B=\SI{1.26}{\milli\tesla}$}
\resolucao{(a) $n=\dfrac{100}{0{,}10}=\SI{1000}{\per\meter}$.\\
(b) $B=\mu_0nI=(4\pi\times10^{-7})(1000)(1)
=\pot{1.257}{-3}\,\si{\tesla}$.\\
\textbf{Comparação instrutiva:} isso é apenas $25$ vezes o campo terrestre ---
um solenoide modesto produz campo modesto. Para chegar a
$\SI{0.1}{\tesla}$ seriam necessários $\SI{80}{\ampere}$ com a mesma
geometria, o que exigiria refrigeração.\\
\textbf{Caminho alternativo:} inserir núcleo ferromagnético, que multiplica
$B$ por $\mu_r$ --- centenas ou milhares de vezes.}
\end{questao}

\begin{questao}[2]{}
Um fio conduz corrente constante. Compare o campo a $\SI{4}{\centi\meter}$ com
o campo a $\SI{12}{\centi\meter}$.
\resposta{$B_1/B_2=3$}
\resolucao{Como $B\propto1/r$:
\[
\frac{B_1}{B_2}=\frac{r_2}{r_1}=\frac{12}{4}=3.
\]
\textbf{Contraste importante com o campo elétrico:} o campo de uma carga
puntiforme cai com $1/r^{2}$, mas o de um fio \emph{infinito} cai com
$1/r$.\\
A razão é a mesma da eletrostática: o fio é uma fonte \emph{estendida} em uma
dimensão, e a superfície amperiana (um cilindro) tem área proporcional a $r$,
não a $r^{2}$.}
\end{questao}

\begin{questao}[2]{}
Uma espira circular de raio $\SI{5}{\centi\meter}$ deve produzir
$\SI{100}{\micro\tesla}$ em seu centro. Determine a corrente.
\resposta{$I\approx\SI{7.96}{\ampere}$}
\resolucao{Do campo no centro de uma espira, $B=\dfrac{\mu_0I}{2R}$:
\[
I=\frac{2RB}{\mu_0}=\frac{2(0{,}05)(\pot{1}{-4})}{4\pi\times10^{-7}}
=\SI{7.96}{\ampere}.
\]
\textbf{Comparação com o fio reto:} um fio retilíneo a
$\SI{5}{\centi\meter}$ precisaria de $\SI{25}{\ampere}$ para o mesmo campo. A
espira é mais eficiente porque \emph{todos} os seus elementos contribuem no
mesmo sentido no centro --- a razão entre as duas fórmulas é exatamente
$\pi$.}
\end{questao}

\begin{questao}[2]{}
Um solenoide tem $N$ espiras e comprimento $L$. Determine o que ocorre com o
campo interno se $N$ e $L$ forem \textbf{ambos dobrados}.
\resposta{O campo não se altera}
\resolucao{Como $B=\mu_0\dfrac{N}{L}I$, dobrar $N$ e $L$ mantém a razão
$N/L$ inalterada:
\[
B'=\mu_0\frac{2N}{2L}I=B.
\]
\textbf{Interpretação:} o solenoide dobrado equivale a dois solenoides
originais colocados em sequência. Cada trecho "enxerga" a mesma vizinhança de
espiras, e o campo local não muda.\\
\textbf{O que muda:} a região de campo uniforme fica mais longa, e a resistência
do enrolamento dobra --- exigindo o dobro da tensão para a mesma corrente.}
\end{questao}

\begin{questao}[2]{}
Dois fios longos e paralelos, distantes $\SI{8}{\centi\meter}$, conduzem
$\SI{4}{\ampere}$ cada, no \textbf{mesmo} sentido. Determine o campo no ponto
médio entre eles.
\resposta{$B=0$}
\resolucao{Cada fio produz, no ponto médio,
\[
B=\pot{2}{-7}\frac{4}{0{,}04}=\pot{2}{-5}\,\si{\tesla}.
\]
Pela regra da mão direita, com as correntes no \emph{mesmo} sentido, os dois
campos têm sentidos \textbf{opostos} no ponto médio --- e módulos iguais.
Cancelam-se:
\[
B_{total}=0.
\]
\textbf{Cuidado com a intuição:} correntes de mesmo sentido \emph{atraem-se},
mas seus campos se \emph{cancelam} entre elas. Força e campo são coisas
diferentes.\\
\textbf{Fora do par}, ao contrário, os campos se somam --- e é por isso que dois
fios juntos conduzindo no mesmo sentido produzem, de longe, o campo de um único
fio de corrente $2I$.}
\end{questao}

\begin{questao}[2]{}
Um fio reto é imerso em um meio de permeabilidade relativa $\mu_r=200$.
Determine o efeito sobre o campo, comparado ao ar.
\resposta{O campo fica $200$ vezes maior}
\resolucao{Em meio linear,
\[
B=\frac{\mu_r\mu_0I}{2\pi r},
\]
de modo que $B$ é multiplicado por $\mu_r=200$.\\
\textbf{Mecanismo:} o material se \emph{magnetiza}, e seus dipolos alinhados
somam sua própria contribuição ao campo aplicado. Diferentemente do dielétrico
--- que \emph{reduz} o campo elétrico --- o material ferromagnético o
\textbf{amplifica}.\\
\textbf{Ressalva séria:} $\mu_r$ não é constante. Ele cai drasticamente quando o
material satura, e a relação $B\times H$ deixa de ser linear. A fórmula acima
vale apenas em campos baixos.}
\end{questao}

\begin{questao}[2]{}
Determine a distância a que um fio de $\SI{50}{\ampere}$ produz campo igual ao
terrestre ($\SI{50}{\micro\tesla}$).
\resposta{$r=\SI{20}{\centi\meter}$}
\resolucao{
\[
r=\frac{\pot{2}{-7}I}{B}=\frac{(\pot{2}{-7})(50)}{\pot{5}{-5}}
=\SI{0.20}{\meter}.
\]
\textbf{Consequência prática:} a menos de $\SI{20}{\centi\meter}$ de um
condutor de $\SI{50}{\ampere}$, uma bússola aponta mais para o fio do que para
o norte.\\
É por isso que instrumentos magnéticos sensíveis --- e monitores de tubo, na
época em que existiam --- precisam ficar afastados de cabos de potência.}
\end{questao}

\begin{questao}[2]{}
Um toroide e um solenoide têm a mesma densidade de espiras e a mesma corrente.
Compare seus campos internos e externos.
\resposta{Campo interno praticamente igual; o toroide tem campo externo
\emph{rigorosamente} nulo}
\resolucao{\textbf{Campo interno:} ambos valem $\approx\mu_0nI$. No toroide,
$n=N/(2\pi r)$ varia ligeiramente com o raio, de modo que o campo não é
perfeitamente uniforme na seção --- mas para toroides finos a diferença é
pequena.\\
\textbf{Campo externo:} aqui está a diferença essencial.
\begin{itemize}
\item \textbf{Solenoide:} as linhas precisam retornar por fora, e o campo
externo, embora fraco, \emph{existe}.
\item \textbf{Toroide:} as linhas fecham-se \emph{inteiramente} dentro do
núcleo. Aplicando a lei de Ampère a um círculo externo, a corrente líquida
envolvida é \textbf{zero} --- e o campo também.
\end{itemize}
\textbf{Consequência de projeto:} indutores toroidais praticamente não irradiam
nem captam interferência, e podem ser montados lado a lado sem acoplamento
mútuo. É por isso que dominam aplicações de filtragem em fontes chaveadas.}
\end{questao}

\begin{questao}[2]{}
Dois fios perpendiculares entre si, ambos no mesmo plano da folha, conduzem
$\SI{6}{\ampere}$ e $\SI{8}{\ampere}$. Determine o campo resultante em um
ponto a $\SI{3}{\centi\meter}$ do primeiro e $\SI{4}{\centi\meter}$ do
segundo.
\resposta{$B=\SI{56.6}{\micro\tesla}$}
\resolucao{Cada fio produz campo perpendicular ao plano definido por ele e pelo
ponto. Como os fios são perpendiculares entre si e coplanares, ambos os campos
saem (ou entram) \emph{perpendicularmente à folha} --- portanto são
\textbf{paralelos} e se somam algebricamente:
\[
B_1=\pot{2}{-7}\frac{6}{0{,}03}=\SI{40}{\micro\tesla};
\qquad
B_2=\pot{2}{-7}\frac{8}{0{,}04}=\SI{40}{\micro\tesla}.
\]
Se os sentidos coincidirem: $B=\SI{80}{\micro\tesla}$; se forem opostos,
$B=0$.\\
\textbf{Atenção ao caso geral:} se os fios \emph{não} fossem coplanares, os
campos teriam direções diferentes e a soma seria vetorial --- com
$\SI{40}{\micro\tesla}$ perpendiculares, o resultado seria
$40\sqrt2=\SI{56.6}{\micro\tesla}$. \textbf{A geometria decide se a soma é
algébrica ou vetorial}, e é ela que precisa ser estabelecida antes de qualquer
conta.}
\end{questao}

\blocoaprofundamento

\begin{questao}[3]{}
Um condutor cilíndrico maciço de raio $a=\SI{2}{\milli\meter}$ conduz
$\SI{100}{\ampere}$ uniformemente distribuídos na seção.
\begin{itensq}
\item Determine o campo a $\SI{1}{\milli\meter}$, na superfície e a
      $\SI{5}{\milli\meter}$ do eixo.
\item Identifique onde o campo é máximo.
\end{itensq}
\resposta{$5$, $10$ e $\SI{4}{\milli\tesla}$; máximo na superfície}
\resolucao{\textbf{Dentro ($r<a$):} a superfície amperiana envolve apenas a
fração $(r/a)^{2}$ da corrente:
\[
B=\frac{\mu_0 I r}{2\pi a^{2}}
=\pot{2}{-7}\frac{100\,r}{(0{,}002)^{2}}.
\]
Em $r=\SI{1}{\milli\meter}$: $B=\SI{5}{\milli\tesla}$.\\
\textbf{Fora ($r>a$):} toda a corrente é envolvida, e vale a fórmula do fio
fino:
\[
B=\pot{2}{-7}\frac{100}{r}.
\]
Em $r=\SI{5}{\milli\meter}$: $B=\SI{4}{\milli\tesla}$.\\
\textbf{Na superfície ($r=a$):} as duas expressões coincidem em
$B=\SI{10}{\milli\tesla}$ --- o campo é contínuo.\\
(b) \textbf{O máximo está na superfície.} Dentro, $B$ cresce linearmente com
$r$; fora, decai com $1/r$. O pico ocorre exatamente em $r=a$.\\
\textbf{Analogia com a esfera isolante carregada:} lá o campo elétrico também
cresce linearmente dentro e cai com $1/r^{2}$ fora, com máximo na superfície. A
estrutura matemática é a mesma --- muda apenas a lei de Gauss ou de Ampère
aplicada.}
\end{questao}

\begin{questao}[3]{}
Um cabo coaxial tem condutor central de raio $a$ conduzindo $I$ e blindagem
externa conduzindo $I$ em sentido \textbf{oposto}, com $I=\SI{5}{\ampere}$.
\begin{itensq}
\item Determine o campo entre os condutores, a $\SI{2}{\milli\meter}$ do eixo.
\item Determine o campo fora do cabo.
\item Explique a consequência prática.
\end{itensq}
\resposta{$B=\SI{0.5}{\milli\tesla}$ entre os condutores; $B=0$ fora}
\resolucao{(a) Entre os condutores, a superfície amperiana envolve apenas a
corrente central:
\[
B=\pot{2}{-7}\frac{5}{0{,}002}=\pot{5}{-4}\,\si{\tesla}=\SI{0.5}{\milli\tesla}.
\]
(b) \textbf{Fora do cabo}, a superfície amperiana envolve $+I$ e $-I$ --- a
corrente líquida é \textbf{zero}:
\[
\oint\vec B\cdot\mathrm{d}\vec\ell=\mu_0(0)=0
\;\Longrightarrow\;
B=0.
\]
(c) \textbf{Consequência prática: o cabo coaxial não irradia nem capta campo
magnético.} Todo o campo fica confinado entre os condutores.\\
Isso o torna a escolha padrão para transmitir sinais em ambientes com
interferência --- e explica por que um cabo paralelo comum, cujos condutores
estão afastados, é muito pior: ali as correntes opostas não se cancelam
perfeitamente e o par forma uma espira que irradia.\\
\textbf{Mesma ideia, outra forma:} o par trançado obtém efeito semelhante ao
inverter a orientação da espira a cada torção, de modo que as contribuições
sucessivas se anulem.}
\end{questao}

\begin{questao}[3]{}
Uma espira de raio $R=\SI{10}{\centi\meter}$ conduz $\SI{5}{\ampere}$.
\begin{itensq}
\item Determine o campo no centro.
\item Determine o campo no eixo, a $\SI{10}{\centi\meter}$ e a
      $\SI{20}{\centi\meter}$ do centro.
\item Comente a taxa de queda.
\end{itensq}
\resposta{$31{,}4$; $11{,}1$ e $\SI{2.81}{\micro\tesla}$}
\resolucao{(a) $B_0=\dfrac{\mu_0I}{2R}=\dfrac{(4\pi\times10^{-7})(5)}{0{,}20}
=\pot{3.14}{-5}\,\si{\tesla}$.\\
(b) No eixo,
\[
B=\frac{\mu_0IR^{2}}{2\left(R^{2}+x^{2}\right)^{3/2}}.
\]
Em $x=R=\SI{0.10}{\meter}$:
$B=\pot{1.11}{-5}\,\si{\tesla}$, ou seja
$B_0/2^{3/2}=0{,}354\,B_0$.\\
Em $x=2R=\SI{0.20}{\meter}$: $B=\pot{2.81}{-6}\,\si{\tesla}$, cerca de
$0{,}089\,B_0$.\\
(c) \textbf{A queda é rápida.} A apenas um raio de distância, o campo já caiu
para $35\%$; a dois raios, para $9\%$.\\
Para $x\gg R$, a expressão tende a
$B\to\dfrac{\mu_0IR^{2}}{2x^{3}}$ --- decaimento com $1/x^{3}$, típico de
\textbf{dipolo}. É o mesmo expoente do dipolo elétrico, e pela mesma razão: a
espira não tem "monopolo magnético" para produzir termo $1/x^{2}$.}
\end{questao}

\begin{questao}[3]{}
Compare o campo no centro de um solenoide \emph{finito} com o do solenoide
ideal, e determine o campo em sua extremidade.
\begin{itensq}
\item Indique a condição para a aproximação ideal.
\item Determine o campo na extremidade.
\item Discuta a consequência para projeto.
\end{itensq}
\resposta{Na extremidade, $B\approx\mu_0nI/2$ --- metade do valor central}
\resolucao{(a) \textbf{Condição:} $L\gg D$, isto é, comprimento muito maior que
o diâmetro. Na prática, $L/D>10$ garante erro inferior a $1\%$ no centro.\\
(b) \textbf{Na extremidade}, apenas "metade" do solenoide contribui --- as
espiras estão todas de um lado. O resultado exato para solenoide longo é
\[
B_{extremidade}=\frac{\mu_0nI}{2},
\]
exatamente \textbf{metade} do valor central.\\
(c) \textbf{Consequências de projeto.}
\begin{itemize}
\item A região de campo verdadeiramente uniforme é \emph{menor} que o
solenoide: ela se restringe à parte central, tipicamente à metade do
comprimento.
\item Para uniformidade melhor que $1\%$ em uma região de comprimento
$\ell$, o solenoide precisa ter comprimento bem maior que $\ell$ --- ou usar
enrolamento com densidade \emph{corrigida} nas extremidades.
\item A alternativa clássica é o par de \textbf{bobinas de Helmholtz}: duas
bobinas iguais separadas por uma distância igual ao raio, arranjo que anula a
segunda derivada do campo no centro e produz região uniforme com pouco
material.
\end{itemize}}
\end{questao}

\begin{questao}[3]{}
Três fios longos e paralelos, coplanares e igualmente espaçados de
$\SI{5}{\centi\meter}$, conduzem $\SI{10}{\ampere}$ cada, todos no mesmo
sentido. Determine o campo no fio central.
\resposta{$B=0$ no fio central}
\resolucao{O fio central não produz campo sobre si mesmo. Restam as
contribuições dos dois vizinhos, ambos a $\SI{5}{\centi\meter}$:
\[
B_1=B_2=\pot{2}{-7}\frac{10}{0{,}05}=\SI{40}{\micro\tesla}.
\]
Pela regra da mão direita, com os dois vizinhos conduzindo no mesmo sentido e
estando em \emph{lados opostos} do fio central, seus campos ali têm sentidos
\textbf{contrários}:
\[
B=40-40=0.
\]
\textbf{Mas a força não é nula --- ela também é.} Como $\vec F=I\vec\ell\times
\vec B$ e $\vec B=\vec 0$ no fio central, a força sobre ele é zero: as atrações
dos dois vizinhos se equilibram.\\
\textbf{Nos fios externos, porém}, a situação é diferente: cada um sofre atração
resultante para o centro. O conjunto tende a se \emph{comprimir} --- efeito real
em barramentos de alta corrente, que precisam de espaçadores mecânicos.}
\end{questao}

\begin{questao}[3]{}
Um eletroímã deve produzir $\SI{0.1}{\tesla}$ com núcleo de ar, usando
solenoide de $\SI{20}{\centi\meter}$.
\begin{itensq}
\item Determine o produto $NI$ necessário.
\item Avalie a viabilidade com $N=1000$ espiras.
\item Proponha alternativa.
\end{itensq}
\resposta{$NI=\SI{15.9}{\kilo\ampere}$-espira; exigiria $\SI{15.9}{\ampere}$
com $1000$ espiras}
\resolucao{(a) De $B=\mu_0NI/L$:
\[
NI=\frac{BL}{\mu_0}=\frac{(0{,}1)(0{,}20)}{4\pi\times10^{-7}}
=\pot{1.59}{4}\ \text{A-espira}.
\]
(b) Com $N=1000$: $I=\SI{15.9}{\ampere}$.\\
\textbf{Viável, mas problemático.} Mil espiras conduzindo
$\SI{16}{\ampere}$ exigem fio grosso (a densidade de corrente admissível limita
a seção), o que aumenta o volume do enrolamento. E a dissipação é considerável:
com resistência de apenas $\SI{1}{\ohm}$, seriam
$\SI{256}{\watt}$ --- exigindo refrigeração.\\
(c) \textbf{Alternativa: núcleo ferromagnético.} Com $\mu_r=1000$, a mesma
corrente produziria $\SI{100}{\tesla}$ --- valor impossível, porque o material
\textbf{satura} em torno de $1{,}5$ a $\SI{2}{\tesla}$.\\
\textbf{O que o núcleo realmente entrega:} os $\SI{0.1}{\tesla}$ pedidos com
corrente cerca de mil vezes menor ($\SI{16}{\milli\ampere}$), enquanto o
material permanecer na região linear.\\
\textbf{Limite físico:} acima da saturação, o núcleo deixa de ajudar e o
eletroímã volta a se comportar como se fosse de ar. É por isso que campos acima
de $\SI{2}{\tesla}$ exigem bobinas supercondutoras, sem núcleo.}
\end{questao}

\begin{questao}[3]{}
Um cabo de par paralelo tem condutores separados por $\SI{2}{\milli\meter}$,
conduzindo $\SI{1}{\ampere}$ em sentidos opostos.
\begin{itensq}
\item Determine o campo no ponto médio entre eles.
\item Determine o campo a $\SI{10}{\centi\meter}$ do par.
\item Compare com um fio único.
\end{itensq}
\resposta{$\SI{400}{\micro\tesla}$ no meio; $\approx\SI{0.04}{\micro\tesla}$ a
$\SI{10}{\centi\meter}$}
\resolucao{(a) No ponto médio, cada fio dista $\SI{1}{\milli\meter}$, e com
correntes \emph{opostas} os campos se \textbf{somam}:
\[
B=2\times\pot{2}{-7}\frac{1}{0{,}001}=\SI{400}{\micro\tesla}.
\]
(b) A $\SI{10}{\centi\meter}$, os dois fios estão praticamente à mesma
distância e seus campos quase se cancelam. O resultado é aproximadamente o de um
\textbf{dipolo}:
\[
B\approx\frac{\mu_0 I d}{2\pi r^{2}}
=\pot{2}{-7}\frac{(1)(0{,}002)}{(0{,}10)^{2}}
=\pot{4}{-8}\,\si{\tesla}.
\]
(c) \textbf{Um fio único} de $\SI{1}{\ampere}$ produziria, a
$\SI{10}{\centi\meter}$:
$B=\pot{2}{-6}\,\si{\tesla}$ --- \textbf{cinquenta vezes mais}.\\
\textbf{Conclusão:} manter os condutores de ida e volta \emph{juntos} reduz
drasticamente o campo irradiado, porque o par se comporta como dipolo
($1/r^{2}$) em vez de fio isolado ($1/r$). Quanto menor a separação $d$, melhor
--- e é exatamente por isso que se recomenda não separar os condutores de um
mesmo circuito.}
\end{questao}

\begin{questao}[3]{}
Um sensor Hall indica $\SI{2.5}{\milli\volt}$ quando exposto a
$\SI{100}{\milli\tesla}$, com resposta linear e tensão de repouso nula.
\begin{itensq}
\item Determine a sensibilidade.
\item Determine o campo correspondente a uma leitura de
      $\SI{0.6}{\milli\volt}$.
\item Discuta as fontes de erro.
\end{itensq}
\resposta{$\SI{25}{\micro\volt\per\milli\tesla}$; $B=\SI{24}{\milli\tesla}$}
\resolucao{(a)
\[
S=\frac{2{,}5\,\si{\milli\volt}}{100\,\si{\milli\tesla}}
=\SI{25}{\micro\volt\per\milli\tesla}.
\]
(b) $B=\dfrac{0{,}6}{0{,}025}=\SI{24}{\milli\tesla}$.\\
(c) \textbf{Fontes de erro, em ordem de importância.}
\begin{itemize}
\item \textbf{Tensão de \emph{offset}:} sensores reais indicam algo mesmo com
$B=0$, por assimetria construtiva. Valores de dezenas de microvolts são comuns
--- comparáveis ao sinal de alguns militesla. \emph{Correção:} medir com o
sensor blindado e subtrair.
\item \textbf{Deriva térmica:} tanto o \emph{offset} quanto a sensibilidade
variam com a temperatura.
\item \textbf{Orientação:} o sensor responde apenas à componente
\emph{perpendicular} à sua face. Um erro angular $\theta$ introduz fator
$\cos\theta$ --- desprezível para ângulos pequenos, mas $13\%$ a
$\ang{30}$.
\item \textbf{Campo terrestre} somado à medida, se não for compensado.
\end{itemize}
\textbf{Procedimento robusto:} medir com o sensor em uma orientação, depois
girá-lo $\ang{180}$ e medir de novo. A média das duas leituras cancela o
\emph{offset}, e a semidiferença dá o campo.}
\end{questao}

\blocodesafio

\begin{questao}[4]{}
\textbf{(Dedução por Biot--Savart.)} Deduza o campo magnético no centro de uma
espira circular de raio $R$ percorrida por corrente $I$.
\begin{itensq}
\item Aplique a lei de Biot--Savart a um elemento de corrente.
\item Integre e obtenha o resultado.
\item Explique por que a integração é simples neste caso.
\end{itensq}
\resposta{$B=\dfrac{\mu_0 I}{2R}$}
\resolucao{(a) A lei de Biot--Savart estabelece que um elemento
$I\,\mathrm{d}\vec\ell$ produz, a uma distância $\vec r$,
\[
\mathrm{d}\vec B=\frac{\mu_0}{4\pi}\,
\frac{I\,\mathrm{d}\vec\ell\times\hat r}{r^{2}}.
\]
Para o \textbf{centro} da espira, todo elemento está à mesma distância $R$, e
$\mathrm{d}\vec\ell$ é sempre \emph{perpendicular} a $\hat r$ --- logo
$|\mathrm{d}\vec\ell\times\hat r|=\mathrm{d}\ell$.\\
(b) Além disso, todos os $\mathrm{d}\vec B$ apontam na \emph{mesma} direção
(perpendicular ao plano da espira), de modo que a integral vetorial vira
escalar:
\[
B=\frac{\mu_0I}{4\pi R^{2}}\oint\mathrm{d}\ell
=\frac{\mu_0I}{4\pi R^{2}}(2\pi R)
=\frac{\mu_0 I}{2R}.
\]
(c) \textbf{Três simplificações coincidiram:} distância constante, produto
vetorial de módulo máximo (ângulo reto) e contribuições todas paralelas. É a
combinação mais favorável possível.\\
\textbf{Fora do centro isso se perde:} no eixo, a distância deixa de ser $R$ e
as contribuições ganham componentes radiais que se cancelam --- exigindo
projeção. Fora do eixo, a integral não tem forma fechada elementar e recorre-se a
métodos numéricos ou a funções elípticas.}
\end{questao}

\begin{questao}[4]{}
\textbf{(Ampère contra Biot--Savart.)} Compare os dois métodos de cálculo de
campo magnético.
\begin{itensq}
\item Enuncie cada lei e sua condição de uso.
\item Determine quando a lei de Ampère é praticável.
\item Classifique quatro configurações quanto ao método adequado.
\end{itensq}
\resposta{Ampère exige simetria suficiente para tirar $B$ da integral;
Biot--Savart é geral mas trabalhoso}
\resolucao{(a) \textbf{Biot--Savart:}
$\mathrm{d}\vec B=\dfrac{\mu_0}{4\pi}
\dfrac{I\,\mathrm{d}\vec\ell\times\hat r}{r^{2}}$ --- vale \emph{sempre}, para
qualquer distribuição de corrente. É a lei fundamental.\\
\textbf{Ampère:} $\oint\vec B\cdot\mathrm{d}\vec\ell=\mu_0I_{env}$ --- também
sempre verdadeira, mas só \emph{útil} para calcular $B$ em casos especiais.\\
(b) \textbf{Ampère é praticável quando} existe uma curva fechada ao longo da
qual $|\vec B|$ é constante e $\vec B$ é paralelo (ou perpendicular) ao
caminho. Só assim $B$ pode ser retirado da integral:
\[
\oint\vec B\cdot\mathrm{d}\vec\ell=B\oint\mathrm{d}\ell=B\,\ell.
\]
Isso exige \textbf{simetria} --- cilíndrica, plana ou toroidal --- que precisa
ser \emph{conhecida de antemão}, por argumento físico.\\
(c)
\begin{center}
\begin{tabular}{lll}
\hline
Configuração & Método & Motivo\\
\hline
Fio retilíneo infinito & Ampère & simetria cilíndrica\\
Solenoide/toroide ideal & Ampère & simetria evidente\\
Espira circular (centro) & Biot--Savart & sem curva de $B$ constante\\
Segmento finito de fio & Biot--Savart & sem simetria alguma\\
\hline
\end{tabular}
\end{center}
\textbf{Erro comum:} tentar aplicar Ampère a uma espira circular usando a
própria espira como caminho. O campo \emph{não} é constante nem tangente ao
longo dela --- a integral existe, mas não permite isolar $B$.\\
\textbf{Analogia com a eletrostática:} Gauss está para Coulomb assim como
Ampère está para Biot--Savart. Em ambos os casos, a lei integral é elegante e
poderosa \emph{apenas} quando a simetria coopera.}
\end{questao}

\begin{questao}[4]{}
\textbf{(Projeto de solenoide.)} Projete um solenoide de $500$ espiras e
$\SI{25}{\centi\meter}$ para produzir $\SI{1}{\milli\tesla}$.
\begin{itensq}
\item Determine a corrente necessária.
\item Determine a potência dissipada, para fio de resistência total
      $\SI{2}{\ohm}$.
\item Avalie a uniformidade do campo.
\end{itensq}
\resposta{$I=\SI{0.398}{\ampere}$; $P=\SI{0.32}{\watt}$}
\resolucao{(a) $n=\dfrac{500}{0{,}25}=\SI{2000}{\per\meter}$, e
\[
I=\frac{B}{\mu_0 n}=\frac{\pot{1}{-3}}{(4\pi\times10^{-7})(2000)}
=\SI{0.398}{\ampere}.
\]
(b) $P=RI^{2}=2(0{,}398)^{2}=\SI{0.32}{\watt}$ --- desprezível, sem necessidade
de refrigeração.\\
(c) \textbf{Uniformidade.} Supondo diâmetro de $\SI{4}{\centi\meter}$, a razão
$L/D=6{,}25$ --- razoável, mas não excelente. Consequências:
\begin{itemize}
\item No \textbf{centro}, o campo fica cerca de $1$ a $2\%$ abaixo do valor
ideal $\mu_0nI$.
\item Nas \textbf{extremidades}, cai para aproximadamente metade.
\item A região com uniformidade melhor que $1\%$ limita-se ao terço central,
cerca de $\SI{8}{\centi\meter}$.
\end{itemize}
\textbf{Se for necessária melhor uniformidade:} alongar o solenoide (mantendo
$n$, o que exige mais espiras e mais tensão), ou adotar bobinas de Helmholtz.\\
\textbf{Verificação de coerência:} $\SI{1}{\milli\tesla}$ é vinte vezes o campo
terrestre --- o experimento precisa ser orientado ou compensado, sob pena de o
campo ambiente introduzir erro de $5\%$.}
\end{questao}

\begin{questao}[4]{}
\textbf{(Dipolo magnético.)} Uma espira de raio $R$ e corrente $I$ é observada
no eixo, a distância $x\gg R$.
\begin{itensq}
\item Obtenha a forma assintótica do campo.
\item Defina o momento de dipolo magnético e reescreva o resultado.
\item Compare com o dipolo elétrico e comente a ausência de monopolos.
\end{itensq}
\resposta{$B\to\dfrac{\mu_0 m}{2\pi x^{3}}$, com $m=I\pi R^{2}$}
\resolucao{(a) Partindo de
$B=\dfrac{\mu_0IR^{2}}{2(R^{2}+x^{2})^{3/2}}$ e fazendo $x\gg R$:
\[
B\to\frac{\mu_0IR^{2}}{2x^{3}}.
\]
(b) Definindo o \textbf{momento de dipolo magnético}
$m=I\,A=I\pi R^{2}$:
\[
B=\frac{\mu_0 m}{2\pi x^{3}}.
\]
\textbf{Verificação numérica:} para $R=\SI{10}{\centi\meter}$ e
$I=\SI{5}{\ampere}$, tem-se $m=\SI{0.157}{\ampere\meter\squared}$ e, a
$\SI{1}{\meter}$, $B=\pot{3.14}{-8}\,\si{\tesla}$. O valor exato é
$\pot{3.10}{-8}$ --- a aproximação erra $1{,}5\%$ a dez raios de distância.\\
(c) \textbf{Comparação com o dipolo elétrico.} Ambos decaem com $1/r^{3}$, e as
expressões são formalmente análogas ($p=2qa$ lá, $m=IA$ aqui).\\
\textbf{Mas há uma diferença fundamental.} O dipolo elétrico é feito de duas
cargas que \emph{podem} ser separadas --- e então cada uma produz campo
$1/r^{2}$. Já o dipolo magnético \textbf{não pode} ser separado: não existem
monopolos magnéticos.\\
\textbf{Consequência:} o dipolo é o termo \emph{dominante} de qualquer fonte
magnética estacionária. Não há termo $1/r^{2}$ em magnetostática --- e é por isso
que $\nabla\cdot\vec B=0$ figura entre as equações de Maxwell, sem análogo
elétrico.}
\end{questao}

\begin{questao}[4]{}
\textbf{(Blindagem magnética.)} Compare a blindagem de campos magnéticos com a
de campos elétricos.
\begin{itensq}
\item Explique por que uma gaiola metálica não blinda campo magnético estático.
\item Descreva os dois mecanismos disponíveis.
\item Discuta a dependência com a frequência.
\end{itensq}
\resposta{Não há "cargas magnéticas" a redistribuir; usa-se desvio por material
de alto $\mu_r$ ou correntes induzidas}
\resolucao{(a) \textbf{Por que a gaiola falha.} A blindagem elétrica funciona
porque o condutor possui \emph{cargas livres} que se redistribuem até anular o
campo interno. Não existe equivalente magnético: \textbf{não há monopolos
magnéticos livres} para migrar.\\
Um campo magnético \emph{estático} atravessa o cobre e o alumínio praticamente
sem atenuação.\\
(b) \textbf{Dois mecanismos disponíveis.}
\begin{itemize}
\item \textbf{Desvio de fluxo} (campos estáticos ou lentos). Um invólucro de
material de alto $\mu_r$ --- \emph{mu-metal}, com $\mu_r$ de dezenas de
milhares --- oferece caminho de baixa relutância. O fluxo prefere circular pelo
material e "contorna" a região protegida. Não há cancelamento: há \emph{desvio}.
\item \textbf{Correntes induzidas} (campos variáveis). Pela lei de Lenz, um
campo alternado induz correntes em um invólucro condutor, e essas correntes
geram campo oposto. A atenuação cresce com a frequência e com a condutividade.
\end{itemize}
(c) \textbf{Dependência com a frequência.}
\begin{center}
\begin{tabular}{lll}
\hline
Faixa & Mecanismo eficaz & Material\\
\hline
Estático (CC) & desvio de fluxo & mu-metal\\
Baixa (dezenas de Hz) & desvio, pouca indução & mu-metal\\
Alta (kHz e acima) & correntes induzidas & cobre, alumínio\\
\hline
\end{tabular}
\end{center}
\textbf{Consequência prática:} blindar campo elétrico é fácil e barato --- basta
qualquer metal aterrado. Blindar campo magnético estático é caro, exige material
especial, e a atenuação típica fica em uma ou duas ordens de grandeza, contra
muitas ordens no caso elétrico.\\
É por isso que salas de ressonância magnética têm blindagem magnética de
toneladas, enquanto a blindagem elétrica de um cabo é uma malha fina.}
\end{questao}

\begin{questao}[4]{}
\textbf{(Mapeamento experimental.)} Elabore um procedimento para mapear o campo
magnético de uma bobina usando sensor Hall.
\begin{itensq}
\item Descreva o procedimento e o tratamento do \emph{offset}.
\item Indique como obter as três componentes do campo.
\item Estabeleça o critério de validação dos dados.
\end{itensq}
\resposta{Medir com inversão de orientação para cancelar \emph{offset};
três orientações ortogonais para as componentes}
\resolucao{(a) \textbf{Procedimento.}
\begin{enumerate}
\item \textbf{Caracterizar o zero:} com a bobina \emph{desligada}, registrar a
leitura em cada ponto. Isso captura tanto o \emph{offset} do sensor quanto o
campo ambiente (terrestre e de equipamentos próximos).
\item \textbf{Medir com a bobina ligada} nos mesmos pontos.
\item \textbf{Subtrair} ponto a ponto. O resultado é o campo da bobina,
livre de \emph{offset} e de campo ambiente.
\end{enumerate}
\textbf{Alternativa mais robusta:} inverter o sentido da corrente na bobina e
tomar a \emph{semidiferença} das duas leituras. Isso cancela tudo o que não
depende da corrente.\\
(b) \textbf{Três componentes.} O sensor Hall responde apenas à componente
perpendicular à sua face. Para obter $\vec B$ completo, meça em três orientações
ortogonais no mesmo ponto:
\[
|\vec B|=\sqrt{B_x^{2}+B_y^{2}+B_z^{2}}.
\]
Um sensor triaxial faz isso simultaneamente e evita erros de reposicionamento.\\
(c) \textbf{Critérios de validação.}
\begin{itemize}
\item \textbf{Linearidade com a corrente:} $B$ deve ser proporcional a $I$.
Meça em três correntes e verifique --- desvio indica saturação de núcleo ou
não linearidade do sensor.
\item \textbf{Simetria:} pontos simétricos em relação ao eixo devem dar valores
iguais. Discrepância revela desalinhamento do sensor ou da bobina.
\item \textbf{Valor central:} deve concordar com $\mu_0nI$ dentro da incerteza,
descontada a correção de comprimento finito.
\item \textbf{Reprodutibilidade:} repita o primeiro ponto ao final. Diferença
indica deriva térmica do sensor.
\end{itemize}
\textbf{Registre sempre a temperatura} --- a sensibilidade Hall varia
tipicamente algumas centenas de ppm por grau.}
\end{questao}

\begin{questao}[4]{}
\textbf{(Limites de eletroímãs.)} Analise o que limita, na prática, o campo
magnético alcançável.
\begin{itensq}
\item Determine o limite imposto pela saturação do núcleo.
\item Determine o limite térmico de bobinas de cobre.
\item Descreva as soluções empregadas para superá-los.
\end{itensq}
\resposta{Núcleos saturam entre $1{,}5$ e $\SI{2}{\tesla}$; acima disso é
preciso bobina sem núcleo, com corrente muito alta}
\resolucao{(a) \textbf{Saturação.} Materiais ferromagnéticos alinham seus
domínios completamente em torno de $1{,}5$ a $\SI{2}{\tesla}$ (ferro-silício,
$\approx\SI{2}{\tesla}$; ferrites, bem menos). Acima disso, $\mu_r$ cai para
próximo de $1$ e o núcleo deixa de contribuir --- \textbf{acrescentar corrente
passa a render tão pouco quanto no ar}.\\
Este é um limite \emph{de material}, não de projeto: nenhuma geometria o
contorna.\\
(b) \textbf{Limite térmico.} Sem núcleo, $B=\mu_0nI$ exige corrente enorme. A
potência dissipada vai com $I^{2}$, e a densidade de corrente admissível em
cobre refrigerado a água fica em torno de
$\SI{20}{\ampere\per\milli\meter\squared}$.\\
Eletroímãs resistivos de laboratório atingem $1$ a $\SI{2}{\tesla}$ consumindo
\emph{megawatts} --- e boa parte do custo operacional é a refrigeração.\\
(c) \textbf{Soluções empregadas.}
\begin{itemize}
\item \textbf{Supercondutores:} resistência nula elimina o limite térmico.
Bobinas de NbTi e $\mathrm{Nb_3Sn}$ chegam a $\SI{20}{\tesla}$ e são a base de
equipamentos de ressonância magnética e de aceleradores. O custo passa a ser a
criogenia.
\item \textbf{Campos pulsados:} correntes altíssimas por milissegundos, antes
que o calor se acumule. Alcançam $\SI{100}{\tesla}$, com a bobina
frequentemente destruída no processo.
\item \textbf{Ímãs híbridos:} bobina supercondutora externa somada a uma
resistiva interna --- combinação que detém os recordes de campo contínuo, acima
de $\SI{45}{\tesla}$.
\end{itemize}
\textbf{Limite último:} acima de algumas dezenas de tesla, a \emph{pressão
magnética} $B^{2}/2\mu_0$ supera a resistência mecânica de qualquer material. A
$\SI{100}{\tesla}$ ela vale cerca de $\SI{4}{\giga\pascal}$ --- a bobina
literalmente se rompe.}
\end{questao}

\begin{questao}[4]{}
\textbf{(Superposição e simetria.)} Um condutor cilíndrico de raio $a$ possui
uma cavidade cilíndrica de raio $b$, paralela ao eixo e deslocada de $d$.
A corrente $I$ distribui-se uniformemente na seção restante.
\begin{itensq}
\item Formule a estratégia de solução por superposição.
\item Determine o campo dentro da cavidade.
\item Interprete o resultado.
\end{itensq}
\resposta{O campo na cavidade é \textbf{uniforme}, de módulo
$\mu_0Jd/2$}
\resolucao{(a) \textbf{Estratégia.} O condutor furado equivale à
\emph{superposição} de:
\begin{itemize}
\item um cilindro \textbf{maciço} de raio $a$ com densidade de corrente
$+J$; \emph{mais}
\item um cilindro de raio $b$, na posição da cavidade, com densidade
$-J$.
\end{itemize}
A soma reproduz corrente $J$ na região sólida e zero na cavidade --- exatamente
a situação real.\\
(b) Dentro de um cilindro maciço de densidade $J$, a lei de Ampère dá
$B=\dfrac{\mu_0 J r}{2}$, ou vetorialmente
$\vec B=\dfrac{\mu_0}{2}\,\vec J\times\vec r$, com $\vec r$ medido a partir do
\emph{eixo daquele} cilindro.\\
Somando as duas contribuições em um ponto da cavidade, com $\vec r_1$ e
$\vec r_2$ medidos a partir de cada eixo:
\[
\vec B=\frac{\mu_0}{2}\vec J\times\vec r_1
-\frac{\mu_0}{2}\vec J\times\vec r_2
=\frac{\mu_0}{2}\vec J\times(\vec r_1-\vec r_2)
=\frac{\mu_0}{2}\vec J\times\vec d.
\]
\textbf{O vetor $\vec d$ é constante} --- ele liga os dois eixos e não depende do
ponto.\\
(c) \textbf{Interpretação: o campo na cavidade é rigorosamente uniforme}, de
módulo $\mu_0Jd/2$ e direção perpendicular a $\vec d$.\\
Isso é notável: uma geometria sem simetria evidente produz, na região vazia, o
campo mais simples possível.\\
\textbf{Aplicação:} o mesmo resultado, com $\vec J$ substituído por densidade de
carga, dá campo elétrico uniforme na cavidade de uma esfera carregada --- e é
usado no projeto de aceleradores e de ímãs de deflexão, em que se busca campo
uniforme em uma região de acesso.}
\end{questao}

% END BANCO COMPLEMENTAR

\gabarito[12.2 Campo magnético]

% =====================================================================
\subsection{Fluxo magnético}
% =====================================================================

\blocofixacao

\begin{questao}[1]{}
O fluxo magnético através de uma espira é máximo quando o plano da espira está:
\begin{alternativas}
\item paralelo ao campo.
\item perpendicular ao campo (campo atravessando a espira frontalmente).
\item inclinado a $\SI{45}{\degree}$ com o campo.
\item em qualquer orientação (o fluxo não depende do ângulo).
\item girando rapidamente.
\end{alternativas}
\resposta{(b)}
\resolucao{Como $\Phi = B\,A\cos\theta$, sendo $\theta$ o ângulo entre o campo e a
\emph{normal} à espira, o fluxo é máximo quando $\theta = 0$ --- ou seja, quando o
campo atravessa a espira perpendicularmente ao seu plano. Se o plano fica paralelo
ao campo, $\theta = \SI{90}{\degree}$ e o fluxo é nulo.}
\end{questao}

\begin{questao}[1]{}
Uma espira de área $\SI{20}{\centi\meter\squared}$ é atravessada
perpendicularmente por um campo uniforme de $\SI{0.5}{\tesla}$. Calcule o fluxo
magnético.
\resposta{$\Phi = \SI{1}{\milli\weber}$}
\resolucao{Com $\theta = 0$ ($\cos\theta = 1$) e
$A = \SI{20}{\centi\meter\squared} = \pot{20}{-4}\,\si{\meter\squared}$:
\[
\Phi = B\,A = 0{,}5\cdot\pot{20}{-4} = \pot{1}{-3}\,\si{\weber}=\SI{1}{\milli\weber}.
\]}
\end{questao}

\blocoaplicacao

\begin{questao}[2]{}
Uma espira circular de raio $\SI{8}{\centi\meter}$ está imersa em um campo
uniforme de $\SI{0.3}{\tesla}$, perpendicular ao seu plano.
\begin{itensq}
\item Calcule o fluxo magnético através da espira.
\item Se o campo variar linearmente de $\SI{0.3}{}$ para $\SI{0.1}{\tesla}$ em
      $\SI{0.5}{\second}$, qual é a fem induzida?
\end{itensq}
\resposta{(a) $\Phi \approx \SI{6.03}{\milli\weber}$; (b) $\varepsilon \approx \SI{8.04}{\milli\volt}$}
\resolucao{(a) A área é $A = \pi R^2 = \pi(0{,}08)^2 \approx \SI{0.0201}{\meter\squared}$:
\[
\Phi = B\,A = 0{,}3\cdot0{,}0201 \approx \SI{6.03}{\milli\weber}.
\]
(b) A fem é o módulo da taxa de variação do fluxo. Como só $B$ muda:
\[
\varepsilon = \left|\frac{\Delta\Phi}{\Delta t}\right|
= A\,\frac{|\Delta B|}{\Delta t}
= 0{,}0201\cdot\frac{0{,}2}{0{,}5}\approx\SI{8.04}{\milli\volt}.
\]}
\end{questao}

\begin{questao}[2]{Apostila original revisada}
Um campo magnético uniforme de módulo $B=\SI{1e-4}{\tesla}$ atravessa uma espira de
área $A=\SI{1e-5}{\metre\squared}$. O campo forma $\SI{45}{\degree}$ com a normal
ao plano da espira. Calcule o fluxo magnético.
\resposta{$\Phi=\SI{7.07e-10}{\weber}$}
\resolucao{$\Phi=BA\cos\theta=(10^{-4})(10^{-5})\cos45^\circ
=7{,}07\times10^{-10}\,\si{\weber}$.}
\end{questao}

\blocoaprofundamento

\begin{questao}[3]{}
O gráfico mostra a variação do fluxo magnético $\Phi(t)$ através de uma espira.
Determine a fem induzida em cada trecho e esboce mentalmente a forma de onda da
fem.

\circuitoqq{%
\begin{tikzpicture}
 \begin{axis}[width=10cm,height=5cm,xlabel={$t$ (s)},ylabel={$\Phi$ (mWb)},
   xmin=0,xmax=6,ymin=0,ymax=5,xtick={0,1,2,3,4,5,6},ytick={0,2,4},
   grid=both,axis lines=left]
   \addplot[Icol,very thick] coordinates {(0,0)(2,4)(4,4)(6,0)};
 \end{axis}
\end{tikzpicture}}{Fluxo magnético em função do tempo.}

\resposta{$0$--$2\,$s: $\varepsilon=\SI{2}{\milli\volt}$; $2$--$4\,$s:
$\varepsilon=0$; $4$--$6\,$s: $\varepsilon=\SI{-2}{\milli\volt}$}
\resolucao{A fem é a \emph{inclinação} do gráfico $\Phi \times t$ (com sinal
trocado, por Lenz):
\begin{itemize}[leftmargin=1.6em,itemsep=1pt,topsep=2pt]
\item \textbf{$0$ a $\SI{2}{\second}$:} $\Phi$ sobe de $0$ a
      \SI{4}{\milli\weber}, inclinação $\dfrac{4}{2} = \SI{2}{\milli\weber\per\second}$,
      logo $|\varepsilon| = \SI{2}{\milli\volt}$.
\item \textbf{$2$ a $\SI{4}{\second}$:} $\Phi$ constante, inclinação nula,
      $\varepsilon = 0$.
\item \textbf{$4$ a $\SI{6}{\second}$:} $\Phi$ desce de \SI{4}{} a $0$, inclinação
      $-\SI{2}{\milli\weber\per\second}$, $\varepsilon = \SI{-2}{\milli\volt}$
      (sentido oposto).
\end{itemize}
A fem só existe quando há \emph{variação} de fluxo, e seu sinal se inverte entre a
subida e a descida --- exatamente o princípio de um gerador.}
\end{questao}

\blocodesafio

\begin{questao}[4]{}
Uma espira retangular de área $A = \SI{200}{\centi\meter\squared}$ gira com
velocidade angular constante $\omega$ em um campo magnético uniforme
$B = \SI{0.4}{\tesla}$, completando $\SI{1800}{}$ rotações por minuto.
\begin{itensq}
\item Escreva a expressão do fluxo $\Phi(t)$ e calcule seu valor máximo.
\item Obtenha a fem induzida $\varepsilon(t)$ e seu valor de pico.
\item Determine a frequência da tensão gerada e explique por que a saída é
      \emph{senoidal}.
\item Se a espira tivesse $N = 100$ voltas, o que mudaria?
\end{itensq}
\resposta{(a) $\Phi_{\max}=\SI{8}{\milli\weber}$;
(b) $\varepsilon_{\text{pico}}\approx\SI{1.51}{\volt}$; (c) \SI{30}{\hertz};
(d) tudo $\times100$}
\resolucao{\textbf{(a)} Com a espira girando, o ângulo entre a normal e o campo
varia: $\theta = \omega t$. Logo
\[
\Phi(t)=B\,A\cos(\omega t),
\qquad
\Phi_{\max}=B\,A=0{,}4\cdot\pot{200}{-4}=\SI{8}{\milli\weber}.
\]
\textbf{(b)} A velocidade angular é
\[
\omega = \frac{2\pi\cdot1800}{60}=60\pi\approx\SI{188.5}{\radian\per\second}.
\]
Pela lei de Faraday,
\[
\varepsilon(t)=-\frac{\mathrm{d}\Phi}{\mathrm{d}t}
= B A\,\omega\sin(\omega t),
\qquad
\varepsilon_{\text{pico}}=BA\omega
=(\pot{8}{-3})(188{,}5)\approx\SI{1.51}{\volt}.
\]
\textbf{(c)} A frequência é
$f = \dfrac{1800}{60} = \SI{30}{\hertz}$ (uma volta $=$ um ciclo completo).\\
\textbf{Por que senoidal?} Porque o \emph{fluxo} varia cossenoidalmente com o
ângulo (projeção da área sobre a direção do campo), e a fem é sua
\emph{derivada}. A derivada do cosseno é o seno --- daí a defasagem de
$\SI{90}{\degree}$ entre fluxo e tensão. Este é literalmente o princípio de
funcionamento do \textbf{alternador}: a forma de onda senoidal da rede elétrica
nasce da geometria circular do movimento.

\textbf{(d)} Com $N$ espiras, todas concatenam o mesmo fluxo e suas fem se somam
em série:
\[
\varepsilon_{\text{pico}} = N B A \omega \approx 100\cdot1{,}51 = \SI{151}{\volt}.
\]
O fluxo \emph{por espira} não muda; o que muda é o \emph{fluxo concatenado}
$N\Phi$. É assim que geradores reais atingem centenas de volts com campos
modestos: aumentando $N$, $B$, $A$ ou $\omega$ --- os quatro parâmetros de
projeto de uma máquina elétrica.}
\end{questao}

% BEGIN BANCO COMPLEMENTAR Fluxo magnético
% =====================================================================
%  Banco complementar --- Fluxo Magnetico  (REVISADO)
%  Substitui a versao gerada por template (6 clones em N1, 10 em N2,
%  11 em N3 e 7 questoes conceituais marcadas como N4).
%  O cap11 ja cobre: condicao de fluxo maximo (MC), espira de 20 cm2
%  perpendicular a 0,5 T, espira circular de R=8 cm em 0,3 T, campo de
%  1e-4 T em area de 1e-5 m2, grafico Phi(t) com fem por trecho e --
%  como N4 -- a ESPIRA RETANGULAR GIRANDO com velocidade angular
%  constante. Esta ultima motivou a remocao da N4 antiga sobre fluxo em
%  espira girante, que a duplicava.
%  Este banco explora: campo nao uniforme por integracao, fluxo de um
%  fio reto atraves de espira retangular, CIRCUITO MAGNETICO com
%  relutancia e entreferro, dispersao e coeficiente de acoplamento,
%  independencia da superficie e nucleo de secao variavel.
%  A INDUCAO propriamente dita tem banco proprio (Lenz); aqui a fem so
%  aparece onde o fluxo e o objeto de interesse.
%  Distribuicao mantida: 6 N1 + 10 N2 + 11 N3 + 7 N4 = 34
% =====================================================================

\subsubsection*{Banco complementar --- Fluxo magnético}

\begin{atencao}[Organização do banco]
O fluxo é $\Phi=\int\vec B\cdot\mathrm{d}\vec A$, que se reduz a
$BA\cos\theta$ apenas quando o campo é \textbf{uniforme} e a superfície
\textbf{plana} --- $\theta$ sendo o ângulo entre $\vec B$ e a \emph{normal}, e
não entre $\vec B$ e o plano. Duas propriedades estruturam o tópico: o fluxo por
qualquer superfície \emph{fechada} é nulo, e o fluxo por superfícies de mesma
\emph{borda} é o mesmo. O N3 trata de campos não uniformes, circuitos magnéticos
e dispersão; o N4 exige demonstrações e projeto.
\end{atencao}

\blocofixacao

\begin{questao}[1]{}
Uma superfície plana de $\SI{0.01}{\meter\squared}$ está em campo uniforme de
$\SI{0.1}{\tesla}$, com a normal formando $\ang{60}$ com o campo. Determine o
fluxo.
\resposta{$\Phi=\SI{0.5}{\milli\weber}$}
\resolucao{
\[
\Phi=BA\cos\theta=(0{,}1)(0{,}01)\cos\ang{60}
=(0{,}1)(0{,}01)(0{,}5)=\pot{5}{-4}\,\si{\weber}.
\]
\textbf{Atenção ao ângulo:} $\theta$ é medido entre $\vec B$ e a \textbf{normal}
à superfície. Se o enunciado der o ângulo com o \emph{plano}, use o complemento
--- é a troca de $\cos$ por $\sin$ que mais gera erro neste tópico.}
\end{questao}

\begin{questao}[1]{}
Determine o fluxo quando o campo é \textbf{paralelo} ao plano da superfície.
\resposta{$\Phi=0$}
\resolucao{Campo paralelo ao plano significa campo \emph{perpendicular} à
normal, logo $\theta=\ang{90}$ e $\cos\theta=0$:
\[
\Phi=0.
\]
\textbf{Interpretação geométrica:} nenhuma linha de campo \emph{atravessa} a
superfície --- elas passam rentes a ela.\\
\textbf{Casos-limite úteis:} $\Phi$ é máximo com $\vec B$ perpendicular ao plano
(paralelo à normal) e nulo com $\vec B$ contido no plano.}
\end{questao}

\begin{questao}[1]{}
Um fluxo de $\SI{2}{\milli\weber}$ atravessa perpendicularmente uma superfície
em campo de $\SI{0.4}{\tesla}$. Determine a área.
\resposta{$A=\SI{50}{\centi\meter\squared}$}
\resolucao{
\[
A=\frac{\Phi}{B}=\frac{\pot{2}{-3}}{0{,}4}=\pot{5}{-3}\,\si{\meter\squared}
=\SI{50}{\centi\meter\squared}.
\]
\textbf{Lembre da conversão:}
$\SI{1}{\meter\squared}=\SI{e4}{\centi\meter\squared}$, pois a área é o
\emph{quadrado} de um comprimento.}
\end{questao}

\begin{questao}[1]{}
Uma espira de $\SI{30}{\centi\meter\squared}$ é atravessada
perpendicularmente por fluxo de $\SI{1.5}{\milli\weber}$. Determine a densidade
de fluxo.
\resposta{$B=\SI{0.5}{\tesla}$}
\resolucao{
\[
B=\frac{\Phi}{A}=\frac{\pot{1.5}{-3}}{\pot{3}{-3}}=\SI{0.5}{\tesla}.
\]
\textbf{Nomenclatura:} $B$ é chamado \emph{densidade de fluxo} justamente por
ser fluxo por unidade de área --- e $\SI{1}{\tesla}=\SI{1}{\weber\per\meter\squared}$.\\
As duas grandezas descrevem a mesma realidade em escalas diferentes: $B$ é
local, $\Phi$ é integral.}
\end{questao}

\begin{questao}[1]{}
O fluxo magnético líquido através de qualquer superfície \textbf{fechada} é:
\begin{alternativas}[2]
\item proporcional à corrente envolvida
\item sempre nulo
\item proporcional ao volume
\item máximo quando a superfície é esférica
\end{alternativas}
\resposta{(b)}
\resolucao{As linhas de campo magnético são \textbf{fechadas} --- não têm início
nem fim. Toda linha que entra em uma superfície fechada necessariamente sai dela,
e as contribuições se cancelam:
\[
\oint\vec B\cdot\mathrm{d}\vec A=0.
\]
\textbf{Contraste com o campo elétrico:} lá, a lei de Gauss dá
$\oint\vec E\cdot\mathrm{d}\vec A=Q_{env}/\varepsilon_0$ --- diferente de zero
sempre que houver carga envolvida.\\
\textbf{A diferença tem uma causa única:} existem cargas elétricas isoladas
(monopolos elétricos), mas \textbf{não existem monopolos magnéticos}. Esta é uma
das equações de Maxwell, e sua forma diferencial é $\nabla\cdot\vec B=0$.}
\end{questao}

\begin{questao}[1]{}
Converta $\SI{5}{\milli\weber}$ para maxwells.
\dados{$\SI{1}{\weber}=10^{8}\,\si{\maxwell}$}
\resposta{$\pot{5}{5}\,\si{\maxwell}$}
\resolucao{
\[
\Phi=\pot{5}{-3}\times10^{8}=\pot{5}{5}\,\si{\maxwell}.
\]
\textbf{Origem da unidade:} o maxwell pertence ao sistema CGS, e vale
$\SI{1}{\gauss\centi\meter\squared}$. A conversão decorre diretamente de
$\SI{1}{\tesla}=\SI{e4}{\gauss}$ e
$\SI{1}{\meter\squared}=\SI{e4}{\centi\meter\squared}$:
\[
10^{4}\times10^{4}=10^{8}.
\]
Embora obsoleto no SI, o maxwell ainda aparece em catálogos de ímãs e em
literatura mais antiga.}
\end{questao}

\blocoaplicacao

\begin{questao}[2]{}
Uma bobina de $50$ espiras é atravessada por fluxo de
$\SI{1}{\milli\weber}$ por espira.
\begin{itensq}
\item Determine o fluxo concatenado.
\item Determine a indutância, se a corrente for $\SI{2}{\ampere}$.
\end{itensq}
\resposta{$\lambda=\SI{50}{\milli\weber}$-espira; $L=\SI{25}{\milli\henry}$}
\resolucao{(a)
\[
\lambda=N\Phi=50(\pot{1}{-3})=\pot{5}{-2}\,\si{\weber}\text{-espira}.
\]
(b)
\[
L=\frac{\lambda}{I}=\frac{0{,}05}{2}=\SI{25}{\milli\henry}.
\]
\textbf{Distinção essencial.} $\Phi$ é o fluxo através de \emph{uma} espira;
$\lambda=N\Phi$ é o \emph{fluxo concatenado}, que conta cada espira
separadamente. É $\lambda$, e não $\Phi$, que entra na definição de indutância e
na lei de Faraday para bobinas ($\varepsilon=-\mathrm{d}\lambda/\mathrm{d}t$).}
\end{questao}

\begin{questao}[2]{}
Uma espira de $\SI{200}{\centi\meter\squared}$ está em campo de
$\SI{0.25}{\tesla}$, com a normal a $\ang{60}$ do campo. Determine o fluxo e o
valor que ele teria na posição de fluxo máximo.
\resposta{$\Phi=\SI{2.5}{\milli\weber}$; máximo de $\SI{5}{\milli\weber}$}
\resolucao{
\[
\Phi=BA\cos\ang{60}=(0{,}25)(\pot{2}{-2})(0{,}5)=\pot{2.5}{-3}\,\si{\weber};
\]
\[
\Phi_{\max}=BA=\pot{5}{-3}\,\si{\weber}\quad(\theta=0).
\]
\textbf{Aplicação:} girar a espira entre essas duas posições varia o fluxo de
$\SI{2.5}{\milli\weber}$. Feito em $\SI{10}{\milli\second}$, isso induziria
$\varepsilon=\SI{0.25}{\volt}$ por espira --- princípio do gerador.}
\end{questao}

\begin{questao}[2]{}
Um solenoide de $\SI{1000}{}$ espiras por metro conduz $\SI{2}{\ampere}$ e tem
seção de $\SI{10}{\centi\meter\squared}$. Determine o fluxo por espira.
\resposta{$\Phi=\SI{2.51}{\micro\weber}$}
\resolucao{
\[
B=\mu_0nI=(4\pi\times10^{-7})(1000)(2)=\pot{2.51}{-3}\,\si{\tesla};
\]
\[
\Phi=BA=(\pot{2.51}{-3})(\pot{1}{-3})=\pot{2.51}{-6}\,\si{\weber}.
\]
\textbf{Ordem de grandeza:} apenas $\SI{2.5}{\micro\weber}$ --- fluxos em
núcleo de ar são pequenos. Com núcleo ferromagnético de
$\mu_r=1000$, o mesmo solenoide produziria $\SI{2.5}{\milli\weber}$.\\
\textbf{É por isso que máquinas elétricas usam núcleo:} sem ele, o fluxo (e
portanto a f.e.m. induzida e o torque) seria mil vezes menor.}
\end{questao}

\begin{questao}[2]{}
O fluxo através de uma bobina de $100$ espiras varia linearmente de zero a
$\SI{4}{\milli\weber}$ em $\SI{20}{\milli\second}$. Determine a f.e.m.
induzida.
\resposta{$\varepsilon=\SI{20}{\volt}$}
\resolucao{
\[
|\varepsilon|=N\frac{\Delta\Phi}{\Delta t}
=100\,\frac{\pot{4}{-3}}{\pot{2}{-2}}=\SI{20}{\volt}.
\]
\textbf{Como a variação é linear}, a f.e.m. é \emph{constante} durante todo o
intervalo. Se o fluxo variasse de outra forma, seria preciso derivar ponto a
ponto.\\
\textbf{Note o papel de $N$:} cem espiras multiplicam a f.e.m. por cem, embora o
fluxo \emph{por espira} seja o mesmo. É o fluxo concatenado que importa.}
\end{questao}

\begin{questao}[2]{}
Duas superfícies diferentes apoiam-se na \textbf{mesma} espira: uma plana e
outra em forma de calota. Compare o fluxo através delas.
\resposta{São iguais}
\resolucao{O fluxo através de qualquer superfície depende apenas de sua
\textbf{borda}.\\
\textbf{Demonstração:} as duas superfícies, juntas, formam uma superfície
\emph{fechada}. Como
$\oint\vec B\cdot\mathrm{d}\vec A=0$, o fluxo que entra por uma sai pela outra
--- e portanto os fluxos através delas são iguais (atentando para a orientação
das normais).\\
\textbf{Consequência prática:} na lei de Faraday, pode-se escolher
\emph{qualquer} superfície apoiada no circuito. Escolha a que tornar a
integral mais simples --- geralmente a plana, mas às vezes uma superfície
curva evita regiões de campo complicado.\\
\textbf{Contraste:} isso \emph{não} vale para o fluxo elétrico, que depende da
carga envolvida e portanto da superfície escolhida.}
\end{questao}

\begin{questao}[2]{}
Uma espira quadrada de lado $\SI{20}{\centi\meter}$ é girada de
$\ang{0}$ a $\ang{90}$ (normal em relação ao campo) em campo de
$\SI{0.6}{\tesla}$. Determine a variação de fluxo.
\resposta{$\Delta\Phi=\SI{-24}{\milli\weber}$}
\resolucao{
\[
\Phi_i=BA=(0{,}6)(0{,}04)=\SI{24}{\milli\weber};
\qquad
\Phi_f=BA\cos\ang{90}=0.
\]
\[
\Delta\Phi=0-0{,}024=\SI{-24}{\milli\weber}.
\]
\textbf{Duas formas de variar o fluxo aparecem no tópico:} mudar $B$, mudar
$A$, ou mudar a \emph{orientação}. Esta última é a base do gerador rotativo, em
que $\theta=\omega t$ e o fluxo varia senoidalmente.\\
\textbf{Sinal:} o fluxo diminuiu. Pela lei de Lenz, a corrente induzida
circularia no sentido de \emph{manter} o fluxo original.}
\end{questao}

\begin{questao}[2]{}
Um núcleo toroidal tem seção de $\SI{4}{\centi\meter\squared}$ e conduz fluxo
de $\SI{0.6}{\milli\weber}$. Determine a densidade de fluxo e avalie a
proximidade da saturação.
\resposta{$B=\SI{1.5}{\tesla}$ --- próximo da saturação típica}
\resolucao{
\[
B=\frac{\Phi}{A}=\frac{\pot{6}{-4}}{\pot{4}{-4}}=\SI{1.5}{\tesla}.
\]
\textbf{Avaliação:} aços-silício saturam entre $1{,}5$ e $\SI{2}{\tesla}$;
ferrites, bem antes (tipicamente $0{,}3$ a $\SI{0.5}{\tesla}$).\\
Portanto, para aço-silício este núcleo opera \textbf{no limite}. Qualquer
aumento de corrente levaria à saturação, com queda abrupta de $\mu_r$ e
distorção severa da forma de onda.\\
\textbf{Prática de projeto:} dimensione para $B_{\max}$ entre $60$ e
$80\%$ da saturação, deixando margem para transitórios e para a variação com a
temperatura.}
\end{questao}

\begin{questao}[2]{}
Uma bobina de $200$ espiras e área $\SI{15}{\centi\meter\squared}$ está em
campo de $\SI{0.4}{\tesla}$ perpendicular. Determine o fluxo concatenado.
\resposta{$\lambda=\SI{120}{\milli\weber}$-espira}
\resolucao{
\[
\Phi=BA=(0{,}4)(\pot{1.5}{-3})=\pot{6}{-4}\,\si{\weber};
\]
\[
\lambda=N\Phi=200(\pot{6}{-4})=\pot{0.12}{}\,\si{\weber}\text{-espira}.
\]
\textbf{Se essa bobina fosse retirada do campo em $\SI{50}{\milli\second}$:}
\[
|\varepsilon|=\frac{\Delta\lambda}{\Delta t}=\frac{0{,}12}{0{,}05}
=\SI{2.4}{\volt}.
\]
É esse o princípio do \emph{fluxímetro}: integrando a f.e.m. medida, recupera-se
$\Delta\lambda$ e, com $N$ e $A$ conhecidos, o campo.}
\end{questao}

\begin{questao}[2]{}
Uma espira está totalmente fora de uma região de campo uniforme e é empurrada
até ficar totalmente dentro. Descreva a variação do fluxo.
\resposta{Cresce de zero a $BA$, com taxa constante enquanto a espira atravessa
a fronteira}
\resolucao{\textbf{Três fases:}
\begin{itemize}
\item \textbf{Totalmente fora:} $\Phi=0$ e $\mathrm{d}\Phi/\mathrm{d}t=0$ ---
nenhuma f.e.m.
\item \textbf{Atravessando a fronteira:} a área \emph{dentro} do campo cresce
linearmente com o deslocamento. Se a velocidade for constante,
$\Phi$ cresce linearmente e a f.e.m. é \textbf{constante}.
\item \textbf{Totalmente dentro:} $\Phi=BA$ e volta a ser constante ---
nenhuma f.e.m., mesmo que a espira continue se movendo.
\end{itemize}
\textbf{Ponto conceitual importante:} não é o \emph{movimento} que induz f.e.m.,
e sim a \textbf{variação de fluxo}. Uma espira movendo-se rapidamente dentro de
um campo uniforme não gera nada.\\
Com largura $\ell$ e velocidade $v$, a f.e.m. na fase intermediária vale
$\varepsilon=B\ell v$.}
\end{questao}

\begin{questao}[2]{}
Explique a diferença entre fluxo magnético e densidade de fluxo, e indique
quando cada um é a grandeza adequada.
\resposta{$\Phi$ é integral (weber); $B$ é local (tesla). $\Phi$ governa a
indução; $B$ governa forças e saturação}
\resolucao{\textbf{Densidade de fluxo $\vec B$:} grandeza \emph{local} e
vetorial, medida em tesla. Ela é o que determina:
\begin{itemize}
\item a força sobre cargas e condutores ($\vec F=q\vec v\times\vec B$);
\item a saturação do material ($B$ próximo de $B_{sat}$);
\item o que um sensor Hall mede.
\end{itemize}
\textbf{Fluxo $\Phi$:} grandeza \emph{integral} e escalar, medida em weber.
Ela é o que determina:
\begin{itemize}
\item a f.e.m. induzida ($\varepsilon=-\mathrm{d}\Phi/\mathrm{d}t$);
\item a indutância ($L=\lambda/I$);
\item o acoplamento entre bobinas.
\end{itemize}
\textbf{Relação:} $\Phi=\int\vec B\cdot\mathrm{d}\vec A$, e no caso uniforme
$B=\Phi/A$.\\
\textbf{Erro comum:} falar em "fluxo" quando se quer dizer "campo". Um ímã
pequeno e um grande podem ter o mesmo $B$ na superfície e fluxos totais muito
diferentes --- e é o fluxo que determina o desempenho em um circuito magnético.}
\end{questao}

\blocoaprofundamento

\begin{questao}[3]{}
Um campo não uniforme $B(x)=B_0\dfrac{x}{L}$, perpendicular ao plano, atravessa
uma superfície retangular de largura $L$ (na direção $x$) e altura $h$, com
$B_0=\SI{0.5}{\tesla}$, $L=\SI{20}{\centi\meter}$ e
$h=\SI{10}{\centi\meter}$.
\begin{itensq}
\item Determine o fluxo por integração.
\item Compare com as estimativas usando $B_0$ e usando o valor médio.
\end{itensq}
\resposta{$\Phi=\SI{5}{\milli\weber}$ --- igual à estimativa pelo valor médio}
\resolucao{(a) Dividindo em faixas verticais de largura $\mathrm{d}x$ e área
$h\,\mathrm{d}x$:
\[
\Phi=\int_0^{L}B_0\frac{x}{L}\,h\,\mathrm{d}x
=\frac{B_0h}{L}\left[\frac{x^{2}}{2}\right]_0^{L}
=\frac{B_0hL}{2}.
\]
\[
\Phi=\frac{(0{,}5)(0{,}10)(0{,}20)}{2}=\pot{5}{-3}\,\si{\weber}.
\]
(b)
\begin{center}
\begin{tabular}{lc}
\hline
Estimativa & $\Phi$\\
\hline
Usando $B_0$ (valor máximo) & $\SI{10}{\milli\weber}$\\
\textbf{Integração exata} & $\SI{5}{\milli\weber}$\\
Usando $\bar B=B_0/2$ & $\SI{5}{\milli\weber}$\\
\hline
\end{tabular}
\end{center}
\textbf{A média coincide com o exato} --- mas apenas porque $B$ varia
\emph{linearmente}. Para $B\propto x^{2}$, a integração daria $B_0hL/3$,
enquanto a média aritmética dos extremos daria $B_0hL/2$ --- erro de
$50\%$.\\
\textbf{Regra:} substituir o campo pela média só é legítimo em variação linear.
Fora disso, integre.}
\end{questao}

\begin{questao}[3]{}
Um fio retilíneo longo conduz $\SI{10}{\ampere}$. Uma espira retangular de
$\SI{20}{\centi\meter}$ de comprimento (paralelo ao fio) está no mesmo plano,
com os lados a $\SI{5}{\centi\meter}$ e $\SI{20}{\centi\meter}$ do fio.
\begin{itensq}
\item Deduza a expressão do fluxo.
\item Calcule seu valor.
\end{itensq}
\resposta{$\Phi=\dfrac{\mu_0Ia}{2\pi}\ln\dfrac{d+b}{d}=\SI{0.55}{\micro\weber}$}
\resolucao{(a) O campo do fio \textbf{não é uniforme} na região da espira:
$B(r)=\dfrac{\mu_0I}{2\pi r}$. Dividindo a espira em faixas de largura
$\mathrm{d}r$ e comprimento $a$:
\[
\Phi=\int_d^{d+b}\frac{\mu_0I}{2\pi r}\,a\,\mathrm{d}r
=\frac{\mu_0Ia}{2\pi}\ln\!\left(\frac{d+b}{d}\right).
\]
(b) Com $a=\SI{0.20}{\meter}$, $d=\SI{0.05}{\meter}$ e
$b=\SI{0.15}{\meter}$ (borda externa em $\SI{0.20}{\meter}$):
\[
\Phi=\pot{2}{-7}(10)(0{,}20)\ln 4
=(\pot{4}{-7})(1{,}386)=\pot{5.55}{-7}\,\si{\weber}.
\]
\textbf{Observe o logaritmo:} afastar a espira reduz o fluxo muito devagar.
Dobrar $d$ e $b$ juntos (mantendo a razão) \emph{não muda nada} --- o fluxo
depende apenas da razão $(d+b)/d$ e do comprimento $a$.\\
\textbf{Consequência:} a indutância mútua entre um fio e uma espira próxima é
notavelmente insensível à distância absoluta. É por isso que interferência
magnética de cabos de potência é difícil de eliminar apenas afastando-os.}
\end{questao}

\begin{questao}[3]{}
Um núcleo magnético de ferro tem comprimento médio $\SI{50}{\centi\meter}$,
seção de $\SI{4}{\centi\meter\squared}$ e $\mu_r=2000$, com enrolamento de
$\SI{1000}{}$ ampère-espiras.
\begin{itensq}
\item Determine a relutância do núcleo.
\item Determine o fluxo e a densidade de fluxo.
\item Avalie o resultado.
\end{itensq}
\resposta{$\mathcal{R}=\pot{4.97}{5}\,\si{\per\henry}$;
$\Phi=\SI{2}{\milli\weber}$; $B=\SI{5}{\tesla}$ --- \textbf{impossível}}
\resolucao{(a) Em analogia ao circuito elétrico, a \textbf{relutância} vale
\[
\mathcal{R}=\frac{\ell}{\mu_r\mu_0A}
=\frac{0{,}50}{(2000)(4\pi\times10^{-7})(\pot{4}{-4})}
=\pot{4.97}{5}\,\si{\per\henry}.
\]
(b) Com a força magnetomotriz $\mathcal{F}=NI=\SI{1000}{}$ ampère-espiras:
\[
\Phi=\frac{\mathcal{F}}{\mathcal{R}}=\frac{1000}{\pot{4.97}{5}}
=\pot{2.01}{-3}\,\si{\weber};
\qquad
B=\frac{\Phi}{A}=\SI{5.03}{\tesla}.
\]
(c) \textbf{O resultado é impossível.} Nenhum material ferromagnético comum
sustenta $\SI{5}{\tesla}$ --- todos saturam entre $1{,}5$ e
$\SI{2}{\tesla}$.\\
\textbf{O que realmente aconteceria:} conforme $B$ se aproxima de
$\SI{2}{\tesla}$, $\mu_r$ despenca de $2000$ para valores próximos de $1$. A
relutância dispara, e o fluxo estaciona em torno de
$B_{sat}A=\SI{0.8}{\milli\weber}$.\\
\textbf{Lição de método:} o modelo linear com $\mu_r$ constante é válido
\emph{apenas} abaixo da saturação. Calcular e depois \textbf{verificar $B$}
contra $B_{sat}$ é etapa obrigatória --- exatamente como se verifica a rigidez
dielétrica em problemas de campo elétrico.}
\end{questao}

\begin{questao}[3]{}
O mesmo núcleo da questão anterior recebe um \textbf{entreferro} de
$\SI{1}{\milli\meter}$.
\begin{itensq}
\item Determine a relutância do entreferro e a total.
\item Determine o novo fluxo e $B$.
\item Explique por que o entreferro domina.
\end{itensq}
\resposta{$\mathcal{R}_g=\pot{1.99}{6}$, quatro vezes a do ferro;
$B=\SI{1.0}{\tesla}$}
\resolucao{(a) No entreferro, $\mu_r=1$:
\[
\mathcal{R}_g=\frac{g}{\mu_0A}
=\frac{\pot{1}{-3}}{(4\pi\times10^{-7})(\pot{4}{-4})}
=\pot{1.99}{6}\,\si{\per\henry}.
\]
Em série com o ferro:
$\mathcal{R}_{tot}=\pot{4.97}{5}+\pot{1.99}{6}=\pot{2.49}{6}$.\\
(b)
\[
\Phi=\frac{1000}{\pot{2.49}{6}}=\pot{4.02}{-4}\,\si{\weber};
\qquad
B=\SI{1.005}{\tesla}.
\]
\textbf{Agora sim, valor fisicamente sustentável} --- e com margem para
saturação.\\
(c) \textbf{Por que o entreferro domina.} Ele tem apenas
$\SI{1}{\milli\meter}$ contra $\SI{500}{\milli\meter}$ de ferro --- é
$500$ vezes mais curto. Mas sua permeabilidade é $2000$ vezes menor. O saldo:
\[
\frac{\mathcal{R}_g}{\mathcal{R}_{fe}}=\frac{2000}{500}=4.
\]
\textbf{Um milímetro de ar vale quatro vezes meio metro de ferro.}\\
\textbf{Consequências de projeto.}
\begin{itemize}
\item O entreferro \emph{lineariza} o circuito magnético: como ele domina, a
relação $\Phi\times I$ torna-se quase reta, mesmo com o ferro não linear.
\item Ele permite armazenar muito mais energia (que se concentra no ar) --- daí
seu uso em indutores de conversores chaveados.
\item Em contrapartida, reduz a indutância e exige mais ampère-espiras.
\item \textbf{Entreferros parasitas} (juntas mal ajustadas de núcleos em E ou U)
degradam severamente o desempenho, mesmo com dezenas de micrômetros.
\end{itemize}}
\end{questao}

\begin{questao}[3]{}
Duas bobinas compartilham parte do fluxo. A bobina $1$ produz
$\SI{5}{\milli\weber}$, dos quais $\SI{4}{\milli\weber}$ atravessam a bobina
$2$.
\begin{itensq}
\item Determine o fluxo de dispersão.
\item Determine o coeficiente de acoplamento aproximado.
\item Discuta como aumentá-lo.
\end{itensq}
\resposta{$\SI{1}{\milli\weber}$ de dispersão; $k\approx0{,}8$}
\resolucao{(a) \textbf{Dispersão:} a parcela que não atravessa a outra bobina:
\[
\Phi_{disp}=5-4=\SI{1}{\milli\weber}\quad(20\%).
\]
(b) O coeficiente de acoplamento é a fração de fluxo compartilhada:
\[
k\approx\frac{\Phi_{12}}{\Phi_1}=\frac{4}{5}=0{,}8.
\]
\emph{(A definição rigorosa é $k=M/\sqrt{L_1L_2}$, que coincide com esta razão
quando as bobinas são simétricas.)}\\
(c) \textbf{Como aumentar $k$.}
\begin{itemize}
\item \textbf{Núcleo fechado} de alta permeabilidade: o fluxo prefere o caminho
de baixa relutância e "não escapa". Toroides atingem $k>0{,}99$.
\item \textbf{Enrolamentos concêntricos ou intercalados}, em vez de lado a lado
--- reduz drasticamente o caminho de fuga.
\item \textbf{Eliminar entreferros} desnecessários, que forçam o fluxo a
divergir.
\end{itemize}
\textbf{Compromisso:} $k$ alto é desejável em transformadores (transferência
eficiente), mas dispersão \emph{controlada} é útil em alguns casos --- ela
funciona como indutância série natural, limitando corrente de curto em
transformadores de solda e em reatores.}
\end{questao}

\begin{questao}[3]{}
Um núcleo tem seção variável: $\SI{6}{\centi\meter\squared}$ em um trecho e
$\SI{2}{\centi\meter\squared}$ em outro, sem dispersão.
\begin{itensq}
\item Compare o fluxo nos dois trechos.
\item Compare a densidade de fluxo.
\item Identifique onde ocorre a saturação primeiro.
\end{itensq}
\resposta{Mesmo fluxo; $B$ três vezes maior na seção estreita, que satura
primeiro}
\resolucao{(a) \textbf{O fluxo é o mesmo.} Sem dispersão, todas as linhas que
passam por um trecho passam pelo outro --- é a consequência direta de
$\oint\vec B\cdot\mathrm{d}\vec A=0$ aplicada a uma superfície que envolve a
junção.\\
\textbf{Analogia elétrica:} é o equivalente magnético da LKC. O fluxo se
conserva ao longo de um circuito magnético sem ramificação, como a corrente em
um circuito série.\\
(b) Como $B=\Phi/A$:
\[
\frac{B_2}{B_1}=\frac{A_1}{A_2}=\frac{6}{2}=3.
\]
(c) \textbf{A seção estreita satura primeiro}, por ter $B$ três vezes maior.\\
\textbf{Consequência de projeto:} o dimensionamento de um núcleo é governado
pela sua \textbf{menor} seção. Aumentar a seção larga não adianta nada ---
é preciso alargar o gargalo.\\
\textbf{Onde isso aparece na prática:} cantos de núcleos em E, junções mal
projetadas e regiões com furos de fixação. São esses pontos que saturam
primeiro e produzem aquecimento e distorção localizados.}
\end{questao}

\begin{questao}[3]{}
Uma espira de $\SI{25}{\centi\meter\squared}$ está parcialmente dentro de uma
região de campo de $\SI{0.8}{\tesla}$, com $60\%$ de sua área na região.
\begin{itensq}
\item Determine o fluxo.
\item Determine a f.e.m. se a espira for retirada completamente em
      $\SI{40}{\milli\second}$.
\end{itensq}
\resposta{$\Phi=\SI{1.2}{\milli\weber}$; $\varepsilon=\SI{30}{\milli\volt}$}
\resolucao{(a) Apenas a área \emph{dentro} do campo contribui:
\[
\Phi=B\,A_{efetiva}=(0{,}8)(0{,}60)(\pot{2.5}{-3})
=\pot{1.2}{-3}\,\si{\weber}.
\]
(b)
\[
|\varepsilon|=\frac{\Delta\Phi}{\Delta t}
=\frac{\pot{1.2}{-3}}{\pot{4}{-2}}=\pot{3}{-2}\,\si{\volt}.
\]
\textbf{Observação importante:} o fluxo depende da área \emph{efetivamente}
imersa no campo, não da área geométrica total da espira. Fora da região,
$B=0$ e a contribuição é nula.\\
\textbf{Se a espira tivesse $N$ espiras}, a f.e.m. seria $N$ vezes maior --- e é
assim que se aumenta a sensibilidade de sensores de fluxo sem aumentar a área.}
\end{questao}

\begin{questao}[3]{}
Um transformador tem núcleo de seção $\SI{10}{\centi\meter\squared}$ e opera
com $B_{\max}=\SI{1.2}{\tesla}$.
\begin{itensq}
\item Determine o fluxo máximo.
\item Determine o número de espiras do primário para
      $\SI{220}{\volt}$ a $\SI{60}{\hertz}$.
\end{itensq}
\resposta{$\Phi_{\max}=\SI{1.2}{\milli\weber}$; $N\approx688$ espiras}
\resolucao{(a) $\Phi_{\max}=B_{\max}A=(1{,}2)(\pot{1}{-3})
=\pot{1.2}{-3}\,\si{\weber}$.\\
(b) Para fluxo senoidal, a f.e.m. eficaz vale
\[
V_{ef}=4{,}44\,N\,f\,\Phi_{\max},
\]
\emph{(o fator $4{,}44=2\pi/\sqrt2$ decorre de derivar uma senoide e converter
para valor eficaz)}. Logo
\[
N=\frac{220}{4{,}44(60)(\pot{1.2}{-3})}=\frac{220}{0{,}3197}\approx688.
\]
\textbf{Compromisso central do projeto:} menos espiras significam $B$ maior (e
risco de saturação); mais espiras significam mais cobre, mais resistência e mais
perdas.\\
\textbf{Por que a frequência importa:} $N\propto1/f$. Em $\SI{400}{\hertz}$
(aviação) ou em dezenas de quilohertz (fontes chaveadas), o mesmo transformador
precisa de muito menos espiras --- e o núcleo pode ser muito menor. É essa a
razão física de fontes chaveadas serem tão mais compactas que as lineares.}
\end{questao}

\begin{questao}[3]{}
Uma superfície fechada envolve completamente um ímã permanente. Determine o
fluxo magnético líquido através dela e justifique.
\resposta{Nulo, independentemente da posição e da intensidade do ímã}
\resolucao{\textbf{Resultado:} $\oint\vec B\cdot\mathrm{d}\vec A=0$, sempre.\\
\textbf{Justificativa:} as linhas de campo de um ímã \emph{não começam nem
terminam} nele --- elas saem do polo norte, percorrem o espaço externo, entram
pelo sul e \textbf{continuam dentro do material}, fechando-se sobre si mesmas.\\
Toda linha que atravessa a superfície para fora atravessa também para dentro,
em algum outro ponto. O balanço é rigorosamente nulo.\\
\textbf{Contraste com a carga elétrica.} Se a superfície envolvesse uma carga
$+Q$, o fluxo elétrico seria $Q/\varepsilon_0\ne0$ --- porque as linhas
elétricas \emph{nascem} na carga.\\
\textbf{O que isso significa fisicamente:} não existe "carga magnética". Cortar
um ímã ao meio não separa os polos --- produz dois ímãs completos, cada um com
seus dois polos. Não há como isolar um norte de um sul.\\
\textbf{Consequência para a teoria:} $\nabla\cdot\vec B=0$ é uma equação de
Maxwell, e é ela que garante a existência do potencial vetor $\vec A$ com
$\vec B=\nabla\times\vec A$.}
\end{questao}

\begin{questao}[3]{}
Uma bobina de $500$ espiras e área $\SI{8}{\centi\meter\squared}$ é usada como
sensor de campo. Ela é retirada rapidamente do campo, e um integrador registra
$\SI{6}{\milli\volt\second}$.
\begin{itensq}
\item Determine o fluxo concatenado inicial.
\item Determine o campo.
\item Explique a vantagem do método.
\end{itensq}
\resposta{$\lambda=\SI{6}{\milli\weber}$-espira; $B=\SI{15}{\milli\tesla}$}
\resolucao{(a) Integrando a lei de Faraday:
\[
\int\varepsilon\,\mathrm{d}t=\Delta\lambda
\;\Longrightarrow\;
\lambda_i=\pot{6}{-3}\,\si{\weber}\text{-espira}
\]
(pois o fluxo final é zero).\\
(b)
\[
\Phi=\frac{\lambda}{N}=\frac{\pot{6}{-3}}{500}=\pot{1.2}{-5}\,\si{\weber};
\qquad
B=\frac{\Phi}{A}=\frac{\pot{1.2}{-5}}{\pot{8}{-4}}=\pot{1.5}{-2}\,\si{\tesla}.
\]
(c) \textbf{Vantagem do método integral.} O resultado depende \emph{apenas} da
variação total de fluxo --- \textbf{não} de \emph{como} a bobina foi retirada.
Rápido ou devagar, com solavancos ou suavemente, a integral é a mesma.\\
Isso torna a medição notavelmente robusta: não é preciso controlar a velocidade
nem a trajetória.\\
\textbf{Limitações:} o integrador acumula qualquer \emph{offset} do
amplificador, o que limita o tempo de medição (problema já discutido na
integração de sinais ruidosos). E o método mede apenas \emph{variações} --- para
campos estáticos permanentes, é preciso mover a bobina ou usar sensor Hall.}
\end{questao}

\begin{questao}[3]{}
Avalie o erro de assumir campo uniforme ao calcular o fluxo através de uma
espira de $\SI{10}{\centi\meter}$ de lado, situada entre $\SI{5}{}$ e
$\SI{15}{\centi\meter}$ de um fio retilíneo.
\begin{itensq}
\item Determine o fluxo exato e o estimado pelo campo no centro.
\item Determine o erro.
\end{itensq}
\resposta{Erro de cerca de $+2{,}5\%$ na estimativa pelo centro}
\resolucao{\textbf{Exato} (com $I$ genérico e $a=\SI{0.10}{\meter}$):
\[
\Phi_{ex}=\frac{\mu_0Ia}{2\pi}\ln\frac{0{,}15}{0{,}05}
=\pot{2}{-8}I\ln3=\pot{2.197}{-8}\,I.
\]
\textbf{Estimativa pelo centro} ($r=\SI{0.10}{\meter}$):
\[
\Phi_{est}=\frac{\mu_0I}{2\pi(0{,}10)}\,a^{2}
=\pot{2}{-7}\frac{I}{0{,}10}(0{,}01)=\pot{2}{-8}\,I.
\]
\[
\text{erro}=\frac{2{,}00-2{,}197}{2{,}197}=-9{,}0\%.
\]
\textbf{A estimativa pelo centro \emph{subestima}} o fluxo --- porque
$B\propto1/r$ é uma função \textbf{convexa}, e a média de uma função convexa é
maior que a função da média (desigualdade de Jensen).\\
\textbf{Quando a aproximação é aceitável:} quando a variação de $B$ ao longo da
espira for pequena. Aqui, $B$ varia por um fator $3$ --- muito. Se a espira
estivesse entre $95$ e $\SI{105}{\centi\meter}$, a variação seria de apenas
$10\%$ e o erro cairia para frações de por cento.\\
\textbf{Critério prático:} use campo uniforme se
$B_{\max}/B_{\min}<1{,}1$; caso contrário, integre.}
\end{questao}

\blocodesafio

\begin{questao}[4]{}
\textbf{(Demonstração --- fluxo por superfície fechada.)} Prove que
$\oint\vec B\cdot\mathrm{d}\vec A=0$ e explore suas consequências.
\begin{itensq}
\item Demonstre a partir da inexistência de monopolos.
\item Deduza a conservação do fluxo em circuitos magnéticos.
\item Compare com a lei de Gauss elétrica.
\end{itensq}
\resposta{Linhas fechadas $\Rightarrow$ fluxo líquido nulo $\Rightarrow$ o fluxo
se conserva ao longo de um circuito magnético}
\resolucao{(a) \textbf{Demonstração.} As linhas de $\vec B$ são curvas
\emph{fechadas} --- fato experimental, equivalente à inexistência de monopolos
magnéticos. Uma curva fechada que penetre uma superfície fechada precisa
atravessá-la um número \emph{par} de vezes, alternando entrada e saída.\\
Cada entrada contribui com fluxo negativo e cada saída com positivo, de módulos
iguais. A soma se cancela:
\[
\oint\vec B\cdot\mathrm{d}\vec A=0.
\]
Pelo teorema da divergência, isso equivale a $\nabla\cdot\vec B=0$ em todo
ponto.\\
(b) \textbf{Conservação do fluxo.} Considere uma junção em um núcleo magnético,
onde um trecho de seção $A_1$ encontra outro de seção $A_2$. Envolvendo a junção
com uma superfície fechada:
\[
-\Phi_1+\Phi_2=0
\;\Longrightarrow\;
\Phi_1=\Phi_2.
\]
\textbf{O fluxo se conserva} --- exatamente como a corrente em um nó. Se o
núcleo se ramificar em dois caminhos:
\[
\Phi_{total}=\Phi_a+\Phi_b,
\]
que é a \textbf{lei dos nós magnética}, análoga à LKC.\\
(c) \textbf{Comparação com Gauss elétrica.}
\begin{center}
\begin{tabular}{ll}
\hline
Elétrico & Magnético\\
\hline
$\oint\vec E\cdot\mathrm{d}\vec A=Q/\varepsilon_0$
 & $\oint\vec B\cdot\mathrm{d}\vec A=0$\\
Linhas nascem e morrem em cargas & Linhas são fechadas\\
Existem monopolos (cargas) & Não existem monopolos\\
$\nabla\cdot\vec E=\rho/\varepsilon_0$ & $\nabla\cdot\vec B=0$\\
\hline
\end{tabular}
\end{center}
\textbf{Consequência profunda:} a ausência de monopolos permite escrever
$\vec B=\nabla\times\vec A$, base de toda a formulação moderna do
eletromagnetismo. Se monopolos existissem, essa construção seria impossível.}
\end{questao}

\begin{questao}[4]{}
\textbf{(Independência da superfície.)} Prove que o fluxo através de qualquer
superfície apoiada em uma dada curva fechada é o mesmo.
\begin{itensq}
\item Demonstre.
\item Explique a importância para a lei de Faraday.
\item Indique o que ocorre se as normais forem orientadas incoerentemente.
\end{itensq}
\resposta{Duas superfícies de mesma borda formam superfície fechada, cujo fluxo
é nulo}
\resolucao{(a) \textbf{Demonstração.} Sejam $S_1$ e $S_2$ duas superfícies
apoiadas na mesma curva $C$. Juntas, elas delimitam um \textbf{volume fechado}.
Aplicando o resultado anterior:
\[
\oint\vec B\cdot\mathrm{d}\vec A
=\Phi_{S_2}^{(saindo)}-\Phi_{S_1}^{(saindo)}=0
\;\Longrightarrow\;
\Phi_{S_1}=\Phi_{S_2}.
\]
(b) \textbf{Importância para Faraday.} A lei
$\varepsilon=-\dfrac{\mathrm{d}\Phi}{\mathrm{d}t}$ só faz sentido se
$\Phi$ estiver \emph{bem definido} --- ou seja, se não depender da superfície
escolhida.\\
Se dependesse, a f.e.m. de um mesmo circuito teria valores diferentes conforme a
superfície imaginada --- absurdo, pois a f.e.m. é medível.\\
\textbf{Liberdade prática:} escolha a superfície que simplificar a integral. Ao
calcular o fluxo através de uma espira que envolve um solenoide, por exemplo,
convém uma superfície que corte o solenoide perpendicularmente.\\
(c) \textbf{Orientação das normais.} A igualdade exige que as normais sejam
orientadas \emph{coerentemente} com o sentido de percurso da borda, pela regra
da mão direita.\\
Se uma delas for invertida, aparece um sinal negativo --- e a conclusão errônea
de que os fluxos são opostos. \textbf{Fixe o sentido de percurso da curva
primeiro}, e derive as normais dele.}
\end{questao}

\begin{questao}[4]{}
\textbf{(Circuito magnético.)} Estabeleça a analogia formal entre circuitos
magnéticos e elétricos, e discuta seus limites.
\begin{itensq}
\item Construa a tabela de correspondências.
\item Deduza a expressão da relutância.
\item Aponte três diferenças que quebram a analogia.
\end{itensq}
\resposta{$\mathcal{F}=\Phi\mathcal{R}$, com
$\mathcal{R}=\ell/\mu A$; a analogia falha por não linearidade, dispersão e
ausência de dissipação}
\resolucao{(a)
\begin{center}
\begin{tabular}{ll}
\hline
Circuito elétrico & Circuito magnético\\
\hline
F.e.m. $\varepsilon$ & Força magnetomotriz $\mathcal{F}=NI$\\
Corrente $I$ & Fluxo $\Phi$\\
Resistência $R=\ell/\sigma A$ & Relutância
 $\mathcal{R}=\ell/\mu A$\\
$\varepsilon=RI$ & $\mathcal{F}=\mathcal{R}\Phi$\\
Condutividade $\sigma$ & Permeabilidade $\mu$\\
LKC nos nós & Conservação do fluxo\\
Resistências em série somam & Relutâncias em série somam\\
\hline
\end{tabular}
\end{center}
(b) \textbf{Dedução da relutância.} Da lei de Ampère aplicada ao caminho médio
do núcleo, $NI=H\ell$. Com $B=\mu H$ e $\Phi=BA$:
\[
NI=\frac{B\ell}{\mu}=\frac{\Phi\,\ell}{\mu A}
\;\Longrightarrow\;
\mathcal{R}=\frac{\ell}{\mu A}.
\]
(c) \textbf{Três diferenças que quebram a analogia.}
\begin{enumerate}
\item \textbf{Não linearidade.} $\sigma$ é constante em metais, mas
$\mu$ varia enormemente com $B$ --- e desaba na saturação. A "resistência
magnética" depende da própria corrente que a atravessa.
\item \textbf{Dispersão.} A corrente elétrica fica confinada no condutor porque
a razão de condutividades metal/ar é de $10^{20}$. Já a razão de
permeabilidades ferro/ar é de apenas $10^{3}$ --- e por isso parte do fluxo
\emph{escapa} do núcleo. Não existe "isolante magnético".
\item \textbf{Ausência de dissipação.} A corrente dissipa $RI^{2}$ em calor.
O fluxo \emph{não dissipa} nada --- manter $\Phi$ constante não custa energia
(as perdas por histerese e Foucault ocorrem apenas quando o fluxo \emph{varia},
e por outros mecanismos).
\end{enumerate}
\textbf{Conclusão:} a analogia é excelente como \emph{ferramenta de cálculo}
abaixo da saturação, e enganosa como \emph{modelo físico}. Use-a para
dimensionar; não a use para raciocinar sobre energia.}
\end{questao}

\begin{questao}[4]{}
\textbf{(Projeto de entreferro.)} Um indutor de $\SI{100}{\micro\henry}$ deve
conduzir $\SI{5}{\ampere}$ de pico, com núcleo de ferrite de seção
$\SI{1}{\centi\meter\squared}$, caminho médio $\SI{6}{\centi\meter}$,
$\mu_r=2000$ e $B_{sat}=\SI{0.3}{\tesla}$.
\begin{itensq}
\item Determine a energia armazenada e o número mínimo de espiras.
\item Determine o entreferro necessário.
\item Determine a fração de energia armazenada no entreferro.
\end{itensq}
\resposta{$W=\SI{1.25}{\milli\joule}$; $N=17$ espiras;
$g\approx\SI{0.33}{\milli\meter}$; $92\%$ da energia no entreferro}
\resolucao{(a)
\[
W=\tfrac12LI^{2}=\tfrac12(\pot{1}{-4})(25)=\SI{1.25}{\milli\joule}.
\]
O fluxo concatenado é $\lambda=LI=\pot{5}{-4}\,\si{\weber}$-espira, e o fluxo
máximo por espira é limitado pela saturação:
\[
\Phi_{\max}=B_{sat}A=(0{,}3)(\pot{1}{-4})=\SI{30}{\micro\weber}.
\]
\[
N\ge\frac{\lambda}{\Phi_{\max}}=\frac{\pot{5}{-4}}{\pot{3}{-5}}=16{,}7
\;\Longrightarrow\;
N=17\ \text{espiras}.
\]
\emph{(Verificação: com $N=17$, $B=\lambda/(NA)=\SI{0.294}{\tesla}$ ---
logo abaixo do limite.)}\\
(b) De $L=N^{2}/\mathcal{R}$:
\[
\mathcal{R}_{tot}=\frac{17^{2}}{\pot{1}{-4}}=\pot{2.89}{6}\,\si{\per\henry}.
\]
A parcela do ferrite é
$\mathcal{R}_{fe}=\dfrac{0{,}06}{(2000)\mu_0(\pot{1}{-4})}
=\pot{2.39}{5}$, restando para o entreferro
$\mathcal{R}_g=\pot{2.65}{6}$:
\[
g=\mathcal{R}_g\,\mu_0A=(\pot{2.65}{6})(4\pi\times10^{-7})(\pot{1}{-4})
=\SI{0.33}{\milli\meter}.
\]
\textbf{Note que o entreferro responde por $92\%$ da relutância total}, embora
seja $180$ vezes mais curto que o caminho de ferrite.\\
(c) Como $B$ é o mesmo nos dois meios e $u=B^{2}/2\mu$:
\[
\frac{W_g}{W_{fe}}=\frac{g/\mu_0}{\ell_{fe}/(\mu_r\mu_0)}
=\frac{g\,\mu_r}{\ell_{fe}}
=\frac{(\pot{3.33}{-4})(2000)}{0{,}06}=11{,}1,
\]
ou seja, $\SI{91.7}{\percent}$ da energia está no \textbf{entreferro}.\\
\textbf{Por que isso importa.} A densidade de energia
$u=B^{2}/2\mu$ é \emph{inversamente} proporcional a $\mu$ --- para o mesmo
$B$, o ar armazena duas mil vezes mais energia por volume que o ferrite.\\
\textbf{O entreferro é o reservatório; o núcleo apenas conduz o fluxo até
ele.} Um indutor sem entreferro satura com corrente mínima e armazena quase
nada --- ao contrário de um transformador, que \emph{não deve} armazenar
energia e por isso dispensa entreferro.}
\end{questao}

\begin{questao}[4]{}
\textbf{(Dispersão e acoplamento.)} Analise o efeito da dispersão de fluxo no
desempenho de um transformador.
\begin{itensq}
\item Relacione $k$ com as indutâncias própria e mútua.
\item Determine o efeito sobre a tensão de saída.
\item Discuta quando a dispersão é desejável.
\end{itensq}
\resposta{$k=M/\sqrt{L_1L_2}$; a dispersão aparece como indutância série e
causa queda de tensão com a carga}
\resolucao{(a) Por definição,
\[
k=\frac{M}{\sqrt{L_1L_2}},
\qquad 0\le k\le1,
\]
resultado já demonstrado a partir da não negatividade da energia do par de
bobinas.\\
Decompondo cada indutância em parcela \emph{magnetizante} (fluxo compartilhado)
e parcela de \emph{dispersão}:
\[
L_1=L_{m1}+L_{d1},
\qquad
k^{2}=\frac{L_{m1}L_{m2}}{L_1L_2}.
\]
(b) \textbf{Efeito na saída.} A indutância de dispersão comporta-se como uma
\emph{impedância em série} com o enrolamento. Sob carga, ela produz queda de
tensão proporcional à corrente:
\begin{itemize}
\item \textbf{Em vazio}, a relação de tensões é praticamente $N_2/N_1$ --- a
dispersão não se manifesta.
\item \textbf{Com carga}, a tensão de saída \emph{cai}, e a regulação piora
proporcionalmente à dispersão.
\end{itemize}
Um transformador com $k=0{,}98$ tem cerca de $2\%$ de dispersão e regulação
tipicamente na faixa de $2$ a $5\%$.\\
(c) \textbf{Quando a dispersão é desejável.}
\begin{itemize}
\item \textbf{Transformadores de solda:} a dispersão elevada (às vezes
$k\approx0{,}8$) limita naturalmente a corrente de arco, dispensando reator
externo. A característica "mole" é justamente o que se quer.
\item \textbf{Transformadores de forno e de descarga:} mesma lógica --- a carga
é essencialmente um curto controlado.
\item \textbf{Limitação de curto-circuito:} em transformadores de distribuição,
uma dispersão mínima é \emph{exigida} por norma, para limitar a corrente de
falta a valores que a proteção consiga interromper.
\end{itemize}
\textbf{Conclusão:} $k\to1$ não é universalmente melhor. A dispersão é um
parâmetro de projeto, e não apenas um defeito a minimizar.}
\end{questao}

\begin{questao}[4]{}
\textbf{(Fluxo e energia.)} Deduza a energia armazenada em um circuito
magnético e mostre onde ela reside.
\begin{itensq}
\item Deduza $W=\frac12\mathcal{R}\Phi^{2}$.
\item Determine a densidade de energia em função de $B$ e $\mu$.
\item Determine a fração de energia no entreferro do exemplo anterior.
\end{itensq}
\resposta{$u=B^{2}/2\mu$; praticamente toda a energia fica no entreferro}
\resolucao{(a) Com $L=N^{2}/\mathcal{R}$ e $\lambda=N\Phi=LI$:
\[
W=\frac12LI^{2}=\frac{\lambda^{2}}{2L}
=\frac{(N\Phi)^{2}}{2N^{2}/\mathcal{R}}
=\frac{1}{2}\mathcal{R}\,\Phi^{2}.
\]
\textbf{Analogia:} é a forma magnética de $W=\frac12 CU^{2}$ escrita como
$Q^{2}/2C$ --- o fluxo faz o papel da carga e a relutância, o do inverso da
capacitância.\\
(b) Escrevendo $\mathcal{R}=\ell/\mu A$ e $\Phi=BA$:
\[
W=\frac{1}{2}\frac{\ell}{\mu A}(BA)^{2}=\frac{B^{2}}{2\mu}(A\ell)
\;\Longrightarrow\;
u=\frac{B^{2}}{2\mu}.
\]
\textbf{Ponto decisivo:} a densidade de energia é \emph{inversamente}
proporcional a $\mu$. Para o mesmo $B$, o ar armazena $\mu_r$ vezes \emph{mais}
energia por volume que o ferro.\\
(c) \textbf{Fração no entreferro.} Como $B$ é o mesmo nos dois meios (fluxo
conservado e mesma seção), a razão das energias é a razão dos produtos
$\ell/\mu$:
\[
\frac{W_g}{W_{fe}}
=\frac{g/\mu_0}{\ell_{fe}/(\mu_r\mu_0)}
=\frac{g\,\mu_r}{\ell_{fe}}
=\frac{(\pot{0.56}{-3})(2000)}{0{,}50}=2{,}24.
\]
Ou seja, cerca de $69\%$ da energia está no entreferro, contra $31\%$ no
ferro --- e a proporção cresce rapidamente com $g$.\\
\textbf{Por isso o entreferro é o elemento de projeto de um indutor de
potência:} ele é o \emph{reservatório}. O ferro apenas conduz o fluxo até ele.}
\end{questao}

\begin{questao}[4]{}
\textbf{(Núcleo de seção variável.)} Formule o cálculo do fluxo em um núcleo
cuja seção varia ao longo do caminho, sem dispersão.
\begin{itensq}
\item Escreva a relutância total.
\item Determine o critério de dimensionamento.
\item Analise o caso de um núcleo com estreitamento localizado.
\end{itensq}
\resposta{$\mathcal{R}=\displaystyle\int\frac{\mathrm{d}\ell}{\mu A(\ell)}$;
dimensiona-se pela \emph{menor} seção}
\resolucao{(a) Cada trecho infinitesimal contribui como uma relutância em
série:
\[
\mathcal{R}=\int_0^{L}\frac{\mathrm{d}\ell}{\mu\,A(\ell)}.
\]
Para trechos de seção constante, isso vira a soma
$\sum\ell_i/(\mu A_i)$.\\
(b) \textbf{Critério de dimensionamento.} Sem dispersão, o fluxo é o
\emph{mesmo} em toda a extensão. Logo
\[
B(\ell)=\frac{\Phi}{A(\ell)},
\]
e $B$ é máximo onde $A$ é \textbf{mínimo}. A condição de não saturação é
\[
\Phi\le B_{sat}\,A_{\min}.
\]
\textbf{Toda a capacidade do núcleo é ditada pela sua menor seção} --- alargar
o resto não ajuda em nada.\\
(c) \textbf{Estreitamento localizado.} Suponha um núcleo de
$\SI{6}{\centi\meter\squared}$ com um trecho curto de
$\SI{2}{\centi\meter\squared}$:
\begin{itemize}
\item \textbf{Efeito na relutância:} se o trecho estreito for curto, sua
contribuição a $\mathcal{R}$ é pequena --- o fluxo total quase não muda.
\item \textbf{Efeito na saturação:} ali $B$ é \emph{três vezes} maior, e o
trecho satura com um terço do fluxo que o resto suportaria.
\item \textbf{Efeito local:} saturado, aquele trecho aquece por perdas
concentradas e sua permeabilidade cai --- criando um entreferro
\emph{efetivo} não projetado.
\end{itemize}
\textbf{Onde isso ocorre na prática:} furos de fixação, cantos internos de
núcleos em E, e emendas de lâminas. São defeitos que passam despercebidos no
cálculo de relutância (pequeno efeito) mas dominam o comportamento na saturação
(efeito grande).\\
\textbf{Regra:} calcule $\mathcal{R}$ pelo caminho todo, mas verifique
$B$ na \emph{seção crítica}.}
\end{questao}


 
 

% END BANCO COMPLEMENTAR

\gabarito[12.3 Fluxo magnético]

% =====================================================================
\subsection{Força magnética}
% =====================================================================

\blocofixacao

\begin{questao}[1]{}
A força magnética sobre uma carga em movimento é \textbf{nula} quando:
\begin{alternativas}
\item a velocidade é perpendicular ao campo.
\item a velocidade é paralela (ou antiparalela) ao campo.
\item a carga é positiva.
\item o campo é muito intenso.
\item a velocidade é máxima.
\end{alternativas}
\resposta{(b)}
\resolucao{Como $F = qvB\sin\theta$, a força se anula quando $\sin\theta = 0$, ou
seja, quando $\vec{v}$ é paralela (ou antiparalela) a $\vec{B}$. É máxima quando
$\vec{v}\perp\vec{B}$ ($\theta = \SI{90}{\degree}$).}
\end{questao}

\begin{questao}[1]{}
Um condutor retilíneo de $\SI{30}{\centi\meter}$ conduz $\SI{2}{\ampere}$
perpendicularmente a um campo de $\SI{0.5}{\tesla}$. Determine a força sobre ele.
\resposta{$F = \SI{0.3}{\newton}$}
\resolucao{$F = B\,I\,\ell = 0{,}5\cdot2\cdot0{,}3 = \SI{0.3}{\newton}$
(com $\sin\theta = 1$, pois são perpendiculares).}
\end{questao}

\begin{questao}[1]{}
Uma carga positiva move-se horizontalmente para a direita, em uma região onde o
campo magnético aponta para dentro do plano da página. A força magnética sobre a
carga aponta:
\begin{alternativas}[2]
\item para cima.
\item para baixo.
\item para a direita.
\item para a esquerda.
\item para dentro da página.
\end{alternativas}
\resposta{(a)}
\resolucao{Pela regra da mão direita para $\vec F = q\vec v\times\vec B$: dedos
apontando no sentido de $\vec v$ (direita), curvando-se para $\vec B$ (para dentro
da página), o polegar aponta para \emph{cima}. Como a carga é positiva, a força
tem esse mesmo sentido. (Se fosse negativa, seria para baixo.)}
\end{questao}

\blocoaplicacao

\begin{questao}[2]{}
Uma carga de $\SI{1.6e-19}{\coulomb}$ move-se a
$\SI{e6}{\meter\per\second}$, perpendicularmente a um campo de $\SI{0.2}{\tesla}$.
Determine o módulo da força magnética e descreva a trajetória resultante.
\resposta{$F = \SI{3.2e-14}{\newton}$; trajetória circular}
\resolucao{$F = q\,v\,B = \pot{1.6}{-19}\cdot\pot{1}{6}\cdot0{,}2
= \pot{3.2}{-14}\,\si{\newton}$.\\
Como a força é sempre \emph{perpendicular} à velocidade, ela não altera o módulo
de $v$ (não realiza trabalho), apenas curva a trajetória: o resultado é um
\textbf{movimento circular uniforme}. É o princípio dos cíclotrons e dos
espectrômetros de massa.}
\end{questao}

\begin{questao}[2]{}
Uma espira retangular percorrida por corrente é colocada em um campo magnético
uniforme. Explique por que ela tende a \textbf{girar}, e não a se deslocar, e
relacione isso com o funcionamento de um motor elétrico.
\resposta{Ver comentário}
\resolucao{Em um campo uniforme, os dois lados opostos da espira conduzem corrente
em sentidos \emph{contrários}, sofrendo forças magnéticas iguais e opostas
($\vec F = B I \ell$). Essas forças formam um \textbf{binário} (conjugado): a
resultante é nula (não há translação), mas o \emph{torque} não é --- a espira
gira. Esse é exatamente o princípio do \textbf{motor de corrente contínua}: o
torque $\tau = N B I A \sin\theta$ faz o rotor girar, e o comutador inverte a
corrente a cada meia-volta para manter a rotação no mesmo sentido.}
\end{questao}

\blocoaprofundamento

\begin{questao}[3]{}
Um próton ($q = \SI{1.6e-19}{\coulomb}$, $m = \SI{1.67e-27}{\kilogram}$) entra
perpendicularmente em um campo uniforme de $\SI{0.5}{\tesla}$ com velocidade
$\SI{2e6}{\meter\per\second}$. Determine o raio da trajetória circular descrita.
\resposta{$r \approx \SI{4.2}{\centi\meter}$}
\resolucao{A força magnética atua como resultante centrípeta:
$qvB = \dfrac{mv^2}{r}$. Isolando $r$:
\[
r=\frac{mv}{qB}
=\frac{(\pot{1.67}{-27})(\pot{2}{6})}{(\pot{1.6}{-19})(0{,}5)}
\approx\pot{4.2}{-2}\,\si{\meter}=\SI{4.2}{\centi\meter}.
\]
O raio cresce com a quantidade de movimento ($mv$) e diminui com a carga e o
campo. Medindo $r$ conhece-se a razão carga/massa da partícula --- é assim que
funciona o espectrômetro de massa.}
\end{questao}

\begin{questao}[3]{}
Um seletor de velocidades é formado por campos elétrico e magnético
\textbf{cruzados} (perpendiculares entre si e à velocidade). Partículas
carregadas entram com velocidades variadas e apenas as de uma velocidade
específica atravessam sem desviar.

\circuitoq{%
\begin{tikzpicture}[scale=1,line width=0.8pt,>=Stealth]
 \draw[cinza!50,fill=cinza!8] (0,-1.2) rectangle (5.4,1.2);
 % campo E (setas para baixo)
 \foreach \x in {0.7,1.6,2.5,3.4,4.3}
   \draw[->,Vcol!70] (\x,1.05)--(\x,-1.05);
 \node[Vcol,font=\footnotesize] at (5.05,0.9) {$\vec E$};
 % campo B (para dentro da pagina)
 \foreach \x in {1.15,2.05,2.95,3.85}
   \foreach \y in {-0.6,0,0.6}
     {\node[Icol,font=\scriptsize] at (\x,\y) {$\otimes$};}
 \node[Icol,font=\footnotesize] at (0.35,-0.9) {$\vec B$};
 % particula
 \draw[->,laranja,line width=1.2pt] (-1.3,0)--(0,0)
       node[midway,above,font=\footnotesize]{$v$};
 \draw[->,laranja,line width=1.2pt] (5.4,0)--(6.5,0);
 \fill[laranja] (-1.3,0) circle (0.1);
\end{tikzpicture}}{Seletor de velocidades com campos cruzados.}

\begin{itensq}
\item Deduza a condição para que a partícula atravesse sem desvio.
\item Mostre que essa condição \emph{não depende} da carga nem da massa da
      partícula.
\item Com $E = \SI{3e5}{\volt\per\meter}$ e $B = \SI{0.15}{\tesla}$, qual a
      velocidade selecionada?
\item Após o seletor, coloca-se um campo magnético $B'$ que curva a trajetória.
      Explique como essa combinação permite medir a razão $q/m$.
\end{itensq}
\resposta{(a) $qE = qvB \Rightarrow v = E/B$; (c) $v = \SI{2e6}{\meter\per\second}$;
(b) e (d) ver comentário}
\resolucao{\textbf{(a)} Sobre a partícula atuam duas forças perpendiculares à
trajetória e \emph{opostas} entre si:
\[
F_e = qE \quad(\text{para baixo, se } q>0),
\qquad
F_m = qvB \quad(\text{para cima, pela regra da mão direita}).
\]
Para não haver desvio, a resultante deve ser nula:
\[
qE = qvB
\;\Longrightarrow\;
\boxed{\,v=\frac{E}{B}\,}
\]
\textbf{(b)} A carga $q$ \textbf{cancela-se} nos dois lados da igualdade, e a
massa sequer aparece --- não há aceleração, logo a segunda lei de Newton não é
invocada. Consequência notável: o seletor filtra por \emph{velocidade}, e apenas
por ela. Elétrons, prótons e íons pesados com a mesma $v$ atravessam igualmente.
Partículas mais rápidas são desviadas pela força magnética (que cresce com $v$);
mais lentas, pela força elétrica.

\textbf{(c)} $v = \dfrac{E}{B} = \dfrac{\pot{3}{5}}{0{,}15}
= \pot{2}{6}\,\si{\meter\per\second}$.

\textbf{(d) O espectrômetro de massa.} Com a velocidade agora \emph{conhecida e
uniforme}, a partícula entra em uma região de campo $B'$ puro, onde a força
magnética atua como centrípeta:
\[
qvB' = \frac{mv^{2}}{r}
\;\Longrightarrow\;
\frac{q}{m}=\frac{v}{B'r}=\frac{E}{B\,B'\,r}.
\]
Medindo o raio $r$ da trajetória (pela marca deixada no detector) e conhecendo
$E$, $B$ e $B'$, obtém-se $q/m$ diretamente. Íons de massas diferentes descrevem
raios diferentes e se separam espacialmente --- daí o nome
\emph{espectrômetro}.\\
\textbf{Impacto histórico:} foi com esse arranjo que Thomson mediu $q/m$ do
elétron (1897) e que Aston descobriu os \textbf{isótopos} (1919), rendendo-lhe o
Nobel. Hoje é ferramenta padrão em química analítica, datação e antidoping.}
\end{questao}

% BEGIN BANCO COMPLEMENTAR Força magnética
% =====================================================================
%  Banco complementar --- Forca Magnetica  (REVISADO)
%  Substitui a versao gerada por template (5 clones em N1, 10 em N2,
%  10 em N3 e 8 questoes conceituais marcadas como N4).
%  O cap11 ja cobre: condicao de forca nula (MC), condutor de 30 cm com
%  2 A em 0,5 T, regra da mao para carga movendo-se com B entrando no
%  plano, carga elementar a 1e6 m/s em 0,2 T, espira retangular girando
%  (qualitativo), proton em campo uniforme (raio) e o SELETOR DE
%  VELOCIDADES com campos cruzados -- motivo pelo qual é a N4 antiga sobre
%  projeto de seletor foi removida deste banco. Ela e substituida pelo
%  ESPECTROMETRO DE MASSA, que usa o seletor como etapa e vai alem.
%  Este banco explora: forca entre fios e a definicao do ampere,
%  trajetoria helicoidal, torque e momento magnetico, efeito Hall
%  quantitativo, campo nao uniforme, projeto de atuador e a demonstracao
%  de que a forca magnetica nao realiza trabalho.
%  ACRESCIMO POSTERIOR: tres questoes sobre a FORCA DE LORENTZ COMPLETA,
%  F = q(E + v x B). As duas parcelas ja apareciam separadamente e os
%  campos cruzados surgiam no seletor de velocidades e na deriva ExB,
%  mas a expressao unificada nao estava escrita em lugar nenhum.
%  Distribuicao: 5 N1 + 11 N2 + 11 N3 + 9 N4 = 36
% =====================================================================

\subsubsection*{Banco complementar --- Força magnética}

\begin{atencao}[Organização do banco]
Duas expressões governam tudo: $\vec F=q\,\vec v\times\vec B$ para cargas e
$\vec F=I\,\vec\ell\times\vec B$ para condutores. Ambas são \textbf{produtos
vetoriais} --- a força é \emph{perpendicular} tanto à velocidade quanto ao
campo, e se anula quando eles são paralelos. Daí decorre o fato central: a força
magnética \textbf{nunca realiza trabalho}. O N3 trata de trajetórias,
espectrometria e efeito Hall; o N4 exige demonstrações e projeto.
\end{atencao}

\blocofixacao

\begin{questao}[1]{}
Uma carga $q=+\SI{2}{\micro\coulomb}$ move-se a $\SI{500}{\meter\per\second}$
perpendicularmente a um campo de $\SI{0.2}{\tesla}$. Determine a força.
\resposta{$F=\SI{0.2}{\milli\newton}$}
\resolucao{
\[
F=qvB\sin\theta=(\pot{2}{-6})(500)(0{,}2)(1)=\pot{2}{-4}\,\si{\newton}.
\]
\textbf{Direção:} perpendicular \emph{simultaneamente} a $\vec v$ e a
$\vec B$, dada pela regra da mão direita aplicada a
$\vec v\times\vec B$ (e invertida, se a carga for negativa).}
\end{questao}

\begin{questao}[1]{}
Um condutor retilíneo de $\SI{40}{\centi\meter}$ conduz $\SI{3}{\ampere}$
perpendicularmente a $B=\SI{0.25}{\tesla}$. Determine a força.
\resposta{$F=\SI{0.3}{\newton}$}
\resolucao{
\[
F=BIL\sin\theta=(0{,}25)(3)(0{,}40)(1)=\SI{0.3}{\newton}.
\]
\textbf{Origem da fórmula:} ela é a soma das forças sobre todos os portadores do
condutor. Escrevendo $I=nqv_dA$ e $L$ o comprimento, o produto
$nqv_dA\cdot L$ reproduz $qv_d$ somado sobre $nAL$ portadores.\\
A força é transmitida ao \emph{fio} pelas colisões dos portadores com a rede
cristalina --- não são os elétrons que "empurram" o fio diretamente.}
\end{questao}

\begin{questao}[1]{}
A mesma carga da primeira questão move-se agora formando $\ang{30}$ com o
campo. Determine a força.
\resposta{$F=\SI{0.1}{\milli\newton}$}
\resolucao{
\[
F=qvB\sin\ang{30}=(\pot{2}{-4})(0{,}5)=\pot{1}{-4}\,\si{\newton}.
\]
\textbf{Metade} do valor perpendicular.\\
\textbf{Casos-limite:} $\theta=\ang{90}$ dá força máxima; $\theta=0$ ou
$\ang{180}$ (velocidade \emph{paralela} ao campo) dá força \textbf{nula} --- a
carga segue em linha reta, sem sentir o campo.}
\end{questao}

\begin{questao}[1]{}
Uma carga positiva move-se para a \textbf{direita} em uma região onde
$\vec B$ aponta para \textbf{cima}. Determine o sentido da força.
\begin{alternativas}[2]
\item para dentro do plano
\item para fora do plano
\item para cima
\item para a direita
\end{alternativas}
\resposta{(a)}
\resolucao{Pela regra da mão direita para $\vec v\times\vec B$: dedos apontando
no sentido de $\vec v$ (direita), curvando-se para $\vec B$ (cima) --- o polegar
aponta para \textbf{dentro} do plano.\\
\textbf{Verificação:} a força deve ser perpendicular a ambos. "Direita" e
"cima" definem o plano da folha, logo a força só pode ser perpendicular a ele
--- o que elimina (c) e (d) imediatamente.\\
\textbf{Se a carga fosse negativa}, a força seria para fora do plano.}
\end{questao}

\begin{questao}[1]{}
A partir de $F=qvB$, mostre que $\SI{1}{\tesla}$ equivale a
$\SI{1}{\newton\second\per\coulomb\per\meter}$.
\resposta{$\si{\tesla}=\dfrac{\si{\newton}}{\si{\coulomb}\cdot
\si{\meter\per\second}}$}
\resolucao{
\[
B=\frac{F}{qv}
\;\Longrightarrow\;
[\si{\tesla}]=\frac{\si{\newton}}{\si{\coulomb}\cdot\si{\meter\per\second}}
=\frac{\si{\newton}\cdot\si{\second}}{\si{\coulomb}\cdot\si{\meter}}
=\frac{\si{\newton}}{\si{\ampere}\cdot\si{\meter}}.
\]
A última forma, $\si{\newton\per\ampere\per\meter}$, decorre diretamente de
$F=BIL$ --- e é a mais prática para conferir contas com condutores.\\
\textbf{O tesla é uma unidade grande:} $\SI{1}{\tesla}$ exerce
$\SI{1}{\newton}$ sobre um metro de fio conduzindo $\SI{1}{\ampere}$.}
\end{questao}

\blocoaplicacao

\begin{questao}[2]{}
Um fio de $\SI{10}{\centi\meter}$ conduz $\SI{2}{\ampere}$ em um campo de
$\SI{0.2}{\tesla}$, formando $\ang{30}$ com ele. Determine a força.
\resposta{$F=\SI{20}{\milli\newton}$}
\resolucao{
\[
F=BIL\sin\ang{30}=(0{,}2)(2)(0{,}10)(0{,}5)=\pot{2}{-2}\,\si{\newton}.
\]
\textbf{Interpretação do seno:} apenas a componente de $\vec\ell$
\emph{perpendicular} a $\vec B$ contribui. Equivalentemente, o comprimento
\emph{efetivo} é $L\sin\theta=\SI{5}{\centi\meter}$.\\
\textbf{Direção da força:} perpendicular ao plano definido pelo fio e pelo
campo --- e \emph{não} perpendicular apenas ao fio.}
\end{questao}

\begin{questao}[2]{}
Um próton entra perpendicularmente em um campo de $\SI{0.5}{\tesla}$ com
$v=\pot{2}{6}\,\si{\meter\per\second}$.
\begin{itensq}
\item Determine o raio da trajetória.
\item Determine o período.
\end{itensq}
\dados{$e=\pot{1.6}{-19}\,\si{\coulomb}$; $m_p=\pot{1.67}{-27}\,\si{\kilogram}$}
\resposta{$r=\SI{4.18}{\centi\meter}$; $T=\SI{131}{\nano\second}$}
\resolucao{(a) A força magnética é centrípeta:
\[
qvB=\frac{mv^{2}}{r}
\;\Longrightarrow\;
r=\frac{mv}{qB}=\frac{(\pot{1.67}{-27})(\pot{2}{6})}{(\pot{1.6}{-19})(0{,}5)}
=\SI{0.0418}{\meter}.
\]
(b)
\[
T=\frac{2\pi r}{v}=\frac{2\pi m}{qB}
=\frac{2\pi(\pot{1.67}{-27})}{(\pot{1.6}{-19})(0{,}5)}
=\pot{1.31}{-7}\,\si{\second}.
\]
\textbf{Fato notável:} o período \textbf{não depende da velocidade} nem do raio.
Partículas rápidas descrevem círculos maiores, mas levam o mesmo tempo para
completá-los.\\
É esse fato que torna possível o \textbf{cíclotron}: a frequência de aceleração
pode ser fixa, embora a partícula ganhe energia a cada volta.}
\end{questao}

\begin{questao}[2]{}
Um elétron e um próton, com a mesma velocidade, entram no mesmo campo
magnético perpendicularmente. Compare seus raios e períodos.
\resposta{$r_p/r_e=1836$; $T_p/T_e=1836$}
\resolucao{Como $r=mv/(qB)$ e $T=2\pi m/(qB)$, ambos são
\emph{proporcionais à massa} (as cargas têm o mesmo módulo):
\[
\frac{r_p}{r_e}=\frac{T_p}{T_e}=\frac{m_p}{m_e}
=\frac{\pot{1.67}{-27}}{\pot{9.11}{-31}}=1836.
\]
\textbf{Sentidos opostos:} tendo cargas de sinais contrários, as duas partículas
curvam-se para lados opostos --- é assim que se identifica o sinal da carga em
câmaras de detecção.\\
\textbf{Consequência experimental:} para conter elétrons e prótons no mesmo
aparelho, o raio do próton é quase dois mil vezes maior. Aceleradores de prótons
são, por isso, muito maiores que os de elétrons para a mesma energia.}
\end{questao}

\begin{questao}[2]{}
Dois fios paralelos longos, separados por $\SI{5}{\centi\meter}$, conduzem
$\SI{10}{\ampere}$ cada. Determine a força por unidade de comprimento.
\dados{$\mu_0=4\pi\times10^{-7}\,\si{\tesla\meter\per\ampere}$}
\resposta{$F/L=\SI{0.4}{\milli\newton\per\meter}$}
\resolucao{Um fio cria campo $B=\dfrac{\mu_0I_1}{2\pi d}$ na posição do outro,
que sofre $F=BI_2L$:
\[
\frac{F}{L}=\frac{\mu_0I_1I_2}{2\pi d}
=\pot{2}{-7}\frac{(10)(10)}{0{,}05}=\pot{4}{-4}\,\si{\newton\per\meter}.
\]
\textbf{Sentido:} correntes de \emph{mesmo} sentido \textbf{atraem-se};
de sentidos opostos, \textbf{repelem-se}.\\
\textbf{Cuidado com a intuição elétrica:} cargas de mesmo sinal se repelem, mas
correntes de mesmo sentido se atraem. As duas situações não são análogas.}
\end{questao}

\begin{questao}[2]{}
Uma espira retangular de área $\SI{100}{\centi\meter\squared}$ com $50$ voltas
conduz $\SI{2}{\ampere}$ em um campo de $\SI{0.3}{\tesla}$.
\begin{itensq}
\item Determine o momento de dipolo magnético.
\item Determine o torque máximo.
\end{itensq}
\resposta{$m=\SI{1}{\ampere\meter\squared}$; $\tau=\SI{0.3}{\newton\meter}$}
\resolucao{(a) $A=\SI{100}{\centi\meter\squared}=\pot{1}{-2}\,\si{\meter\squared}$:
\[
m=NIA=50(2)(\pot{1}{-2})=\SI{1}{\ampere\meter\squared}.
\]
(b)
\[
\tau=mB\sin\theta
\;\Longrightarrow\;
\tau_{\max}=mB=(1)(0{,}3)=\SI{0.3}{\newton\meter}.
\]
\textbf{Quando ocorre o máximo:} com o plano da espira \emph{paralelo} ao campo,
isto é, $\vec m$ perpendicular a $\vec B$. O torque é \textbf{nulo} quando
$\vec m$ está alinhado com $\vec B$ --- posição de equilíbrio.\\
\textbf{Analogia útil:} a espira comporta-se como uma agulha de bússola, e
$\vec m$ faz o papel do "norte" dela.}
\end{questao}

\begin{questao}[2]{}
Explique por que a força magnética não altera o módulo da velocidade de uma
partícula carregada.
\resposta{Ela é sempre perpendicular a $\vec v$, logo não realiza trabalho}
\resolucao{A potência entregue por uma força é $P=\vec F\cdot\vec v$. Como
\[
\vec F=q\,\vec v\times\vec B
\;\Longrightarrow\;
\vec F\perp\vec v
\;\Longrightarrow\;
\vec F\cdot\vec v=0.
\]
Sem potência, não há variação de energia cinética --- e portanto o \emph{módulo}
da velocidade permanece constante.\\
\textbf{O que a força magnética faz:} muda a \emph{direção} do movimento, e só
isso. Ela curva a trajetória sem acelerar nem frear.\\
\textbf{Consequência prática:} um campo magnético sozinho jamais acelera
partículas. Aceleradores usam campos \emph{elétricos} para ganhar energia e
campos magnéticos apenas para \emph{guiar} o feixe.}
\end{questao}

\begin{questao}[2]{}
Um elétron é abandonado \textbf{em repouso} em uma região onde coexistem
$E=\SI{500}{\volt\per\meter}$ e $B=\SI{0.02}{\tesla}$, perpendiculares entre
si.
\begin{itensq}
\item Determine a força e a aceleração no instante inicial.
\item Descreva qualitativamente o movimento subsequente.
\end{itensq}
\dados{$e=\pot{1.6}{-19}\,\si{\coulomb}$; $m_e=\pot{9.11}{-31}\,\si{\kilogram}$}
\resposta{$F=\SI{8e-17}{\newton}$ (só elétrica);
$a=\pot{8.8}{13}\,\si{\meter\per\second\squared}$}
\resolucao{(a) A força total é a \textbf{força de Lorentz}:
\[
\vec F=q\left(\vec E+\vec v\times\vec B\right).
\]
Com $\vec v=\vec 0$, a parcela magnética \textbf{se anula} --- resta apenas a
elétrica:
\[
F=eE=(\pot{1.6}{-19})(500)=\pot{8}{-17}\,\si{\newton};
\qquad
a=\frac{F}{m_e}=\pot{8.8}{13}\,\si{\meter\per\second\squared}.
\]
(b) \textbf{Movimento subsequente.} Assim que o elétron adquire velocidade, a
parcela magnética \emph{deixa de ser nula} e passa a encurvar a trajetória. A
força elétrica continua acelerando; a magnética continua desviando.\\
O resultado é uma trajetória \textbf{cicloidal}, com deriva média
\[
v_d=\frac{E}{B}=\frac{500}{0{,}02}=\pot{2.5}{4}\,\si{\meter\per\second},
\]
perpendicular a ambos os campos.\\
\textbf{Ponto essencial:} a força magnética \emph{depende da velocidade} --- ela
não age sobre cargas paradas. Já a elétrica age sempre. É essa assimetria que
faz o movimento começar de um jeito e evoluir para outro.}
\end{questao}

\begin{questao}[2]{}
Um fio em forma de semicircunferência de raio $R$ conduz corrente $I$ em um
campo uniforme $\vec B$ perpendicular ao plano do semicírculo. Determine a
força resultante.
\resposta{$F=2RIB$, como se o fio fosse retilíneo de comprimento $2R$}
\resolucao{Para um fio de forma qualquer em campo \textbf{uniforme},
\[
\vec F=I\left(\int\mathrm{d}\vec\ell\right)\times\vec B
=I\,\vec L_{ef}\times\vec B,
\]
onde $\vec L_{ef}$ é o vetor que liga o \textbf{início ao fim} do fio --- a
\emph{corda}, não o comprimento percorrido.\\
Para a semicircunferência, a corda vale $2R$:
\[
F=2RIB.
\]
\textbf{Poder do resultado:} a forma do fio é irrelevante em campo uniforme.
Um fio retorcido, ondulado ou em zigue-zague sofre a mesma força que o segmento
reto entre suas extremidades.\\
\textbf{Corolário importante:} em uma \emph{espira fechada}, a corda é nula ---
logo a força resultante sobre qualquer espira em campo uniforme é
\textbf{zero}. Só há torque.}
\end{questao}

\begin{questao}[2]{}
Uma balança de corrente tem um fio horizontal de $\SI{20}{\centi\meter}$ em um
campo vertical de $\SI{0.4}{\tesla}$. Determine a corrente necessária para
equilibrar uma massa de $\SI{2}{\gram}$.
\dados{$g=\SI{9.8}{\meter\per\second\squared}$}
\resposta{$I=\SI{0.245}{\ampere}$}
\resolucao{Equilíbrio entre peso e força magnética:
\[
BIL=mg
\;\Longrightarrow\;
I=\frac{mg}{BL}=\frac{(\pot{2}{-3})(9{,}8)}{(0{,}4)(0{,}20)}
=\SI{0.245}{\ampere}.
\]
\textbf{Uso do arranjo:} conhecendo $m$, $L$ e $I$, obtém-se $B$ com boa
exatidão --- é o método clássico de calibração absoluta de campo magnético, sem
depender de sensores.\\
\textbf{Inversamente}, conhecendo $B$, mede-se corrente. Foi assim que o ampère
foi definido durante décadas: pela força entre condutores.}
\end{questao}

\begin{questao}[2]{}
Determine a velocidade de uma partícula que atravessa sem desvio uma região com
$E=\SI{2e4}{\volt\per\meter}$ e $B=\SI{0.1}{\tesla}$ perpendiculares entre si
e à velocidade.
\resposta{$v=\pot{2}{5}\,\si{\meter\per\second}$}
\resolucao{Para não haver desvio, as forças elétrica e magnética devem se
cancelar:
\[
qE=qvB
\;\Longrightarrow\;
v=\frac{E}{B}=\frac{\pot{2}{4}}{0{,}1}=\pot{2}{5}\,\si{\meter\per\second}.
\]
\textbf{Duas propriedades notáveis:}
\begin{itemize}
\item A condição \textbf{não depende da carga} --- $q$ cancela. Íons e elétrons
com a mesma velocidade passam igualmente.
\item Também \textbf{não depende da massa}. É pura seleção de \emph{velocidade}.
\end{itemize}
\textbf{Consequência:} partículas mais rápidas são desviadas pela força
magnética (maior), e mais lentas pela elétrica. Apenas a velocidade $E/B$
atravessa em linha reta.}
\end{questao}

\begin{questao}[2]{}
Uma espira quadrada de lado $\SI{10}{\centi\meter}$ conduz $\SI{5}{\ampere}$
em um campo uniforme de $\SI{0.2}{\tesla}$ paralelo ao seu plano.
\begin{itensq}
\item Determine a força resultante.
\item Determine o torque.
\end{itensq}
\resposta{$\vec F=\vec 0$; $\tau=\SI{10}{\milli\newton\meter}$}
\resolucao{(a) \textbf{Força resultante nula.} Em campo uniforme, os quatro
lados sofrem forças que se cancelam duas a duas --- lados opostos conduzem em
sentidos contrários.\\
(b) \textbf{Mas o torque não é nulo.} As forças nos dois lados perpendiculares
ao campo formam um \emph{binário}:
\[
\tau=mB\sin\ang{90}=IAB=(5)(\pot{1}{-2})(0{,}2)
=\pot{1}{-2}\,\si{\newton\meter}.
\]
\textbf{Situação de torque máximo:} o campo está no plano da espira, ou seja,
$\vec m$ (perpendicular ao plano) está a $\ang{90}$ de $\vec B$.\\
\textbf{Este é o princípio do motor elétrico} --- e também sua limitação: o
torque cai a zero quando a espira gira até o alinhamento, o que exige comutador
para inverter a corrente e manter a rotação.}
\end{questao}

\blocoaprofundamento

\begin{questao}[3]{}
Uma partícula de carga $q$ e massa $m$ entra perpendicularmente em um campo
uniforme $B$ com velocidade $v$.
\begin{itensq}
\item Deduza o raio e o período da trajetória.
\item Determine a frequência angular (frequência de cíclotron).
\item Explique por que o período independe de $v$.
\end{itensq}
\resposta{$r=\dfrac{mv}{qB}$; $T=\dfrac{2\pi m}{qB}$;
$\omega=\dfrac{qB}{m}$}
\resolucao{(a) A força magnética é perpendicular a $\vec v$ e de módulo
constante --- exatamente as condições de um movimento \textbf{circular
uniforme}. Igualando à resultante centrípeta:
\[
qvB=\frac{mv^{2}}{r}
\;\Longrightarrow\;
r=\frac{mv}{qB}.
\]
\[
T=\frac{2\pi r}{v}=\frac{2\pi m}{qB}.
\]
(b)
\[
\omega=\frac{2\pi}{T}=\frac{qB}{m},
\]
a \textbf{frequência angular de cíclotron}.\\
(c) \textbf{Por que $T$ independe de $v$.} Dobrar a velocidade dobra a força
magnética \emph{e} dobra o raio necessário --- os dois efeitos se compensam
exatamente. A partícula percorre um círculo duas vezes maior a velocidade duas
vezes maior, no mesmo tempo.\\
Algebricamente, $v$ aparece em $r$ e é cancelado por $v$ em $T=2\pi r/v$.\\
\textbf{Consequência:} o cíclotron pode operar com frequência \emph{fixa}.
\emph{(Isso vale enquanto $v\ll c$; em regime relativístico a massa efetiva
cresce, $T$ aumenta e é preciso modular a frequência --- o síncrotron.)}}
\end{questao}

\begin{questao}[3]{}
Uma partícula entra em um campo uniforme formando $\ang{30}$ com $\vec B$, com
$v=\pot{1}{6}\,\si{\meter\per\second}$ e $B=\SI{0.1}{\tesla}$ (próton).
\begin{itensq}
\item Decomponha a velocidade e descreva a trajetória.
\item Determine o raio e o passo da hélice.
\end{itensq}
\resposta{Hélice com $r=\SI{5.2}{\centi\meter}$ e passo de
$\SI{0.57}{\meter}$}
\resolucao{(a) \textbf{Decomposição:}
\[
v_\perp=v\sin\ang{30}=\pot{5}{5}\,\si{\meter\per\second};
\qquad
v_\parallel=v\cos\ang{30}=\pot{8.66}{5}\,\si{\meter\per\second}.
\]
A componente \textbf{perpendicular} sofre força e produz movimento circular. A
componente \textbf{paralela} não sofre força alguma ($\vec v\times\vec B=0$
para ela) e produz movimento retilíneo uniforme.\\
A composição é uma \textbf{hélice} de eixo paralelo a $\vec B$.\\
(b)
\[
r=\frac{mv_\perp}{qB}
=\frac{(\pot{1.67}{-27})(\pot{5}{5})}{(\pot{1.6}{-19})(0{,}1)}
=\SI{0.052}{\meter};
\]
\[
T=\frac{2\pi m}{qB}=\pot{6.56}{-7}\,\si{\second};
\qquad
\text{passo}=v_\parallel T=(\pot{8.66}{5})(\pot{6.56}{-7})=\SI{0.568}{\meter}.
\]
\textbf{Aplicação:} é essa trajetória que confina partículas nos cinturões de
Van Allen e em espelhos magnéticos --- as partículas espiralam ao longo das
linhas de campo, e o passo diminui onde o campo se intensifica, até a reversão
do movimento.}
\end{questao}

\begin{questao}[3]{}
Um espectrômetro de massa usa um seletor de velocidades ($v=\pot{1}{5}\,
\si{\meter\per\second}$) seguido de um campo de $\SI{0.1}{\tesla}$. Íons de
neônio de massas $20\,u$ e $22\,u$, todos com carga $+e$, entram no segundo
campo.
\begin{itensq}
\item Determine o raio de cada trajetória.
\item Determine a separação entre os pontos de chegada no detector.
\end{itensq}
\dados{$u=\pot{1.66}{-27}\,\si{\kilogram}$}
\resposta{$r_{20}=\SI{20.8}{\centi\meter}$, $r_{22}=\SI{22.8}{\centi\meter}$;
separação de $\SI{4.2}{\centi\meter}$}
\resolucao{(a) Com $r=mv/(qB)$ e velocidade já selecionada:
\[
r_{20}=\frac{20(\pot{1.66}{-27})(\pot{1}{5})}{(\pot{1.6}{-19})(0{,}1)}
=\SI{0.2075}{\meter};
\qquad
r_{22}=\SI{0.2283}{\meter}.
\]
(b) Após meia volta, cada íon atinge o detector a uma distância $2r$ do ponto
de entrada:
\[
\Delta=2\left(r_{22}-r_{20}\right)=2(0{,}0208)=\SI{4.15}{\centi\meter}.
\]
\textbf{Por que o seletor é indispensável:} sem ele, íons de massas diferentes
mas velocidades diferentes poderiam produzir o \emph{mesmo} raio --- pois
$r\propto mv$. O seletor fixa $v$, tornando $r$ função apenas de $m/q$.\\
\textbf{Resolução do instrumento:} a separação relativa é
$\Delta r/r=\Delta m/m=10\%$ para esses isótopos --- larga. Separar isótopos
com $\Delta m/m$ de $0{,}1\%$ exigiria feixe muito bem colimado e detector de
alta resolução espacial.}
\end{questao}

\begin{questao}[3]{}
A definição do ampère baseava-se na força entre condutores paralelos.
\begin{itensq}
\item Determine a força por metro entre dois fios distantes
      $\SI{1}{\meter}$, conduzindo $\SI{1}{\ampere}$ cada.
\item Explique o papel dessa relação na definição da unidade.
\item Comente a mudança recente.
\end{itensq}
\resposta{$F/L=\pot{2}{-7}\,\si{\newton\per\meter}$}
\resolucao{(a)
\[
\frac{F}{L}=\frac{\mu_0I^{2}}{2\pi d}
=\frac{(4\pi\times10^{-7})(1)}{2\pi(1)}=\pot{2}{-7}\,\si{\newton\per\meter}.
\]
(b) \textbf{Definição antiga (até 2019):} o ampère era \emph{definido} como a
corrente que, em dois condutores paralelos infinitos separados por
$\SI{1}{\meter}$ no vácuo, produzisse exatamente
$\pot{2}{-7}\,\si{\newton}$ por metro.\\
Essa definição \textbf{fixava $\mu_0$} em $4\pi\times10^{-7}$ por construção ---
não era um valor medido, mas uma consequência da escolha.\\
(c) \textbf{Desde 2019}, o ampère é definido fixando-se a \emph{carga
elementar} em $e=\pot{1.602176634}{-19}\,\si{\coulomb}$ exatamente.\\
\textbf{Consequência da mudança:} $\mu_0$ deixou de ser exato e passou a ser uma
grandeza \textbf{medida}, com incerteza experimental --- embora ainda
igual a $4\pi\times10^{-7}$ dentro de partes por bilhão.\\
A troca reflete uma mudança de filosofia: definir unidades por
\emph{constantes fundamentais} em vez de por experimentos difíceis de
reproduzir.}
\end{questao}

\begin{questao}[3]{}
Um sensor Hall de espessura $\SI{0.1}{\milli\meter}$ conduz
$\SI{1}{\ampere}$ em campo de $\SI{0.5}{\tesla}$.
\begin{itensq}
\item Determine a tensão Hall para cobre ($n=\pot{8.5}{28}\,\si{\per\meter\cubed}$).
\item Repita para um semicondutor com $n=\pot{1}{22}\,\si{\per\meter\cubed}$.
\item Explique a diferença.
\end{itensq}
\resposta{$\SI{0.37}{\micro\volt}$ no cobre; $\SI{3.1}{\volt}$ no
semicondutor}
\resolucao{(a) A tensão Hall vale
\[
V_H=\frac{IB}{nqt}
=\frac{(1)(0{,}5)}{(\pot{8.5}{28})(\pot{1.6}{-19})(\pot{1}{-4})}
=\pot{3.68}{-7}\,\si{\volt}.
\]
(b) Com $n=\pot{1}{22}$:
\[
V_H=\frac{0{,}5}{(\pot{1}{22})(\pot{1.6}{-19})(\pot{1}{-4})}=\SI{3.1}{\volt}.
\]
\textbf{Diferença de sete ordens de grandeza.}\\
(c) \textbf{Explicação.} A tensão Hall é inversamente proporcional a $n$. Metais
têm um portador por átomo --- densidade altíssima --- e cada portador precisa se
mover \emph{devagar} para produzir a mesma corrente. Velocidade de deriva baixa
significa força magnética pequena e, portanto, separação de cargas mínima.\\
Semicondutores têm densidade de portadores milhões de vezes menor; eles se movem
muito mais rápido, e o efeito Hall se torna macroscópico.\\
\textbf{Consequência tecnológica:} \emph{todos} os sensores Hall comerciais são
de semicondutor. Um sensor de cobre exigiria amplificação de nanovolts,
inviável na presença de ruído térmico e de tensões termoelétricas.}
\end{questao}

\begin{questao}[3]{}
Um motor elementar tem espira de $200$ voltas, área
$\SI{50}{\centi\meter\squared}$, corrente de $\SI{3}{\ampere}$ e campo de
$\SI{0.4}{\tesla}$.
\begin{itensq}
\item Determine o torque máximo.
\item Determine a potência mecânica a $\SI{1200}{rpm}$, supondo torque médio
      igual a $2/\pi$ do máximo.
\end{itensq}
\resposta{$\tau_{\max}=\SI{1.2}{\newton\meter}$;
$P\approx\SI{96}{\watt}$}
\resolucao{(a)
\[
m=NIA=200(3)(\pot{5}{-3})=\SI{3}{\ampere\meter\squared};
\qquad
\tau_{\max}=mB=3(0{,}4)=\SI{1.2}{\newton\meter}.
\]
\emph{(Atenção à conversão: $\SI{50}{\centi\meter\squared}
=\pot{5}{-3}\,\si{\meter\squared}$, e não $\pot{5}{-2}$.)}\\
(b) O torque varia como $\sin\theta$ ao longo da volta, e seu valor médio sobre
meio ciclo é $\dfrac{2}{\pi}\tau_{\max}=\SI{0.764}{\newton\meter}$.
\[
\omega=1200\ \text{rpm}=\frac{1200(2\pi)}{60}=\SI{125.7}{\radian\per\second};
\]
\[
P=\tau_{\text{méd}}\,\omega=0{,}764(125{,}7)=\SI{96}{\watt}.
\]
\textbf{Por que o comutador é necessário:} sem inverter a corrente a cada meia
volta, o torque mudaria de sinal e a espira apenas oscilaria em torno da posição
de alinhamento.\\
\textbf{Por que motores reais usam muitas bobinas:} com uma só, o torque
\emph{pulsa} entre zero e o máximo. Várias bobinas defasadas produzem torque
quase constante --- e é isso que se busca em qualquer máquina prática.}
\end{questao}

\begin{questao}[3]{}
Uma espira de corrente é colocada em um campo magnético \textbf{não uniforme},
mais intenso de um lado.
\begin{itensq}
\item Determine se há força resultante.
\item Determine o sentido dessa força.
\item Relacione com o comportamento de ímãs.
\end{itensq}
\resposta{Há força resultante, dirigida para a região de campo mais intenso (se
$\vec m$ estiver alinhado com $\vec B$)}
\resolucao{(a) \textbf{Sim.} Em campo \emph{uniforme}, as forças nos lados
opostos da espira se cancelam exatamente e resta apenas torque. Em campo
\emph{não uniforme}, os lados estão em campos de módulos diferentes --- e o
cancelamento é incompleto.\\
(b) \textbf{Sentido.} A força resultante é dada pelo gradiente:
\[
\vec F=\nabla\left(\vec m\cdot\vec B\right).
\]
Se $\vec m$ estiver \emph{alinhado} com $\vec B$, a espira é atraída para a
região de campo mais intenso. Se estiver \emph{antialinhado}, é repelida.\\
(c) \textbf{Relação com ímãs.} Um ímã permanente é, essencialmente, um conjunto
de espiras microscópicas (correntes de magnetização) com momentos alinhados.
Daí:
\begin{itemize}
\item \textbf{Dois ímãs se atraem ou repelem} conforme a orientação relativa dos
momentos --- exatamente o comportamento da espira.
\item \textbf{Um ímã atrai ferro não magnetizado} porque primeiro o magnetiza
(alinhando seus domínios com o campo aplicado) e depois o puxa para a região de
campo mais forte. A atração é sempre no sentido do gradiente.
\item \textbf{Em campo uniforme, um ímã não é atraído} --- ele apenas gira até
se alinhar. É por isso que a agulha de uma bússola aponta, mas não é puxada em
direção ao norte.
\end{itemize}}
\end{questao}

\begin{questao}[3]{}
Compare a força entre dois fios paralelos conduzindo correntes iguais no mesmo
sentido e em sentidos opostos.
\begin{itensq}
\item Determine o sentido em cada caso.
\item Determine o módulo.
\item Explique a diferença em termos de campo.
\end{itensq}
\resposta{Mesmo sentido: atração; opostos: repulsão. Módulos iguais}
\resolucao{(a) e (b) O \textbf{módulo} é o mesmo nos dois casos:
\[
\frac{F}{L}=\frac{\mu_0I_1I_2}{2\pi d}.
\]
O que muda é o \emph{sentido}: correntes paralelas \textbf{atraem-se},
antiparalelas \textbf{repelem-se}.\\
(c) \textbf{Explicação pelo campo.} Considere o fio 2 imerso no campo do fio 1.
\begin{itemize}
\item \textbf{Mesmo sentido:} na posição do fio 2, o campo de 1 aponta em certo
sentido; aplicando $\vec F=I\vec\ell\times\vec B$, a força aponta \emph{para} o
fio 1.
\item \textbf{Sentidos opostos:} $\vec\ell$ inverte, e a força inverte junto.
\end{itemize}
\textbf{Leitura pelas linhas de campo.} Com correntes paralelas, o campo se
\emph{cancela} entre os fios e se \emph{reforça} fora --- as linhas envolvem o
par, e a "tensão" ao longo delas puxa os fios um contra o outro. Com correntes
opostas, ocorre o inverso: campo intenso entre os fios, que os afasta.\\
\textbf{Consequência prática:} em curtos-circuitos, barramentos paralelos
sofrem forças enormes. Com $\SI{20}{\kilo\ampere}$ a $\SI{10}{\centi\meter}$,
a força chega a $\SI{800}{\newton\per\meter}$ --- razão de existirem
espaçadores mecânicos robustos em painéis de potência.}
\end{questao}

\begin{questao}[3]{}
Um próton atravessa uma região com $E=\SI{2e4}{\volt\per\meter}$ e
$B=\SI{0.1}{\tesla}$, perpendiculares entre si e à velocidade, dispostos de
modo que as forças se oponham.
\begin{itensq}
\item Determine as duas parcelas da força de Lorentz para
      $v=\pot{1}{5}$, $\pot{2}{5}$ e $\pot{4}{5}\,\si{\meter\per\second}$.
\item Identifique qual parcela domina em cada caso.
\item Interprete o resultado.
\end{itensq}
\resposta{Domina a elétrica em $\pot{1}{5}$, equilíbrio em $\pot{2}{5}$ e a
magnética em $\pot{4}{5}\,\si{\meter\per\second}$}
\resolucao{(a) A parcela elétrica \textbf{não depende de $v$}:
\[
F_E=qE=(\pot{1.6}{-19})(\pot{2}{4})=\pot{3.2}{-15}\,\si{\newton}.
\]
A magnética é \textbf{proporcional a $v$}: $F_B=qvB$.
\begin{center}
\begin{tabular}{cccl}
\hline
$v$ ($\si{\meter\per\second}$) & $F_E$ ($\si{\newton}$)
 & $F_B$ ($\si{\newton}$) & Resultante\\
\hline
$\pot{1}{5}$ & $\pot{3.2}{-15}$ & $\pot{1.6}{-15}$ & elétrica domina\\
$\pot{2}{5}$ & $\pot{3.2}{-15}$ & $\pot{3.2}{-15}$ & \textbf{nula}\\
$\pot{4}{5}$ & $\pot{3.2}{-15}$ & $\pot{6.4}{-15}$ & magnética domina\\
\hline
\end{tabular}
\end{center}
(b) e (c) \textbf{Interpretação.} O equilíbrio ocorre em
\[
v=\frac{E}{B}=\frac{\pot{2}{4}}{0{,}1}=\pot{2}{5}\,\si{\meter\per\second},
\]
e apenas nessa velocidade a partícula segue reta.\\
\textbf{As três situações da tabela descrevem o mesmo dispositivo}: partículas
lentas são empurradas para um lado pela força elétrica, rápidas para o outro
pela magnética, e só as de velocidade $E/B$ atravessam sem desvio. É a
\emph{seleção de velocidades}, vista agora como competição entre as duas
parcelas de $\vec F=q(\vec E+\vec v\times\vec B)$.\\
\textbf{Observe a estrutura:} uma parcela constante e outra linear em $v$. A
igualdade só pode ocorrer em \emph{um} valor de velocidade --- e é isso que dá
seletividade ao instrumento.}
\end{questao}

\begin{questao}[3]{}
Uma partícula carregada move-se em uma região com campos $\vec E$ e $\vec B$
perpendiculares entre si, mas com velocidade \emph{diferente} de $E/B$.
\begin{itensq}
\item Descreva qualitativamente o movimento.
\item Determine a velocidade média de deriva.
\item Indique uma aplicação.
\end{itensq}
\resposta{Movimento cicloidal, com deriva média
$\vec v_d=\dfrac{\vec E\times\vec B}{B^{2}}$}
\resolucao{(a) \textbf{Movimento.} A partícula descreve uma trajetória
\textbf{cicloidal}: um movimento circular (devido a $\vec B$) superposto a um
deslocamento uniforme. Se ela parte do repouso, a curva é uma cicloide clássica
--- como um ponto na borda de uma roda que rola.\\
(b) \textbf{Velocidade de deriva.} O centro do círculo desloca-se com
\[
\vec v_d=\frac{\vec E\times\vec B}{B^{2}},
\qquad
|v_d|=\frac{E}{B},
\]
perpendicular a \emph{ambos} os campos.\\
\textbf{Duas propriedades notáveis:} a deriva \textbf{não depende da carga}
(íons e elétrons derivam no mesmo sentido) nem da \textbf{massa}.\\
(c) \textbf{Aplicações.}
\begin{itemize}
\item \textbf{Magnetrons} (fornos de micro-ondas): elétrons em campos cruzados
descrevem trajetórias cicloidais que excitam cavidades ressonantes.
\item \textbf{Confinamento de plasma:} a deriva $\vec E\times\vec B$ é um dos
mecanismos que governam o transporte em tokamaks --- e como ela não separa
cargas, não gera corrente própria.
\item \textbf{Magnetosfera terrestre:} explica parte do movimento de partículas
aprisionadas.
\end{itemize}}
\end{questao}

\begin{questao}[3]{}
Um condutor de $\SI{50}{\centi\meter}$ e $\SI{100}{\gram}$ está apoiado sobre
dois trilhos horizontais em campo vertical de $\SI{0.8}{\tesla}$. O
coeficiente de atrito estático é $0{,}25$.
\begin{itensq}
\item Determine a corrente mínima para iniciar o movimento.
\item Determine a aceleração inicial se a corrente for o dobro, com atrito
      cinético de $0{,}20$.
\end{itensq}
\resposta{$I=\SI{0.613}{\ampere}$; $a\approx\SI{2.9}{\meter\per\second\squared}$}
\resolucao{(a) A força magnética deve vencer o atrito estático:
\[
BIL=\mu_e mg
\;\Longrightarrow\;
I=\frac{0{,}25(0{,}100)(9{,}8)}{(0{,}8)(0{,}50)}=\SI{0.613}{\ampere}.
\]
(b) Com $I=\SI{1.226}{\ampere}$:
\[
F=BIL=(0{,}8)(1{,}226)(0{,}50)=\SI{0.490}{\newton};
\]
\[
f_c=\mu_c mg=0{,}20(0{,}100)(9{,}8)=\SI{0.196}{\newton};
\]
\[
a=\frac{0{,}490-0{,}196}{0{,}100}=\SI{2.94}{\meter\per\second\squared}.
\]
\textbf{Este é o princípio do canhão eletromagnético (\emph{railgun})} --- e
também do motor linear. \\
\textbf{Ressalva importante:} conforme o condutor acelera, ele passa a induzir
f.e.m. contrária ($\varepsilon=BLv$), que reduz a corrente. A aceleração
calculada vale apenas no \emph{instante inicial}; o movimento tende a uma
velocidade-limite em que a f.e.m. induzida equilibra a fonte.}
\end{questao}

\blocodesafio

\begin{questao}[4]{}
\textbf{(Demonstração --- trabalho nulo.)} Prove que a força magnética nunca
realiza trabalho sobre uma carga puntiforme.
\begin{itensq}
\item Demonstre formalmente.
\item Concilie com o fato de que motores elétricos realizam trabalho.
\item Explique o que ocorre com a energia em um motor.
\end{itensq}
\resposta{$\vec F\cdot\vec v=0$ sempre; o trabalho no motor vem da fonte, não do
campo}
\resolucao{(a) A potência é
\[
P=\vec F\cdot\vec v=q\left(\vec v\times\vec B\right)\cdot\vec v.
\]
O produto misto $(\vec v\times\vec B)\cdot\vec v$ é \textbf{identicamente
nulo}, pois $\vec v\times\vec B$ é perpendicular a $\vec v$ por construção. Logo
$P=0$ e $\Delta E_c=0$: o \emph{módulo} da velocidade nunca muda.\\
(b) \textbf{O aparente paradoxo do motor.} Um motor claramente realiza trabalho
--- ele levanta cargas e move veículos. Como conciliar?\\
\textbf{Resolução:} a força magnética atua sobre os \emph{portadores} do
condutor, e é perpendicular à velocidade \emph{deles}. Mas o fio se move, e a
velocidade total de um portador é a soma da deriva ao longo do fio com a
velocidade do fio.\\
A força magnética \emph{redistribui} energia: ela empurra o fio (realizando
trabalho positivo sobre ele) e, simultaneamente, freia a deriva dos portadores
(trabalho negativo de mesmo módulo). A soma continua sendo zero.\\
(c) \textbf{De onde vem a energia.} O freio à deriva manifesta-se como
\textbf{f.e.m. contraeletromotriz}, e é a \emph{fonte} que precisa vencê-la:
\[
P_{fonte}=\varepsilon_{fem}\,I=P_{\text{mecânica}}+P_{Joule}.
\]
\textbf{O campo magnético é um intermediário, não uma fonte.} Ele transfere
energia da fonte elétrica para o eixo, sem fornecer nem armazenar nada
líquido.\\
\emph{(É o mesmo papel de uma polia em mecânica: ela redireciona força sem criar
energia.)}}
\end{questao}

\begin{questao}[4]{}
\textbf{(Dedução do torque.)} Deduza a expressão do torque sobre uma espira
retangular em campo uniforme.
\begin{itensq}
\item Analise as forças em cada lado.
\item Obtenha $\tau=NIAB\sin\theta$.
\item Generalize para espira de forma qualquer.
\end{itensq}
\resposta{$\vec\tau=\vec m\times\vec B$, com $\vec m=NI\vec A$}
\resolucao{(a) Seja a espira de lados $a$ e $b$, com $\vec B$ no plano
horizontal e o eixo de rotação vertical. Chamando $\theta$ o ângulo entre
$\vec B$ e a \emph{normal} à espira:
\begin{itemize}
\item Os lados de comprimento $a$ \emph{paralelos} ao eixo sofrem forças
$F=BIa$, perpendiculares tanto ao fio quanto ao campo, em sentidos opostos.
Elas formam um \textbf{binário}.
\item Os lados de comprimento $b$ sofrem forças que se cancelam e cuja linha de
ação passa pelo eixo --- não produzem torque.
\end{itemize}
(b) O braço de alavanca de cada força do binário é
$\dfrac{b}{2}\sin\theta$:
\[
\tau=2\left(BIa\right)\frac{b}{2}\sin\theta=BI(ab)\sin\theta=BIA\sin\theta.
\]
Com $N$ espiras, $\tau=NIAB\sin\theta$.\\
(c) \textbf{Generalização.} Definindo o momento magnético
$\vec m=NI\vec A$ (com $\vec A$ normal à espira, sentido dado pela regra da mão
direita):
\[
\boxed{\vec\tau=\vec m\times\vec B}
\]
Essa forma vale para espira de \textbf{qualquer} formato --- basta decompô-la em
retângulos infinitesimais, cujas correntes internas se cancelam aos pares.\\
\textbf{Analogia formal com o dipolo elétrico:} $\vec\tau=\vec p\times\vec E$.
As duas expressões são idênticas em estrutura, e por isso a espira é
literalmente um \emph{dipolo magnético}.}
\end{questao}

\begin{questao}[4]{}
\textbf{(Projeto de espectrômetro.)} Projete um espectrômetro capaz de separar
isótopos com $\Delta m/m=1\%$.
\begin{itensq}
\item Determine a separação espacial em função dos parâmetros.
\item Estabeleça o requisito de colimação do feixe.
\item Discuta as limitações práticas.
\end{itensq}
\resposta{$\Delta r/r=\Delta m/m$; a resolução exige feixe com dispersão angular
e de velocidade menores que $1\%$}
\resolucao{(a) De $r=\dfrac{mv}{qB}$ com $v$ e $B$ fixos:
\[
\frac{\Delta r}{r}=\frac{\Delta m}{m}=1\%.
\]
Com $r=\SI{20}{\centi\meter}$, a separação de raios é
$\SI{2}{\milli\meter}$, e no detector (após meia volta) a separação é
$2\Delta r=\SI{4}{\milli\meter}$.\\
(b) \textbf{Requisito de colimação.} Para \emph{resolver} duas linhas separadas
por $\SI{4}{\milli\meter}$, a largura de cada linha deve ser bem menor. Duas
fontes de alargamento:
\begin{itemize}
\item \textbf{Dispersão de velocidade:} como $r\propto v$, uma dispersão
$\Delta v/v$ produz alargamento $\Delta r/r$ de mesma magnitude. É preciso
$\Delta v/v\ll1\%$ --- daí a necessidade do seletor de velocidades.
\item \textbf{Dispersão angular:} íons entrando com pequenos ângulos diferentes
descrevem círculos deslocados. \emph{Felizmente}, a geometria de meia volta é
\textbf{autofocalizante}: íons que divergem de um ponto reconvergem após
$\ang{180}$, em primeira ordem. O alargamento residual é de segunda ordem no
ângulo.
\end{itemize}
(c) \textbf{Limitações práticas.}
\begin{itemize}
\item \textbf{Vácuo:} colisões com gás residual espalham o feixe. É preciso
pressão abaixo de $\pot{1}{-5}\,\si{\pascal}$.
\item \textbf{Uniformidade do campo:} variações de $B$ ao longo da trajetória
introduzem erro proporcional. Exige-se uniformidade melhor que a resolução
desejada.
\item \textbf{Carga espacial:} feixes intensos se repelem eletrostaticamente e
se alargam --- limitando a corrente utilizável.
\item \textbf{Íons de mesma razão $m/q$:} não são separáveis. $\mathrm{N_2^+}$ e
$\mathrm{CO^+}$ (ambos $28\,u$) exigem instrumentos de altíssima resolução para
distinguir.
\end{itemize}}
\end{questao}

\begin{questao}[4]{}
\textbf{(Projeto de atuador linear.)} Projete um atuador que produza
$\SI{5}{\newton}$ com curso de $\SI{2}{\centi\meter}$.
\begin{itensq}
\item Determine as combinações de $B$, $I$ e $L$ possíveis.
\item Avalie o compromisso entre corrente e número de espiras.
\item Determine a potência dissipada e o trabalho realizado.
\end{itensq}
\resposta{Com $B=\SI{0.5}{\tesla}$ e $L=\SI{0.2}{\meter}$: $I=\SI{50}{\ampere}$
com uma espira, ou $\SI{2.5}{\ampere}$ com vinte}
\resolucao{(a) De $F=NBIL$:
\[
NIL=\frac{F}{B}=\frac{5}{0{,}5}=\SI{10}{\ampere\meter}.
\]
Qualquer combinação que satisfaça esse produto serve --- há dois graus de
liberdade.\\
(b) \textbf{Compromisso.} Com $L=\SI{0.2}{\meter}$ fixo:
\begin{center}
\begin{tabular}{ccc}
\hline
$N$ & $I$ & Observação\\
\hline
$1$ & $\SI{50}{\ampere}$ & fio muito grosso, fonte de alta corrente\\
$20$ & $\SI{2.5}{\ampere}$ & equilíbrio razoável\\
$200$ & $\SI{0.25}{\ampere}$ & bobina volumosa, alta indutância\\
\hline
\end{tabular}
\end{center}
\textbf{Mais espiras} reduzem a corrente (fio mais fino, fonte mais simples),
mas aumentam volume, massa, resistência e \emph{indutância} --- o que torna o
atuador mais lento a responder.\\
\textbf{Menos espiras} dão resposta rápida, ao custo de corrente elevada.\\
(c) \textbf{Dissipação.} Com $N=20$, $I=\SI{2.5}{\ampere}$ e resistência total
estimada em $\SI{2}{\ohm}$:
\[
P=RI^{2}=2(6{,}25)=\SI{12.5}{\watt}.
\]
\textbf{Trabalho útil:}
\[
W=F\,d=5(0{,}02)=\SI{0.1}{\joule}.
\]
\textbf{Comparação reveladora:} se o curso levar $\SI{0.1}{\second}$, a potência
mecânica é $\SI{1}{\watt}$ contra $\SI{12.5}{\watt}$ dissipados --- rendimento
de $7\%$. Atuadores eletromagnéticos de curso curto são intrinsecamente
ineficientes, e é por isso que costumam operar em regime \emph{pulsado}, e não
contínuo.}
\end{questao}

\begin{questao}[4]{}
\textbf{(Efeito Hall.)} Deduza a expressão da tensão Hall e analise suas
aplicações.
\begin{itensq}
\item Deduza $V_H=\dfrac{IB}{nqt}$.
\item Explique como o sinal revela o tipo de portador.
\item Cite três aplicações e o que cada uma mede.
\end{itensq}
\resposta{Equilíbrio entre força magnética e força elétrica transversal}
\resolucao{(a) \textbf{Dedução.} Os portadores, com velocidade de deriva
$v_d$, sofrem força magnética $qv_dB$ transversal e acumulam-se em uma face.
Isso cria campo elétrico transversal $E_H$, que cresce até equilibrar:
\[
qE_H=qv_dB
\;\Longrightarrow\;
E_H=v_dB.
\]
Com largura $w$, a tensão é $V_H=E_Hw=v_dBw$. Usando
$I=nqv_d(wt)$, elimina-se $v_d$:
\[
v_d=\frac{I}{nqwt}
\;\Longrightarrow\;
V_H=\frac{IB}{nqt}.
\]
\textbf{Note:} a largura $w$ cancela --- só a \emph{espessura} $t$ importa.\\
(b) \textbf{Sinal e tipo de portador.} Elétrons (negativos) e lacunas
(positivas) movem-se em sentidos opostos para produzir a mesma corrente. A força
$q\vec v\times\vec B$ os desvia para o \textbf{mesmo} lado --- mas com cargas de
sinais contrários, a polaridade de $V_H$ \textbf{se inverte}.\\
\textbf{É o único método simples de determinar experimentalmente o sinal dos
portadores} --- e foi assim que se estabeleceu que semicondutores tipo P
conduzem por lacunas.\\
(c) \textbf{Três aplicações.}
\begin{itemize}
\item \textbf{Medição de campo magnético} (gaussímetro): com $I$, $n$ e
$t$ conhecidos, $V_H\propto B$.
\item \textbf{Medição de corrente sem contato:} coloca-se o sensor no entreferro
de um núcleo que envolve o condutor. Mede $B$, e daí $I$ --- com isolação
galvânica completa.
\item \textbf{Caracterização de materiais:} com $B$ e $I$ conhecidos, obtém-se
$n$ (densidade de portadores) e o \emph{sinal} deles. Combinado com a
condutividade, fornece a mobilidade.
\end{itemize}
\textbf{Vantagem sobre alternativas:} sensores Hall respondem a campo
\emph{estático}, ao contrário de bobinas, que só detectam variação.}
\end{questao}

\begin{questao}[4]{}
\textbf{(Estabilidade de espira.)} Analise o equilíbrio de uma espira de
corrente em campo magnético uniforme.
\begin{itensq}
\item Determine as posições de equilíbrio.
\item Classifique cada uma quanto à estabilidade.
\item Deduza o período de pequenas oscilações.
\end{itensq}
\resposta{$\theta=0$ é estável, $\theta=\pi$ é instável;
$T=2\pi\sqrt{J/(mB)}$}
\resolucao{(a) O torque é $\tau=-mB\sin\theta$, com $\theta$ o ângulo entre
$\vec m$ e $\vec B$. Ele se anula em
\[
\theta=0 \quad\text{e}\quad \theta=\pi.
\]
(b) \textbf{Classificação pela energia.} A energia potencial de um dipolo
magnético é
\[
U=-\vec m\cdot\vec B=-mB\cos\theta.
\]
\begin{itemize}
\item $\theta=0$ ($\vec m$ paralelo a $\vec B$): $U=-mB$, \textbf{mínimo}
$\Rightarrow$ equilíbrio \textbf{estável}. Deslocada, a espira sofre torque
restaurador.
\item $\theta=\pi$ (antiparalelo): $U=+mB$, \textbf{máximo} $\Rightarrow$
equilíbrio \textbf{instável}. Qualquer perturbação a faz girar $\ang{180}$.
\end{itemize}
(c) \textbf{Pequenas oscilações.} Para $\theta\ll1$, $\sin\theta\approx\theta$ e
\[
J\ddot\theta=-mB\theta
\;\Longrightarrow\;
\ddot\theta+\frac{mB}{J}\theta=0,
\]
equação do MHS com
\[
\omega=\sqrt{\frac{mB}{J}};
\qquad
T=2\pi\sqrt{\frac{J}{mB}},
\]
onde $J$ é o momento de inércia da espira.\\
\textbf{Aplicação prática:} é assim que funcionam o \textbf{galvanômetro de
bobina móvel} --- em que uma mola acrescenta torque restaurador proporcional a
$\theta$, tornando a deflexão proporcional à corrente --- e a \textbf{bússola},
cujo período de oscilação permite medir o campo terrestre se $J$ e $m$ forem
conhecidos.\\
\textbf{Diferença de $U$ entre os dois equilíbrios:} $2mB$ --- energia
necessária para inverter a espira, e a mesma que um ímã precisa receber para ter
sua polarização revertida.}
\end{questao}

\begin{questao}[4]{}
\textbf{(Força de Lorentz --- estrutura e consequências.)} Analise a expressão
completa $\vec F=q\left(\vec E+\vec v\times\vec B\right)$.
\begin{itensq}
\item Compare as duas parcelas quanto à dependência da velocidade, à direção e
      ao trabalho realizado.
\item Avalie um próton com $v=\pot{3}{4}\,\si{\meter\per\second}$ em
      $E=\SI{500}{\volt\per\meter}$ e $B=\SI{0.02}{\tesla}$
      \emph{perpendiculares entre si}, com $\vec v$ paralelo a $\vec E$.
\item Explique por que aceleradores usam campo elétrico para acelerar e
      magnético para guiar.
\end{itensq}
\resposta{$F_E=\SI{8e-17}{\newton}$, $F_B=\SI{9.6e-17}{\newton}$,
perpendiculares; $|F|=\SI{1.25e-16}{\newton}$ a $\ang{50.2}$ de $\vec E$}
\resolucao{(a) \textbf{Comparação das duas parcelas.}
\begin{center}
\begin{tabular}{lll}
\hline
& Parcela elétrica & Parcela magnética\\
\hline
Expressão & $q\vec E$ & $q\,\vec v\times\vec B$\\
Depende de $v$? & \textbf{não} & \textbf{sim} (linear)\\
Age em repouso? & sim & não\\
Direção & paralela a $\vec E$ & $\perp$ a $\vec v$ \emph{e} a $\vec B$\\
Realiza trabalho? & \textbf{sim} & \textbf{nunca}\\
Altera $|v|$? & sim & não\\
\hline
\end{tabular}
\end{center}
\textbf{A linha decisiva é a do trabalho.} Como
$\left(\vec v\times\vec B\right)\cdot\vec v=0$ identicamente, a parcela
magnética tem potência nula --- ela desvia sem acelerar. A elétrica, ao
contrário, é a única capaz de variar a energia cinética.\\
(b) \textbf{Avaliação numérica.}
\[
F_E=qE=(\pot{1.6}{-19})(500)=\pot{8}{-17}\,\si{\newton};
\]
\[
F_B=qvB=(\pot{1.6}{-19})(\pot{3}{4})(0{,}02)=\pot{9.6}{-17}\,\si{\newton}.
\]
Com $\vec v$ paralelo a $\vec E$ e $\vec B$ perpendicular a ambos, a força
magnética é \textbf{perpendicular} à elétrica:
\[
|\vec F|=\sqrt{F_E^{2}+F_B^{2}}=\pot{1.25}{-16}\,\si{\newton},
\qquad
\theta=\arctan\frac{9{,}6}{8{,}0}=\ang{50.2}
\]
em relação a $\vec E$.\\
\textbf{Razão entre as parcelas:}
$\dfrac{F_B}{F_E}=\dfrac{vB}{E}=1{,}20$ --- adimensional, e igual a $1$
exatamente na velocidade $E/B$ do seletor.\\
\textbf{Trabalho em um deslocamento de $\SI{1}{\centi\meter}$ ao longo de
$\vec E$:} apenas a parcela elétrica contribui,
\[
W=F_E d=(\pot{8}{-17})(0{,}01)=\pot{8}{-19}\,\si{\joule}=\SI{5}{\electronvolt},
\]
que é simplesmente $qEd$ --- ou seja, $q$ vezes a tensão percorrida.\\
(c) \textbf{Divisão de papéis em aceleradores.}
\begin{itemize}
\item \textbf{Campo elétrico acelera.} É o único que realiza trabalho, e o ganho
de energia por passagem é $qU$ --- daí as cavidades de radiofrequência.
\item \textbf{Campo magnético guia.} Ele curva a trajetória sem custo
energético, permitindo fechar o percurso em círculo e reaproveitar as mesmas
cavidades milhares de vezes.
\item \textbf{Campo magnético também focaliza}, com quadrupolos que comprimem o
feixe transversalmente --- de novo, sem alterar a energia.
\end{itemize}
\textbf{Se fosse possível acelerar com campo magnético}, aceleradores lineares
de quilômetros seriam desnecessários. A impossibilidade não é técnica: é
consequência direta de $\left(\vec v\times\vec B\right)\perp\vec v$.\\
\textbf{Nota histórica:} a expressão completa reúne dois fenômenos que foram
descobertos e estudados separadamente. Sua unificação --- e a constatação de que
o que é campo elétrico em um referencial pode ser magnético em outro --- é um
dos pontos de partida da relatividade restrita.}
\end{questao}

\begin{questao}[4]{}
\textbf{(Projeto de cíclotron.)} Um cíclotron acelera prótons até
$\SI{10}{\mega\electronvolt}$ usando campo de $\SI{1.5}{\tesla}$.
\begin{itensq}
\item Determine a frequência de operação.
\item Determine o raio máximo e o número de voltas, com
      $\SI{50}{\kilo\volt}$ de ganho por passagem.
\item Determine se a correção relativística é necessária.
\end{itensq}
\dados{$e=\pot{1.6}{-19}\,\si{\coulomb}$; $m_p=\pot{1.67}{-27}\,\si{\kilogram}$;
$\SI{1}{\mega\electronvolt}=\pot{1.6}{-13}\,\si{\joule}$}
\resposta{$f=\SI{22.9}{\mega\hertz}$; $r=\SI{30.5}{\centi\meter}$;
$100$ voltas; correção de $1{,}1\%$}
\resolucao{(a) A frequência de cíclotron independe da energia:
\[
f=\frac{qB}{2\pi m}=\frac{(\pot{1.6}{-19})(1{,}5)}{2\pi(\pot{1.67}{-27})}
=\SI{22.9}{\mega\hertz}.
\]
\textbf{É isso que torna o aparelho viável:} o oscilador que alimenta os "dês"
opera em frequência \emph{fixa}, embora a partícula ganhe energia
continuamente.\\
(b) Da energia final,
\[
v=\sqrt{\frac{2E}{m}}
=\sqrt{\frac{2(10)(\pot{1.6}{-13})}{\pot{1.67}{-27}}}
=\pot{4.38}{7}\,\si{\meter\per\second};
\]
\[
r=\frac{mv}{qB}=\SI{0.305}{\meter}
\;\Longrightarrow\;
\text{diâmetro}\approx\SI{61}{\centi\meter}.
\]
Com dois cruzamentos do intervalo por volta, cada volta rende
$\SI{100}{\kilo\electronvolt}$:
\[
N=\frac{10\,\si{\mega\electronvolt}}{0{,}1\,\si{\mega\electronvolt}}
=100\ \text{voltas}.
\]
\textbf{Tempo total:} $100/f=\SI{4.4}{\micro\second}$.\\
(c) \textbf{Correção relativística.} Com
$v/c=0{,}146$:
\[
\gamma=\frac{1}{\sqrt{1-(v/c)^{2}}}=1{,}0108,
\]
ou seja, a massa efetiva é $1{,}1\%$ maior no fim da aceleração.\\
\textbf{Consequência:} a frequência de cíclotron cai na mesma proporção, e a
partícula começa a chegar \emph{atrasada} ao intervalo acelerador. Após muitas
voltas o defasamento acumula e a aceleração cessa.\\
\textbf{Aqui, $1{,}1\%$ é tolerável} --- o cíclotron clássico funciona. Mas para
$\SI{100}{\mega\electronvolt}$ o desvio passaria de $10\%$ e seria proibitivo.\\
\textbf{Soluções:} o \textbf{sincrociclotron} modula a frequência ao longo do
ciclo; o \textbf{síncrotron} mantém o raio fixo e varia o campo. Ambos abandonam
a simplicidade do cíclotron justamente porque o período deixa de ser constante
--- e a raiz disso está na expressão $T=2\pi m/qB$, válida apenas enquanto
$m$ o for.}
\end{questao}

\begin{questao}[4]{}
\textbf{(Força entre correntes e o limite de barramentos.)} Um barramento de
subestação tem dois condutores paralelos separados por $\SI{15}{\centi\meter}$,
que devem suportar um curto-circuito de $\SI{40}{\kilo\ampere}$.
\begin{itensq}
\item Determine a força por metro durante o curto.
\item Determine o espaçamento máximo entre suportes, para deflexão admissível.
\item Discuta o caráter dessa solicitação.
\end{itensq}
\resposta{$F/L\approx\SI{2.1}{\kilo\newton\per\meter}$ --- comparável ao peso
de uma pessoa por metro}
\resolucao{(a)
\[
\frac{F}{L}=\frac{\mu_0I^{2}}{2\pi d}
=\pot{2}{-7}\frac{(\pot{4}{4})^{2}}{0{,}15}
=\pot{2}{-7}\frac{\pot{1.6}{9}}{0{,}15}=\SI{2133}{\newton\per\meter}.
\]
Mais de duas toneladas-força por metro de barramento --- e \textbf{repulsiva},
se as correntes forem opostas (que é o caso em um curto entre fases).\\
(b) \textbf{Espaçamento entre suportes.} Tratando o barramento como viga
biapoiada com carga distribuída $w=F/L$, a flecha é
\[
\delta=\frac{5wL_s^{4}}{384EJ}.
\]
Para deflexão admissível $\delta$, o vão máximo cresce apenas com
$L_s\propto\delta^{1/4}$ --- \textbf{dependência muito fraca}. Dobrar a flecha
admissível permite aumentar o vão em apenas $19\%$.\\
Na prática, os suportes ficam a cada $\SI{0.5}{}$ a $\SI{1}{\meter}$.\\
(c) \textbf{Caráter da solicitação.}
\begin{itemize}
\item \textbf{É transitória} --- dura o tempo de atuação da proteção, tipicamente
$\SI{100}{\milli\second}$. Mas a força \emph{de pico} pode ser o dobro da
calculada, devido à componente assimétrica da corrente de curto.
\item \textbf{Vai com $I^{2}$} --- dobrar a corrente de curto quadruplica a
força. É por isso que a corrente de curto-circuito presumida é um dado de
projeto tão crítico.
\item \textbf{É dinâmica:} se a frequência natural da estrutura estiver próxima
de $\SI{120}{\hertz}$ (dobro da rede, pois $F\propto i^{2}$), pode haver
ressonância mecânica.
\end{itemize}
\textbf{Conclusão:} em regime normal ($\SI{1}{\kilo\ampere}$, digamos), a força
seria $\SI{1.3}{\newton\per\meter}$ --- desprezível. O dimensionamento mecânico
de barramentos é governado inteiramente pela condição de \emph{falta}.}
\end{questao}

% END BANCO COMPLEMENTAR

\gabarito[12.4 Força magnética]

% =====================================================================
\subsection{Lei de Lenz e indução eletromagnética}
% =====================================================================

\blocofixacao

\begin{questao}[1]{}
Segundo a lei de Lenz, a corrente induzida em uma espira tem sentido tal que:
\begin{alternativas}
\item reforça a variação do fluxo que a originou.
\item se opõe à variação do fluxo que a originou.
\item é sempre horária.
\item independe do sentido da variação do fluxo.
\item é proporcional ao fluxo, e não à sua variação.
\end{alternativas}
\resposta{(b)}
\resolucao{A lei de Lenz é uma manifestação da conservação de energia: a corrente
induzida cria um campo que \emph{se opõe} à variação de fluxo que a gerou. Se
reforçasse a variação (alternativa a), teríamos energia criada do nada.}
\end{questao}

\begin{questao}[1]{}
Um campo magnético variável atravessa perpendicularmente uma espira, com fluxo
$\Phi(t) = 0{,}02\,t$ (em Wb, com $t$ em segundos). Determine a fem induzida.
\resposta{$\varepsilon = \SI{20}{\milli\volt}$}
\resolucao{A fem é a taxa de variação do fluxo:
\[
\varepsilon=\left|\frac{\mathrm{d}\Phi}{\mathrm{d}t}\right|
=\SI{0.02}{\weber\per\second}=\SI{20}{\milli\volt}.
\]
Como o fluxo cresce linearmente, a fem é \emph{constante}.}
\end{questao}

\blocoaplicacao

\begin{questao}[2]{}
O campo magnético no interior de uma bobina aumenta linearmente de $0$ a
$\SI{0.5}{\tesla}$ em $\SI{2}{\second}$. A bobina tem $\SI{50}{}$ espiras e área de
$\SI{20}{\centi\meter\squared}$. Calcule a fem induzida.
\resposta{$\varepsilon = \SI{25}{\milli\volt}$}
\resolucao{Com $N = 50$ espiras e
$A = \pot{20}{-4}\,\si{\meter\squared}$:
\[
\varepsilon = N\,A\,\frac{\Delta B}{\Delta t}
= 50\cdot\pot{20}{-4}\cdot\frac{0{,}5}{2}
= 50\cdot\pot{20}{-4}\cdot0{,}25 = \SI{25}{\milli\volt}.
\]}
\end{questao}

\begin{questao}[2]{}
Uma bobina de $\SI{100}{}$ espiras e área $\SI{30}{\centi\meter\squared}$ está em
um campo que varia segundo $B(t) = 0{,}5 - 0{,}1\,t$ (em T). Calcule a fem induzida
e diga se a corrente induzida cresce, decresce ou é constante.
\resposta{$\varepsilon = \SI{30}{\milli\volt}$, constante}
\resolucao{A taxa de variação do campo é
$\dfrac{\mathrm{d}B}{\mathrm{d}t} = \SI{-0.1}{\tesla\per\second}$ (constante):
\[
|\varepsilon| = N\,A\,\left|\frac{\mathrm{d}B}{\mathrm{d}t}\right|
= 100\cdot\pot{30}{-4}\cdot0{,}1 = \SI{30}{\milli\volt}.
\]
Como a taxa de variação é constante, a fem é constante --- e, num circuito de
resistência fixa, a corrente induzida também é \emph{constante}. O fluxo está
diminuindo, então a corrente induzida atua no sentido de \emph{reforçá-lo}
(Lenz).}
\end{questao}

\begin{questao}[2]{}
Um ímã é aproximado, com o polo norte à frente, de uma espira condutora, como na
figura. Determine o sentido da corrente induzida e a natureza da força entre o
ímã e a espira.

\circuitoq{%
\begin{tikzpicture}[scale=0.9,line width=0.8pt,>=Stealth]
 % ima
 \draw[fill=Icol!20,draw=Icol] (-4,-0.35) rectangle (-2.4,0.35);
 \node at (-3.5,0){\footnotesize S}; \node at (-2.7,0){\footnotesize N};
 \draw[->,Vcol,thick] (-2.2,0)--(-1.2,0); \node[Vcol,font=\footnotesize] at (-1.7,0.3){$v$};
 % espira
 \draw[azulprof,very thick] (0,0) ellipse (0.35 and 1.1);
 \node[azulprof,font=\footnotesize] at (0,-1.4){espira};
\end{tikzpicture}}{Ímã aproximando-se de uma espira.}

\begin{alternativas}
\item A corrente induzida cria um polo norte voltado para o ímã, repelindo-o.
\item A corrente induzida cria um polo sul voltado para o ímã, atraindo-o.
\item Não há corrente induzida, pois o ímã não toca a espira.
\item A corrente induzida não gera força sobre o ímã.
\item A espira é atraída, acelerando o ímã.
\end{alternativas}
\resposta{(a)}
\resolucao{Ao aproximar o polo norte, o fluxo através da espira \emph{aumenta}.
Pela lei de Lenz, a corrente induzida se opõe a esse aumento, criando um campo que
\emph{repele} o ímã --- ou seja, a face da espira voltada para o ímã se comporta
como um polo \textbf{norte}. Consequentemente, surge uma força de \emph{repulsão}
que dificulta a aproximação. É a manifestação da conservação de energia: para
aproximar o ímã é preciso realizar trabalho, que se converte na energia elétrica
da corrente induzida.}
\end{questao}

\begin{questao}[2]{}
Um transformador ideal tem no primário uma bobina onde o fluxo magnético varia à
taxa de $\SI{0.01}{\weber\per\second}$. Se o enrolamento tem $\SI{200}{}$ espiras,
qual é a fem induzida?
\resposta{$\varepsilon = \SI{2}{\volt}$}
\resolucao{$\varepsilon = N\dfrac{\mathrm{d}\Phi}{\mathrm{d}t}
= 200\cdot0{,}01 = \SI{2}{\volt}$. É assim que o transformador transfere energia:
o fluxo variável criado pelo primário induz tensão no secundário,
proporcionalmente ao número de espiras de cada lado.}
\end{questao}

\blocoaprofundamento

\begin{questao}[3]{}
O fluxo magnético em uma espira varia senoidalmente:
$\Phi(t) = \Phi_0\cos(\omega t)$, com $\Phi_0 = \SI{0.01}{\weber}$ e
$\omega = \SI{100}{\radian\per\second}$.
\begin{itensq}
\item Obtenha a expressão da fem induzida.
\item Qual é o valor máximo da fem?
\end{itensq}
\resposta{(a) $\varepsilon(t) = \Phi_0\omega\sin(\omega t)$;
(b) $\varepsilon_{\max} = \SI{1}{\volt}$}
\resolucao{(a) A fem é a derivada do fluxo (com sinal trocado):
\[
\varepsilon = -\frac{\mathrm{d}\Phi}{\mathrm{d}t}
= -\Phi_0\cdot(-\omega\sin(\omega t)) = \Phi_0\,\omega\sin(\omega t).
\]
(b) O valor máximo ocorre quando $\sin(\omega t) = 1$:
\[
\varepsilon_{\max} = \Phi_0\,\omega = 0{,}01\cdot100 = \SI{1}{\volt}.
\]
\textbf{Observação:} um fluxo cossenoidal gera uma fem \emph{senoidal} --- há uma
defasagem de $\SI{90}{\degree}$ entre eles. É exatamente o que ocorre em um
gerador de corrente alternada: a espira girando em campo uniforme produz fluxo
senoidal e, portanto, tensão senoidal.}
\end{questao}

\begin{questao}[3]{Apostila original revisada}
A corrente induzida em um circuito puramente resistivo é
$i(t)=I_0\sin(\omega t)$. Determine a expressão da fem induzida e indique em quais
instantes a corrente muda de sentido durante um período.
\resposta{$\mathcal{E}(t)=RI_0\sin(\omega t)$; inversões em $t=0$, $\pi/\omega$ e $2\pi/\omega$}
\resolucao{Como o circuito é resistivo, $\mathcal{E}=Ri$. O sinal da senoide muda a
cada múltiplo de $\pi$, portanto o sentido da corrente se inverte a cada meio período.}
\end{questao}

\blocodesafio

\begin{questao}[4]{}
Uma barra condutora de comprimento $\ell = \SI{40}{\centi\meter}$ desliza sem
atrito, com velocidade constante $v = \SI{3}{\meter\per\second}$, sobre dois
trilhos paralelos fechados por um resistor $R = \SI{2}{\ohm}$. O conjunto está
imerso em campo magnético uniforme $B = \SI{0.5}{\tesla}$, perpendicular ao plano
dos trilhos.

\circuitoq{%
\begin{tikzpicture}[scale=1,line width=0.8pt,>=Stealth]
 % trilhos
 \draw[cinza,line width=1.4pt] (0,0)--(6,0);
 \draw[cinza,line width=1.4pt] (0,2.2)--(6,2.2);
 % resistor
 \draw (0,0) to[R,l_={$R$}] (0,2.2);
 % barra
 \draw[Vcol,line width=2.2pt] (3.6,-0.15)--(3.6,2.35);
 \node[Vcol,font=\footnotesize,above] at (3.6,2.4) {barra};
 \draw[->,laranja,line width=1.3pt] (3.9,1.1)--(5.1,1.1)
       node[midway,above,font=\footnotesize]{$v$};
 % campo B saindo
 \foreach \x in {1.0,1.8,2.6}
   \foreach \y in {0.55,1.1,1.65}
     {\node[Icol,font=\scriptsize] at (\x,\y) {$\odot$};}
 \node[Icol,font=\footnotesize] at (1.8,0.15) {$\vec B$ (saindo)};
 \draw[<->,cinza] (6.3,0)--(6.3,2.2) node[midway,right,font=\footnotesize]{$\ell$};
\end{tikzpicture}}{Barra deslizante sobre trilhos --- fem de movimento.}

\begin{itensq}
\item Calcule a fem induzida e a corrente no circuito.
\item Determine a força magnética que atua sobre a barra e seu sentido.
\item Calcule a potência dissipada no resistor e a potência mecânica necessária
      para manter a barra em movimento uniforme. Compare-as.
\item Interprete o resultado à luz da lei de Lenz e da conservação de energia.
\end{itensq}
\resposta{(a) $\varepsilon=\SI{0.6}{\volt}$, $i=\SI{0.3}{\ampere}$;
(b) $F=\SI{0.06}{\newton}$, opondo-se ao movimento;
(c) ambas $=\SI{0.18}{\watt}$; (d) ver comentário}
\resolucao{\textbf{(a) Fem de movimento.} À medida que a barra avança, a área do
circuito aumenta a uma taxa $\dfrac{\mathrm{d}A}{\mathrm{d}t} = \ell v$. Logo
\[
\varepsilon=\left|\frac{\mathrm{d}\Phi}{\mathrm{d}t}\right|
= B\,\ell\,v = 0{,}5\cdot0{,}4\cdot3=\SI{0.6}{\volt},
\qquad
i=\frac{\varepsilon}{R}=\frac{0{,}6}{2}=\SI{0.3}{\ampere}.
\]
\textbf{(b) Força sobre a barra.} A barra, agora percorrida por corrente,
encontra-se no campo $B$ e sofre
\[
F = B\,i\,\ell = 0{,}5\cdot0{,}3\cdot0{,}4=\SI{0.06}{\newton}.
\]
Pela lei de Lenz, essa força \textbf{opõe-se ao movimento} --- ela freia a barra.
(Se favorecesse o movimento, teríamos aceleração espontânea e energia criada do
nada.)

\textbf{(c) Balanço de potência.}
\[
P_{\text{dissipada}} = R\,i^{2} = 2(0{,}3)^{2}=\SI{0.18}{\watt},
\]
\[
P_{\text{mecânica}} = F\,v = 0{,}06\cdot3 = \SI{0.18}{\watt}.
\]
\textbf{São exatamente iguais.}

\textbf{(d) Interpretação.} O resultado não é coincidência: é a
\textbf{conservação de energia} em ação. Para manter a barra em movimento
\emph{uniforme}, um agente externo deve aplicar uma força igual e oposta à força
magnética de frenagem, realizando trabalho a uma taxa $Fv$. Toda essa potência
mecânica é convertida em potência elétrica e, finalmente, dissipada como calor no
resistor. Nada se perde, nada se cria:
\[
\underbrace{P_{\text{mec}}}_{\text{trabalho externo}}
\;\longrightarrow\;
\underbrace{\varepsilon\,i}_{\text{potência elétrica}}
\;\longrightarrow\;
\underbrace{R i^{2}}_{\text{calor}}.
\]
Verificando pela via elétrica: $\varepsilon\,i = 0{,}6\cdot0{,}3
= \SI{0.18}{\watt}$ \checkmark\\
\textbf{Este é o princípio do gerador elétrico} --- e, invertendo o raciocínio, o
do \emph{freio eletromagnético} usado em trens e em esteiras de academia: quanto
maior a velocidade, maior a corrente induzida e maior a frenagem, sem contato
mecânico nem desgaste.}
\end{questao}

% BEGIN BANCO COMPLEMENTAR Lei de Lenz e indução eletromagnética
% =====================================================================
%  Banco complementar --- Lei de Lenz e Inducao Eletromagnetica
%  REVISADO. Substitui a versao gerada por template (8 clones em N1,
%  11 em N2, 13 em N3 e 9 questoes conceituais marcadas como N4).
%  O cap11 ja cobre: enunciado da lei de Lenz (MC), Phi(t)=0,02t,
%  campo crescendo linearmente de 0 a 0,5 T, bobina de 100 espiras com
%  B(t), IMA APROXIMADO COM O POLO NORTE (N2), transformador ideal com
%  dPhi/dt dado, fluxo senoidal, corrente induzida senoidal e -- como
%  N4 -- a BARRA DESLIZANDO SOBRE TRILHOS. As duas ultimas motivaram a
%  remocao de duas N4 do banco antigo, que as duplicavam.
%  Este banco explora: CARGA induzida (que independe do tempo), freio
%  por correntes parasitas, laminacao de nucleos, velocidade-limite,
%  inducao mutua com convencao dos pontos, projeto de bobina e de
%  experimento, e a razao pela qual transformadores nao operam em CC.
%  Distribuicao mantida: 8 N1 + 11 N2 + 13 N3 + 9 N4 = 41
% =====================================================================

\subsubsection*{Banco complementar --- Lei de Lenz e indução eletromagnética}

\begin{atencao}[Organização do banco]
A lei de Faraday é $\varepsilon=-N\dfrac{\mathrm{d}\Phi}{\mathrm{d}t}$: o que
induz f.e.m. é a \textbf{variação} do fluxo, não seu valor. O sinal negativo é a
lei de Lenz, e ela não é um detalhe --- é a conservação da energia em forma
elétrica. O N3 trata de velocidade-limite, carga induzida, correntes parasitas e
indução mútua; o N4 exige demonstrações, projeto e a análise de por que
transformadores exigem corrente alternada.
\end{atencao}

\blocofixacao

\begin{questao}[1]{}
Uma bobina de $200$ espiras sofre variação de fluxo de
$\SI{2}{\milli\weber}$ por espira em $\SI{5}{\milli\second}$. Determine a
f.e.m. média induzida.
\resposta{$\varepsilon=\SI{80}{\volt}$}
\resolucao{
\[
|\varepsilon|=N\frac{\Delta\Phi}{\Delta t}
=200\,\frac{\pot{2}{-3}}{\pot{5}{-3}}=\SI{80}{\volt}.
\]
\textbf{Note o papel de $N$:} cada espira contribui com a mesma f.e.m., e elas
estão em série. Duzentas espiras multiplicam por duzentos --- é assim que se
obtêm tensões altas a partir de variações de fluxo modestas.}
\end{questao}

\begin{questao}[1]{}
Uma barra de $\SI{30}{\centi\meter}$ move-se a $\SI{4}{\meter\per\second}$
perpendicularmente a um campo de $\SI{0.5}{\tesla}$. Determine a f.e.m.
induzida.
\resposta{$\varepsilon=\SI{0.6}{\volt}$}
\resolucao{
\[
\varepsilon=B\ell v=(0{,}5)(0{,}30)(4)=\SI{0.6}{\volt}.
\]
\textbf{Origem da fórmula:} os portadores da barra movem-se com ela e sofrem
força magnética $qvB$ ao longo do comprimento. Eles se acumulam nas
extremidades até que o campo elétrico resultante equilibre a força magnética ---
e a diferença de potencial daí resultante é $B\ell v$.\\
\textbf{Equivalência com Faraday:} se a barra fizer parte de um circuito, a área
varia à taxa $\ell v$ e $\mathrm{d}\Phi/\mathrm{d}t=B\ell v$ --- o mesmo
resultado.}
\end{questao}

\begin{questao}[1]{}
Determine o número de espiras necessário para obter $\SI{12}{\volt}$ com
variação de fluxo de $\SI{0.04}{\weber\per\second}$.
\resposta{$N=300$}
\resolucao{
\[
N=\frac{\varepsilon}{\mathrm{d}\Phi/\mathrm{d}t}=\frac{12}{0{,}04}=300.
\]
\textbf{Compromisso de projeto:} mais espiras dão mais tensão, mas aumentam a
resistência do enrolamento (mais perdas) e o volume. Aumentar
$\mathrm{d}\Phi/\mathrm{d}t$ --- por campo mais intenso ou variação mais rápida
--- costuma ser preferível quando possível.}
\end{questao}

\begin{questao}[1]{}
Uma bobina de $100$ espiras deve gerar $\SI{50}{\volt}$ com variação de fluxo
de $\SI{5}{\milli\weber}$. Determine o tempo em que essa variação deve ocorrer.
\resposta{$\Delta t=\SI{10}{\milli\second}$}
\resolucao{
\[
\Delta t=\frac{N\,\Delta\Phi}{\varepsilon}
=\frac{100(\pot{5}{-3})}{50}=\pot{1}{-2}\,\si{\second}.
\]
\textbf{Leitura:} a f.e.m. é inversamente proporcional ao tempo. A \emph{mesma}
variação de fluxo, feita dez vezes mais rápido, produz dez vezes mais tensão ---
princípio de funcionamento das bobinas de ignição e dos geradores de alta
tensão por interrupção de corrente.}
\end{questao}

\begin{questao}[1]{}
O sinal negativo na lei de Faraday expressa:
\begin{alternativas}[2]
\item a lei de Lenz e a conservação da energia
\item que a f.e.m. é sempre negativa
\item um erro de convenção
\item que a corrente é sempre horária
\end{alternativas}
\resposta{(a)}
\resolucao{O sinal indica que a corrente induzida circula no sentido de
\textbf{opor-se} à variação de fluxo que a produziu.\\
\textbf{Por que precisa ser assim:} se a corrente induzida \emph{reforçasse} a
variação, ela aumentaria o fluxo, que induziria mais corrente, e assim por
diante --- energia surgiria do nada. A lei de Lenz é a garantia de que isso não
ocorre.\\
\textbf{Consequência prática:} qualquer variação de fluxo encontra
\emph{resistência}. Aproximar um ímã de uma espira fechada exige força; girar
um gerador sob carga exige torque. É de onde vem a energia elétrica.}
\end{questao}

\begin{questao}[1]{}
Uma espira está imersa em campo magnético uniforme e \textbf{constante}, e é
deslocada rapidamente dentro da região. Determine a f.e.m. induzida.
\resposta{$\varepsilon=0$}
\resolucao{O fluxo é $\Phi=BA$, com $B$ e $A$ constantes --- ele \textbf{não
varia}, e portanto
\[
\varepsilon=-\frac{\mathrm{d}\Phi}{\mathrm{d}t}=0.
\]
\textbf{Ponto conceitual central:} \emph{movimento não induz f.e.m.}; variação
de fluxo é que induz. Uma espira pode mover-se a grande velocidade sem gerar
nada, desde que o fluxo através dela permaneça o mesmo.\\
\textbf{Onde há f.e.m.:} apenas ao \emph{atravessar a fronteira} da região de
campo, quando a área imersa muda.}
\end{questao}

\begin{questao}[1]{}
Mostre que $\si{\weber\per\second}$ equivale a volt.
\resposta{$\si{\weber}=\si{\volt\second}$}
\resolucao{Da lei de Faraday, $\varepsilon=\mathrm{d}\Phi/\mathrm{d}t$, logo
\[
[\si{\volt}]=\frac{[\si{\weber}]}{[\si{\second}]}
\;\Longrightarrow\;
\si{\weber}=\si{\volt}\cdot\si{\second}.
\]
\textbf{Leitura útil:} o weber é um "volt-segundo". Isso torna imediato o
resultado de que a \emph{integral} da f.e.m. no tempo é a variação de fluxo ---
base do fluxímetro e do método de medição por integração.}
\end{questao}

\begin{questao}[1]{}
Segundo a lei de Lenz, quando o fluxo através de uma espira \textbf{aumenta},
a corrente induzida cria um campo próprio que:
\begin{alternativas}[2]
\item reforça o campo externo
\item opõe-se ao campo externo
\item é perpendicular ao campo externo
\item é nulo
\end{alternativas}
\resposta{(b)}
\resolucao{A corrente induzida produz campo que se \emph{opõe ao aumento} --- ou
seja, no sentido contrário ao campo externo dentro da espira.\\
\textbf{Se o fluxo diminuísse}, a corrente inverteria, criando campo no
\emph{mesmo} sentido do externo, para tentar sustentá-lo.\\
\textbf{Regra prática:} a espira "não gosta de mudanças". Ela sempre reage no
sentido de manter o fluxo como estava --- nunca no de acelerar a mudança.}
\end{questao}

\blocoaplicacao

\begin{questao}[2]{}
Uma barra de $\SI{50}{\centi\meter}$ desliza a $\SI{3}{\meter\per\second}$
sobre trilhos em campo de $\SI{0.4}{\tesla}$, com resistência total de
$\SI{2}{\ohm}$.
\begin{itensq}
\item Determine a f.e.m. e a corrente.
\item Determine a força necessária para manter a velocidade.
\item Verifique o balanço de potência.
\end{itensq}
\resposta{$\varepsilon=\SI{0.6}{\volt}$; $I=\SI{0.3}{\ampere}$;
$F=\SI{60}{\milli\newton}$; $P=\SI{0.18}{\watt}$ pelos dois caminhos}
\resolucao{(a)
\[
\varepsilon=B\ell v=(0{,}4)(0{,}5)(3)=\SI{0.6}{\volt};
\qquad
I=\frac{0{,}6}{2}=\SI{0.3}{\ampere}.
\]
(b) A corrente induzida na barra, imersa no campo, sofre força magnética
\emph{oposta} ao movimento (lei de Lenz):
\[
F=BI\ell=(0{,}4)(0{,}3)(0{,}5)=\SI{0.06}{\newton}.
\]
Para manter a velocidade constante, o agente externo deve aplicar força igual e
contrária.\\
(c) \textbf{Balanço.}
\[
P_{\text{mecânica}}=Fv=(0{,}06)(3)=\SI{0.18}{\watt};
\qquad
P_{\text{elétrica}}=\frac{\varepsilon^{2}}{R}=\frac{0{,}36}{2}=\SI{0.18}{\watt}.
\]
\textbf{Coincidem exatamente} --- toda a energia mecânica fornecida vira calor no
resistor. É a lei de Lenz em forma quantitativa: a força de oposição é
precisamente a necessária para que a conta feche.}
\end{questao}

\begin{questao}[2]{}
Uma bobina de $100$ espiras e área $\SI{200}{\centi\meter\squared}$ gira a
$\SI{377}{\radian\per\second}$ em campo de $\SI{0.5}{\tesla}$.
\begin{itensq}
\item Determine a f.e.m. máxima.
\item Determine a f.e.m. eficaz.
\end{itensq}
\resposta{$\varepsilon_{\max}=\SI{377}{\volt}$;
$\varepsilon_{ef}=\SI{267}{\volt}$}
\resolucao{(a) Com $\Phi=BA\cos\omega t$:
\[
\varepsilon=-N\frac{\mathrm{d}\Phi}{\mathrm{d}t}=NBA\omega\sin\omega t
\;\Longrightarrow\;
\varepsilon_{\max}=NBA\omega.
\]
\[
\varepsilon_{\max}=100(0{,}5)(\pot{2}{-2})(377)=\SI{377}{\volt}.
\]
(b) Para forma senoidal,
\[
\varepsilon_{ef}=\frac{\varepsilon_{\max}}{\sqrt2}=\SI{267}{\volt}.
\]
\textbf{Note:} $\omega=\SI{377}{\radian\per\second}$ corresponde a
$\SI{60}{\hertz}$ --- é o valor $2\pi(60)$, e aparece constantemente em
eletrotécnica.\\
\textbf{Compromisso do gerador:} $\varepsilon\propto\omega$, mas a frequência é
fixada pela rede. Resta aumentar $N$, $B$ ou $A$ --- e é por isso que geradores
de grande porte têm núcleos volumosos e campos próximos da saturação.}
\end{questao}

\begin{questao}[2]{}
Duas bobinas têm indutância mútua $M=\SI{50}{\milli\henry}$. A corrente na
primeira varia à taxa de $\SI{200}{\ampere\per\second}$. Determine a f.e.m.
induzida na segunda.
\resposta{$\varepsilon_2=\SI{10}{\volt}$}
\resolucao{
\[
\varepsilon_2=M\frac{\mathrm{d}i_1}{\mathrm{d}t}=(0{,}05)(200)=\SI{10}{\volt}.
\]
\textbf{Simetria da indução mútua:} $M_{12}=M_{21}$ sempre, mesmo que as bobinas
sejam muito diferentes. Uma corrente variando na bobina pequena induz na grande
exatamente a mesma f.e.m. que o inverso produziria --- resultado de
reciprocidade, análogo ao que vale em redes resistivas.\\
\textbf{Aplicação:} é este o princípio do transformador, e também da
interferência indesejada entre circuitos vizinhos.}
\end{questao}

\begin{questao}[2]{}
Um transformador ideal tem $1000$ espiras no primário e $100$ no secundário,
com $\SI{220}{\volt}$ na entrada.
\begin{itensq}
\item Determine a tensão de saída.
\item Determine a corrente de entrada para carga que exige
      $\SI{5}{\ampere}$.
\end{itensq}
\resposta{$V_2=\SI{22}{\volt}$; $I_1=\SI{0.5}{\ampere}$}
\resolucao{(a) Como o \emph{mesmo} fluxo atravessa os dois enrolamentos,
\[
\frac{V_2}{V_1}=\frac{N_2}{N_1}
\;\Longrightarrow\;
V_2=220\left(\frac{100}{1000}\right)=\SI{22}{\volt}.
\]
(b) No transformador ideal não há perdas, e a potência se conserva:
\[
V_1I_1=V_2I_2
\;\Longrightarrow\;
I_1=\frac{22(5)}{220}=\SI{0.5}{\ampere}.
\]
\textbf{Relação inversa:} a corrente transforma-se na razão \emph{oposta} à da
tensão. Abaixar tensão dez vezes eleva a corrente dez vezes.\\
\textbf{Consequência para os enrolamentos:} o secundário precisa de fio dez
vezes mais grosso, embora tenha dez vezes menos espiras. O volume de cobre acaba
sendo comparável nos dois lados.}
\end{questao}

\begin{questao}[2]{}
O fluxo em uma espira varia segundo $\Phi(t)=0{,}05\,t^{2}$ (weber, com $t$ em
segundos). Determine a f.e.m. em $t=\SI{2}{\second}$ e a f.e.m. média nos dois
primeiros segundos.
\resposta{$\varepsilon(2)=\SI{0.2}{\volt}$; média de $\SI{0.1}{\volt}$}
\resolucao{\textbf{Instantânea:}
\[
\varepsilon=-\frac{\mathrm{d}\Phi}{\mathrm{d}t}=-0{,}10\,t
\;\Longrightarrow\;
|\varepsilon(2)|=\SI{0.2}{\volt}.
\]
\textbf{Média:}
\[
|\bar\varepsilon|=\frac{\Delta\Phi}{\Delta t}
=\frac{0{,}05(4)-0}{2}=\SI{0.1}{\volt}.
\]
\textbf{As duas diferem por um fator $2$} --- e ambas estão corretas, para
perguntas diferentes.\\
\textbf{Quando coincidem:} apenas se o fluxo variar \emph{linearmente}. Aqui a
variação é quadrática, e a f.e.m. cresce ao longo do intervalo, partindo de zero
e chegando a $\SI{0.2}{\volt}$ --- a média é o valor intermediário.}
\end{questao}

\begin{questao}[2]{}
Uma espira de resistência $\SI{10}{\ohm}$ e $50$ espiras sofre variação de
fluxo de $\SI{4}{\milli\weber}$. Determine a \textbf{carga} total que circula.
\resposta{$q=\SI{20}{\milli\coulomb}$}
\resolucao{Integrando a corrente induzida:
\[
q=\int i\,\mathrm{d}t=\int\frac{|\varepsilon|}{R}\,\mathrm{d}t
=\frac{N}{R}\int\left|\frac{\mathrm{d}\Phi}{\mathrm{d}t}\right|\mathrm{d}t
=\frac{N\,\Delta\Phi}{R}.
\]
\[
q=\frac{50(\pot{4}{-3})}{10}=\pot{2}{-2}\,\si{\coulomb}.
\]
\textbf{Resultado notável: a carga \emph{não depende do tempo}.} Rápida ou
lenta, a mesma variação de fluxo faz circular a mesma carga --- porque
f.e.m. e duração variam inversamente e o produto se conserva.\\
\textbf{É o princípio do fluxímetro balístico:} integrando a corrente com um
galvanômetro de bobina móvel, mede-se $\Delta\Phi$ sem controlar a velocidade do
movimento.}
\end{questao}

\begin{questao}[2]{}
Uma espira quadrada de lado $\SI{10}{\centi\meter}$ entra com velocidade
constante de $\SI{2}{\meter\per\second}$ em uma região de campo de
$\SI{0.6}{\tesla}$.
\begin{itensq}
\item Determine a f.e.m. durante a entrada.
\item Determine a f.e.m. quando totalmente dentro.
\end{itensq}
\resposta{$\SI{0.12}{\volt}$ durante a entrada; zero depois}
\resolucao{(a) Enquanto atravessa a fronteira, apenas o lado dianteiro está no
campo e funciona como barra móvel:
\[
\varepsilon=B\ell v=(0{,}6)(0{,}10)(2)=\SI{0.12}{\volt}.
\]
\emph{(Equivalentemente: a área imersa cresce a $\ell v$, e
$\mathrm{d}\Phi/\mathrm{d}t=B\ell v$.)}\\
(b) \textbf{Totalmente dentro}, os lados dianteiro e traseiro estão ambos no
campo e suas f.e.m. se \emph{cancelam}. O fluxo é constante e
\[
\varepsilon=0.
\]
\textbf{Perfil completo:} pulso retangular de $\SI{0.12}{\volt}$ durante a
entrada, zero no percurso interno, e pulso de \emph{polaridade invertida}
durante a saída. Cada pulso dura $\ell/v=\SI{50}{\milli\second}$.}
\end{questao}

\begin{questao}[2]{}
Um disco metálico gira entre os polos de um ímã. Explique o efeito observado
sobre seu movimento.
\resposta{Ele freia --- correntes de Foucault opõem-se ao movimento}
\resolucao{O disco, ao girar, tem diferentes regiões atravessadas por fluxos
variáveis. Isso induz \textbf{correntes de Foucault} (ou parasitas) --- redemoinhos
de corrente no próprio material.\\
Pela lei de Lenz, essas correntes produzem forças que se \textbf{opõem ao
movimento}: o disco desacelera, e a energia cinética vira calor.\\
\textbf{Duas leituras.}
\begin{itemize}
\item \textbf{Como recurso:} é o \emph{freio magnético}, usado em trens, em
esteiras ergométricas e no disco do antigo medidor de energia. Não há contato,
portanto não há desgaste.
\item \textbf{Como problema:} em núcleos de transformadores e motores, as
mesmas correntes representam perda pura. Combate-se \textbf{laminando} o núcleo
--- assunto retomado no bloco de desafio.
\end{itemize}
\textbf{Propriedade importante:} a força de frenagem é proporcional à
\emph{velocidade}. O freio é forte em alta velocidade e desaparece no repouso ---
por isso ele nunca segura um veículo parado.}
\end{questao}

\begin{questao}[2]{}
Uma bobina com núcleo de ferro tem $\SI{500}{}$ espiras. O fluxo varia de
$\SI{0}{}$ a $\SI{1.2}{\milli\weber}$ em $\SI{8}{\milli\second}$. Determine a
f.e.m. média.
\resposta{$\varepsilon=\SI{75}{\volt}$}
\resolucao{
\[
|\varepsilon|=N\frac{\Delta\Phi}{\Delta t}
=500\,\frac{\pot{1.2}{-3}}{\pot{8}{-3}}=\SI{75}{\volt}.
\]
\textbf{Papel do núcleo:} sem ele, o mesmo enrolamento e a mesma corrente
produziriam fluxo centenas de vezes menor --- e f.e.m. proporcionalmente menor.
O ferro é o que torna a indução tecnologicamente útil.\\
\textbf{Ressalva:} se o fluxo variar a partir de um valor já próximo da
saturação, a relação deixa de ser linear e a f.e.m. real fica abaixo da
prevista.}
\end{questao}

\begin{questao}[2]{}
Uma espira é conectada a um resistor e um ímã é afastado dela. Determine o
sentido da força que a espira exerce sobre o ímã.
\resposta{Atrativa --- ela tenta impedir o afastamento}
\resolucao{Ao afastar o ímã, o fluxo através da espira \textbf{diminui}. Pela lei
de Lenz, a corrente induzida circula no sentido de \emph{sustentar} o fluxo ---
criando campo no mesmo sentido do original.\\
A espira comporta-se então como um ímã com o polo oposto voltado ao ímã que se
afasta, e portanto o \textbf{atrai}.\\
\textbf{Pela terceira lei de Newton}, o ímã atrai a espira com força igual e
contrária.\\
\textbf{Conclusão geral:} a espira sempre \emph{resiste} ao que se está fazendo.
Aproximar exige empurrar contra repulsão; afastar exige puxar contra atração.
Em ambos os casos, o trabalho realizado é o que aparece como energia elétrica
--- e o resistor aquece.}
\end{questao}

\begin{questao}[2]{}
Uma bobina de $200$ espiras, área $\SI{100}{\centi\meter\squared}$ e
resistência $\SI{40}{\ohm}$ está em campo de $\SI{0.8}{\tesla}$
perpendicular. O campo é desligado em $\SI{20}{\milli\second}$.
\begin{itensq}
\item Determine a f.e.m. e a corrente induzidas.
\item Determine o sentido da corrente em relação ao campo original.
\item Determine a carga total.
\end{itensq}
\resposta{$\varepsilon=\SI{80}{\volt}$; $I=\SI{2}{\ampere}$;
$q=\SI{40}{\milli\coulomb}$}
\resolucao{(a)
\[
|\varepsilon|=N\frac{\Delta\Phi}{\Delta t}
=200\,\frac{(0{,}8)(\pot{1}{-2})}{\pot{2}{-2}}=\SI{80}{\volt};
\qquad
I=\frac{80}{40}=\SI{2}{\ampere}.
\]
(b) O fluxo \textbf{diminui}, logo a corrente induzida circula no sentido de
\emph{sustentá-lo} --- criando campo próprio no \textbf{mesmo} sentido do campo
original.\\
\emph{(Se o campo estivesse sendo ligado, a corrente circularia ao contrário.)}\\
(c)
\[
q=\frac{N\,\Delta\Phi}{R}=\frac{200(\pot{8}{-3})}{40}
=\pot{4}{-2}\,\si{\coulomb}.
\]
\textbf{Observação de segurança:} $\SI{80}{\volt}$ surgem em uma bobina que
antes não tinha tensão alguma. Desligar bruscamente eletroímãs de grande porte
produz tensões muito superiores --- razão pela qual eles exigem circuitos de
desmagnetização controlada.}
\end{questao}

\blocoaprofundamento

\begin{questao}[3]{}
Uma barra de $\SI{100}{\gram}$ e $\SI{40}{\centi\meter}$ desliza sem atrito em
trilhos \textbf{verticais}, em campo horizontal de $\SI{0.5}{\tesla}$, com
resistência total de $\SI{0.5}{\ohm}$. Ela é solta do repouso.
\begin{itensq}
\item Escreva a equação do movimento.
\item Determine a velocidade-limite.
\item Determine a potência dissipada nesse regime.
\end{itensq}
\dados{$g=\SI{9.8}{\meter\per\second\squared}$}
\resposta{$v_{lim}=\SI{12.25}{\meter\per\second}$; $P=\SI{12}{\watt}$}
\resolucao{(a) A barra sofre o peso e a força magnética de oposição:
\[
m\frac{\mathrm{d}v}{\mathrm{d}t}=mg-\frac{B^{2}\ell^{2}v}{R},
\]
pois $F=BI\ell$ com $I=B\ell v/R$.\\
\textbf{Estrutura da equação:} é idêntica à de um corpo em queda com arrasto
\emph{linear} --- a força de oposição é proporcional à velocidade.\\
(b) Na velocidade-limite, a aceleração se anula:
\[
v_{lim}=\frac{mgR}{B^{2}\ell^{2}}
=\frac{(0{,}1)(9{,}8)(0{,}5)}{(0{,}25)(0{,}16)}
=\frac{0{,}49}{0{,}04}=\SI{12.25}{\meter\per\second}.
\]
(c)
\[
P=mgv_{lim}=(0{,}1)(9{,}8)(12{,}25)=\SI{12}{\watt},
\]
que é exatamente $\dfrac{(B\ell v_{lim})^{2}}{R}$ \checkmark.\\
\textbf{Interpretação:} em regime, toda a energia potencial gravitacional
perdida vira calor no resistor. A barra não ganha mais energia cinética.\\
\textbf{Efeito de $R$:} reduzir $R$ \emph{diminui} a velocidade-limite --- o
freio fica mais forte. Com $R\to0$ (trilhos supercondutores), a barra mal se
moveria.}
\end{questao}

\begin{questao}[3]{}
Uma bobina de $N$ espiras e resistência $R$ é retirada de um campo. Demonstre
que a carga induzida independe da velocidade do movimento.
\begin{itensq}
\item Deduza a expressão.
\item Explique fisicamente.
\item Indique a aplicação em instrumentação.
\end{itensq}
\resposta{$q=\dfrac{N\Delta\Phi}{R}$ --- independe de $\Delta t$}
\resolucao{(a)
\[
q=\int i\,\mathrm{d}t
=\int\frac{N}{R}\frac{\mathrm{d}\Phi}{\mathrm{d}t}\,\mathrm{d}t
=\frac{N}{R}\int\mathrm{d}\Phi
=\frac{N\,\Delta\Phi}{R}.
\]
O tempo \textbf{cancela} na integração.\\
(b) \textbf{Explicação física.} Retirar a bobina depressa produz f.e.m. alta
durante pouco tempo; devagar, f.e.m. baixa durante muito tempo. O
\emph{produto} --- que é a carga --- permanece o mesmo.\\
Algebricamente: $\varepsilon\propto1/\Delta t$ e a duração é
$\Delta t$; o produto não depende dela.\\
(c) \textbf{Aplicação: fluxímetro balístico.} Um galvanômetro balístico responde
à \emph{carga} total, não à corrente instantânea. Conectado à bobina, sua
deflexão máxima é proporcional a $\Delta\Phi$ --- e o operador não precisa
controlar a velocidade com que retira a bobina.\\
\textbf{Robustez notável:} nem a velocidade, nem a suavidade, nem a trajetória
influenciam o resultado. É uma das medições mais tolerantes a erro de
procedimento em todo o eletromagnetismo.}
\end{questao}

\begin{questao}[3]{}
Um núcleo de transformador maciço de $\SI{10}{\milli\meter}$ de espessura é
substituído por $20$ lâminas isoladas de $\SI{0.5}{\milli\meter}$.
\begin{itensq}
\item Determine a redução das perdas por correntes parasitas.
\item Explique o mecanismo.
\item Indique a limitação prática.
\end{itensq}
\resposta{Redução de $400$ vezes}
\resolucao{(a) As perdas por correntes de Foucault em lâminas obedecem a
\[
P\propto\frac{B_{\max}^{2}f^{2}t^{2}}{\rho},
\]
com $t$ a espessura da lâmina. Reduzindo $t$ de $10$ para
$\SI{0.5}{\milli\meter}$:
\[
\frac{P'}{P}=\left(\frac{0{,}5}{10}\right)^{2}=\frac{1}{400}.
\]
(b) \textbf{Mecanismo.} As correntes parasitas circulam em laços
\emph{perpendiculares} ao fluxo. Laços maiores encerram mais fluxo (mais f.e.m.)
e encontram menos resistência (caminho mais largo) --- e por isso a perda cresce
com o \emph{quadrado} da dimensão.\\
Isolar as lâminas quebra os laços grandes em muitos laços pequenos, cada um com
f.e.m. menor e resistência maior.\\
\textbf{Importante:} as lâminas devem ser paralelas ao fluxo e isoladas
\emph{entre si} --- um verniz de poucos micrômetros basta. Se houver rebarba
metálica unindo-as, o efeito se perde.\\
(c) \textbf{Limitações.}
\begin{itemize}
\item \textbf{Fator de empilhamento:} o isolamento ocupa espaço, e a seção
magnética efetiva cai para $95$ a $97\%$ da geométrica.
\item \textbf{Custo:} lâminas mais finas são mais caras de produzir e montar.
\item \textbf{Frequência:} como $P\propto f^{2}t^{2}$, em alta frequência nem
lâminas finíssimas bastam --- daí o uso de \textbf{ferrites}, que são cerâmicas
de altíssima resistividade e praticamente não conduzem correntes parasitas.
\end{itemize}}
\end{questao}

\begin{questao}[3]{}
Um solenoide longo de $n_1$ espiras por metro e seção $A$ tem, em torno de si,
uma bobina de $N_2$ espiras.
\begin{itensq}
\item Deduza a indutância mútua.
\item Avalie para $n_1=\SI{2000}{\per\meter}$, $A=\SI{5}{\centi\meter\squared}$
      e $N_2=100$.
\item Explique por que $M$ não depende do raio da bobina externa.
\end{itensq}
\resposta{$M=\mu_0n_1N_2A$; $M=\SI{1.26}{\milli\henry}$}
\resolucao{(a) A corrente $i_1$ no solenoide cria $B=\mu_0n_1i_1$ em seu
interior, e fluxo $\Phi=BA$ por espira da bobina externa. O fluxo concatenado
nela é
\[
\lambda_2=N_2\mu_0n_1i_1A
\;\Longrightarrow\;
M=\frac{\lambda_2}{i_1}=\mu_0n_1N_2A.
\]
(b)
\[
M=(4\pi\times10^{-7})(2000)(100)(\pot{5}{-4})
=\pot{1.26}{-4}\,\si{\henry}=\SI{0.126}{\milli\henry}.
\]
(c) \textbf{Por que o raio externo não importa.} O campo do solenoide longo é
praticamente nulo do lado de fora. Toda a contribuição ao fluxo vem da região
\emph{interna} ao solenoide, cuja área é $A$ --- não a área da bobina
externa.\\
Enrolar a bobina justa ou folgada dá o \emph{mesmo} $M$.\\
\textbf{Consequência de projeto:} para aumentar o acoplamento, aumente a seção
do solenoide, não o diâmetro do enrolamento secundário. E note que
$M=M_{21}=M_{12}$ --- calcular pelo caminho fácil (do solenoide para a bobina) dá
o mesmo resultado que o caminho difícil.}
\end{questao}

\begin{questao}[3]{}
Uma espira retangular está próxima a um fio retilíneo cuja corrente varia
segundo $i(t)=I_0\sin\omega t$.
\begin{itensq}
\item Escreva o fluxo através da espira.
\item Determine a f.e.m. induzida.
\item Explique por que essa interferência é difícil de eliminar.
\end{itensq}
\resposta{$\varepsilon=-M\omega I_0\cos\omega t$, com
$M=\dfrac{\mu_0a}{2\pi}\ln\dfrac{d+b}{d}$}
\resolucao{(a) Do resultado já obtido para o fluxo de um fio através de espira
retangular:
\[
\Phi(t)=\frac{\mu_0\,i(t)\,a}{2\pi}\ln\!\left(\frac{d+b}{d}\right)
=M\,i(t),
\]
com $M$ a indutância mútua entre fio e espira.\\
(b)
\[
\varepsilon=-\frac{\mathrm{d}\Phi}{\mathrm{d}t}
=-M\,\frac{\mathrm{d}i}{\mathrm{d}t}
=-M\omega I_0\cos\omega t.
\]
A amplitude é $M\omega I_0$ --- \textbf{proporcional à frequência}.\\
(c) \textbf{Por que é difícil eliminar.}
\begin{itemize}
\item $M$ depende do \textbf{logaritmo} da razão de distâncias. Afastar a espira
reduz a interferência muito devagar --- dobrar $d$ e $b$ juntos não muda
\emph{nada}.
\item A amplitude cresce com $\omega$: transitórios rápidos e harmônicos de alta
ordem acoplam-se muito mais que a fundamental.
\item Blindagem magnética é cara e pouco eficaz em baixa frequência, ao
contrário da blindagem elétrica.
\end{itemize}
\textbf{O que realmente funciona:} reduzir a \emph{área} da espira receptora
(cabos trançados, roteamento junto ao retorno) e reduzir a área da fonte (par
trançado no lado do fio). Atacar a geometria, não a distância.}
\end{questao}

\begin{questao}[3]{}
Uma bobina de $\SI{200}{\milli\henry}$ conduz corrente que varia de
$\SI{2}{\ampere}$ a zero em $\SI{5}{\milli\second}$.
\begin{itensq}
\item Determine a f.e.m. autoinduzida.
\item Determine a energia envolvida.
\item Explique o risco associado.
\end{itensq}
\resposta{$\varepsilon=\SI{80}{\volt}$; $W=\SI{0.4}{\joule}$}
\resolucao{(a)
\[
|\varepsilon|=L\left|\frac{\Delta i}{\Delta t}\right|
=0{,}200\,\frac{2}{\pot{5}{-3}}=\SI{80}{\volt}.
\]
(b) $W=\frac12Li^{2}=\frac12(0{,}2)(4)=\SI{0.4}{\joule}$.\\
(c) \textbf{Risco.} A f.e.m. autoinduzida opõe-se à \emph{redução} da corrente
--- ela tenta mantê-la circulando. Se a interrupção for mais rápida, a tensão
sobe proporcionalmente: interromper em $\SI{5}{\micro\second}$ produziria
$\SI{80}{\kilo\volt}$.\\
Na prática, o ar rompe muito antes e forma-se \textbf{arco elétrico} nos
contatos, que os corrói e gera interferência.\\
\textbf{Proteção:} diodo de roda livre em antiparalelo, que oferece caminho à
corrente e limita a tensão. Os $\SI{0.4}{\joule}$ então se dissipam na
resistência do próprio circuito, sem arco.\\
\textbf{Ponto conceitual:} a autoindução é indução mútua da bobina
\emph{consigo mesma} --- a mesma lei de Faraday, aplicada ao próprio fluxo.}
\end{questao}

\begin{questao}[3]{}
Um gerador fornece $\SI{500}{\watt}$ a uma carga. Analise o torque necessário
para acioná-lo e sua relação com a lei de Lenz.
\begin{itensq}
\item Determine o torque a $\SI{1800}{rpm}$, supondo rendimento de
      $\SI{90}{\percent}$.
\item Explique o que acontece se a carga for desligada.
\end{itensq}
\resposta{$\tau=\SI{2.95}{\newton\meter}$; sem carga, o torque cai quase a
zero}
\resolucao{(a) A potência mecânica de entrada é
\[
P_{mec}=\frac{500}{0{,}90}=\SI{556}{\watt};
\qquad
\omega=\frac{1800(2\pi)}{60}=\SI{188.5}{\radian\per\second};
\]
\[
\tau=\frac{P_{mec}}{\omega}=\frac{556}{188{,}5}=\SI{2.95}{\newton\meter}.
\]
(b) \textbf{Sem carga}, não circula corrente no enrolamento. Sem corrente, não
há força magnética de oposição --- e o torque necessário cai a praticamente
zero (resta apenas o atrito e as perdas em vazio).\\
\textbf{É a lei de Lenz em ação.} O torque resistente \emph{surge} porque a
corrente induzida cria campo que se opõe ao movimento. Mais carga significa mais
corrente, mais oposição, mais torque exigido.\\
\textbf{Experiência reveladora:} girar um gerador manualmente com os terminais
abertos é fácil; curto-circuitá-los torna a rotação quase impossível. A energia
elétrica não vem "do campo" --- vem do braço de quem gira.}
\end{questao}

\begin{questao}[3]{}
Um anel condutor é colocado sobre o núcleo de um solenoide alimentado por
corrente alternada. Observa-se que ele é \textbf{arremessado} para cima.
\begin{itensq}
\item Explique o fenômeno.
\item Determine o que muda se o anel for cortado.
\item Determine o que muda com corrente contínua em regime.
\end{itensq}
\resposta{Repulsão por corrente induzida; anel cortado não salta; em CC de
regime, nada acontece}
\resolucao{(a) \textbf{Explicação.} A corrente alternada no solenoide produz
fluxo variável através do anel, induzindo nele corrente. Pela lei de Lenz, essa
corrente cria campo que se \emph{opõe} ao do solenoide --- e dois campos opostos
se \textbf{repelem}.\\
Como a repulsão persiste em ambos os semiciclos (a corrente induzida inverte
junto com a indutora), a força média é sempre para cima.\\
\emph{(É o experimento clássico do "anel saltante" de Elihu Thomson.)}\\
(b) \textbf{Anel cortado:} não há caminho fechado, logo não circula corrente. Sem
corrente, não há campo induzido nem força. \textbf{O anel não salta} --- embora a
f.e.m. continue sendo induzida em seus terminais.\\
Este é o teste decisivo: mostra que a força depende da \emph{corrente}, não da
f.e.m.\\
(c) \textbf{Corrente contínua em regime:} o fluxo é constante, não há indução,
não há força. \textbf{Nada acontece.}\\
\emph{Mas no instante em que a corrente CC é \textbf{ligada}, há um transitório
de fluxo --- e o anel dá um salto único.} Isso reforça o ponto: é a
\emph{variação} que induz, não o campo.}
\end{questao}

\begin{questao}[3]{}
Uma espira de área $\SI{50}{\centi\meter\squared}$ está em campo que varia
segundo $B(t)=0{,}4\sin(100\pi t)$ tesla, e \emph{simultaneamente} sua área
diminui a $\SI{10}{\centi\meter\squared\per\second}$.
\begin{itensq}
\item Escreva a expressão geral da f.e.m.
\item Identifique as duas contribuições.
\end{itensq}
\resposta{$\varepsilon=-\left(A\dfrac{\mathrm{d}B}{\mathrm{d}t}
+B\dfrac{\mathrm{d}A}{\mathrm{d}t}\right)$}
\resolucao{(a) Como $\Phi=BA$ com \emph{ambos} variáveis, aplica-se a regra do
produto:
\[
\varepsilon=-\frac{\mathrm{d}}{\mathrm{d}t}(BA)
=-\left(A\frac{\mathrm{d}B}{\mathrm{d}t}+B\frac{\mathrm{d}A}{\mathrm{d}t}\right).
\]
(b) \textbf{Duas contribuições de naturezas diferentes.}
\begin{itemize}
\item \textbf{Termo $A\,\mathrm{d}B/\mathrm{d}t$:} f.e.m. induzida por
\emph{campo variável} --- o mecanismo do transformador. Aqui,
$\mathrm{d}B/\mathrm{d}t=40\pi\cos(100\pi t)$, com amplitude
$\SI{125.7}{\tesla\per\second}$, dando até
$\SI{0.63}{\volt}$.
\item \textbf{Termo $B\,\mathrm{d}A/\mathrm{d}t$:} f.e.m. por \emph{movimento}
--- o mecanismo do gerador. Com
$\mathrm{d}A/\mathrm{d}t=\SI{-1e-3}{\meter\squared\per\second}$, contribui com
até $\SI{0.4}{\milli\volt}$.
\end{itemize}
\textbf{Neste caso o primeiro termo domina} por três ordens de grandeza.\\
\textbf{Ponto conceitual:} os dois mecanismos são fisicamente distintos --- um
envolve campo elétrico induzido, o outro força magnética sobre portadores --- mas
a lei de Faraday os unifica em uma única expressão. Essa unificação é uma das
observações que levaram Einstein à relatividade restrita.}
\end{questao}

\begin{questao}[3]{}
Uma bobina de $500$ espiras, área $\SI{20}{\centi\meter\squared}$ e resistência
$\SI{25}{\ohm}$ é girada de $\ang{180}$ em campo de $\SI{0.3}{\tesla}$.
\begin{itensq}
\item Determine a variação de fluxo por espira.
\item Determine a carga total induzida.
\end{itensq}
\resposta{$\Delta\Phi=\SI{1.2}{\milli\weber}$; $q=\SI{24}{\milli\coulomb}$}
\resolucao{(a) Girar $\ang{180}$ inverte o sinal do fluxo:
\[
\Phi_i=BA=(0{,}3)(\pot{2}{-3})=\SI{0.6}{\milli\weber};
\qquad
\Phi_f=-\SI{0.6}{\milli\weber};
\]
\[
|\Delta\Phi|=\SI{1.2}{\milli\weber}\quad\text{--- o \textbf{dobro} de }
\Phi_i.
\]
(b)
\[
q=\frac{N|\Delta\Phi|}{R}=\frac{500(\pot{1.2}{-3})}{25}
=\pot{2.4}{-2}\,\si{\coulomb}.
\]
\textbf{Erro comum:} usar $\Delta\Phi=BA$ em vez de $2BA$. A inversão completa
faz o fluxo passar de $+\Phi$ a $-\Phi$, e a variação é o dobro do valor
inicial.\\
\textbf{Vantagem prática do giro de $\ang{180}$:} ele dobra o sinal em relação a
simplesmente retirar a bobina do campo --- técnica padrão em medições com
fluxímetro.}
\end{questao}

\begin{questao}[3]{}
Um transformador com $k=0{,}95$ e $L_1=\SI{100}{\milli\henry}$,
$L_2=\SI{25}{\milli\henry}$.
\begin{itensq}
\item Determine a indutância mútua.
\item Determine a f.e.m. no secundário para
      $\mathrm{d}i_1/\mathrm{d}t=\SI{50}{\ampere\per\second}$.
\item Determine a fração de fluxo que não se acopla.
\end{itensq}
\resposta{$M=\SI{47.5}{\milli\henry}$; $\varepsilon_2=\SI{2.38}{\volt}$;
$5\%$ de dispersão}
\resolucao{(a)
\[
M=k\sqrt{L_1L_2}=0{,}95\sqrt{(0{,}100)(0{,}025)}
=0{,}95(0{,}05)=\SI{47.5}{\milli\henry}.
\]
(b)
\[
\varepsilon_2=M\frac{\mathrm{d}i_1}{\mathrm{d}t}=(0{,}0475)(50)
=\SI{2.38}{\volt}.
\]
(c) A fração acoplada é $k=95\%$; portanto \textbf{$5\%$ do fluxo se
dispersa}, não atravessando o secundário.\\
\textbf{Efeito da dispersão:} ela aparece como indutância em série
(\emph{indutância de dispersão}) e provoca queda de tensão sob carga. Um
transformador com $k=0{,}95$ tem regulação bem pior que um com
$k=0{,}999$.\\
\textbf{Comparação de topologias:} núcleos toroidais fechados atingem
$k>0{,}99$; núcleos em E com entreferro parasita ficam em torno de $0{,}95$;
transformadores com enrolamentos lado a lado, bem menos.}
\end{questao}

\begin{questao}[3]{}
Analise o perfil de f.e.m. produzido por um fluxo que cresce linearmente por
$\SI{2}{\second}$, permanece constante por $\SI{3}{\second}$ e decresce
linearmente por $\SI{1}{\second}$, com amplitude de $\SI{6}{\milli\weber}$ e
$N=100$.
\begin{itensq}
\item Determine a f.e.m. em cada trecho.
\item Verifique a área total sob a curva de f.e.m.
\end{itensq}
\resposta{$+0{,}3$; $0$; $\SI{-0.6}{\volt}$; área líquida nula}
\resolucao{(a)
\begin{center}
\begin{tabular}{lccc}
\hline
Trecho & $\mathrm{d}\Phi/\mathrm{d}t$ & $\varepsilon$ & Duração\\
\hline
Subida & $\SI{3}{\milli\weber\per\second}$ & $\SI{-0.3}{\volt}$
 & $\SI{2}{\second}$\\
Platô & $0$ & $0$ & $\SI{3}{\second}$\\
Descida & $\SI{-6}{\milli\weber\per\second}$ & $\SI{+0.6}{\volt}$
 & $\SI{1}{\second}$\\
\hline
\end{tabular}
\end{center}
(Os sinais decorrem de $\varepsilon=-N\,\mathrm{d}\Phi/\mathrm{d}t$.)\\
\textbf{Observe:} a descida é duas vezes mais rápida, e a f.e.m. é o
\emph{dobro} em módulo.\\
(b) \textbf{Área sob a curva:}
\[
(-0{,}3)(2)+(0)(3)+(0{,}6)(1)=-0{,}6+0{,}6=0.
\]
\textbf{Isso não é coincidência.} A área é $\int\varepsilon\,\mathrm{d}t
=-N\Delta\Phi_{total}$, e o fluxo terminou onde começou. \textbf{Sempre que o
fluxo retorna ao valor inicial, a integral da f.e.m. é nula.}\\
\textbf{Consequência prática:} em regime periódico, a tensão média sobre um
indutor é sempre zero --- princípio do \emph{balanço volt-segundo}, base do
projeto de conversores chaveados.}
\end{questao}

\begin{questao}[3]{}
Um freio magnético produz força proporcional à velocidade, $F=bv$, com
$b=\SI{20}{\newton\second\per\meter}$, atuando sobre um corpo de
$\SI{5}{\kilogram}$ inicialmente a $\SI{10}{\meter\per\second}$.
\begin{itensq}
\item Escreva a equação do movimento e resolva.
\item Determine a constante de tempo e a distância percorrida.
\item Explique por que o freio não para o corpo.
\end{itensq}
\resposta{$v=v_0e^{-t/\tau}$ com $\tau=\SI{0.25}{\second}$;
$d=\SI{2.5}{\meter}$}
\resolucao{(a)
\[
m\frac{\mathrm{d}v}{\mathrm{d}t}=-bv
\;\Longrightarrow\;
v(t)=v_0e^{-bt/m}.
\]
(b) $\tau=m/b=5/20=\SI{0.25}{\second}$.\\
A distância total é
\[
d=\int_0^{\infty}v_0e^{-t/\tau}\mathrm{d}t=v_0\tau=10(0{,}25)=\SI{2.5}{\meter}.
\]
\textbf{Mesma estrutura} da descarga exponencial de um capacitor --- e a mesma
propriedade: a "área" (distância) equivale à velocidade inicial mantida por
$\tau$.\\
(c) \textbf{Por que não para.} A força é proporcional a $v$: quando a velocidade
tende a zero, a força também tende a zero. Matematicamente, o corpo leva tempo
\emph{infinito} para parar --- ele apenas se aproxima assintoticamente do
repouso.\\
\textbf{Consequência prática decisiva:} freios magnéticos são excelentes para
\emph{desacelerar} e péssimos para \emph{segurar}. Trens de alta velocidade os
usam na faixa alta e completam a parada com freios de atrito --- que funcionam
justamente onde o magnético falha.}
\end{questao}

\blocodesafio

\begin{questao}[4]{}
\textbf{(Lenz e conservação da energia.)} Demonstre que a lei de Lenz é uma
consequência necessária da conservação da energia.
\begin{itensq}
\item Suponha o contrário e derive a contradição.
\item Quantifique com o exemplo da barra em trilhos.
\item Discuta o que a lei de Lenz \emph{não} determina.
\end{itensq}
\resposta{Sinal oposto levaria a crescimento ilimitado de energia}
\resolucao{(a) \textbf{Demonstração por absurdo.} Suponha que a corrente
induzida \emph{reforçasse} a variação de fluxo. Ao aproximar um ímã de uma
espira, ela o \textbf{atrairia} --- acelerando a aproximação, o que aumentaria
$\mathrm{d}\Phi/\mathrm{d}t$, que aumentaria a corrente, que aumentaria a
atração.\\
O sistema entraria em fuga: energia cinética \emph{e} energia elétrica cresceriam
simultaneamente, sem fonte. \textbf{Contradição} com a conservação da energia.\\
Portanto a corrente deve \emph{opor-se} --- e o sinal negativo não é convenção,
é necessidade.\\
(b) \textbf{Quantificação.} Na barra em trilhos:
\[
P_{mec}=Fv=\left(\frac{B^{2}\ell^{2}v}{R}\right)v
=\frac{B^{2}\ell^{2}v^{2}}{R};
\qquad
P_{el}=\frac{(B\ell v)^{2}}{R}=\frac{B^{2}\ell^{2}v^{2}}{R}.
\]
\textbf{Idênticas, termo a termo.} A força de oposição tem exatamente o módulo
necessário para que a potência mecânica iguale a elétrica --- nem mais, nem
menos.\\
\textbf{Se a força fosse maior}, sobraria energia mecânica sem destino; se fosse
menor, apareceria energia elétrica do nada.\\
(c) \textbf{O que a lei de Lenz não determina.} Ela fixa apenas o \emph{sentido}
da corrente. O \textbf{módulo} vem da lei de Faraday combinada com a
resistência do circuito.\\
E ela não diz nada sobre \emph{quanta} energia é convertida --- isso depende da
carga. Uma espira aberta induz f.e.m. mas não corrente, e não há oposição
alguma: aproximar um ímã de uma espira cortada não exige força.}
\end{questao}

\begin{questao}[4]{}
\textbf{(Projeto de bobina sensora.)} Projete uma bobina que gere
$\SI{5}{\volt}$ médios quando o fluxo variar $\SI{0.2}{\milli\weber}$ em
$\SI{4}{\milli\second}$.
\begin{itensq}
\item Determine o número de espiras.
\item Determine a resistência admissível para corrente de saída inferior a
      $\SI{1}{\milli\ampere}$ em carga de $\SI{10}{\kilo\ohm}$.
\item Avalie a resposta em frequência.
\end{itensq}
\resposta{$N=100$ espiras}
\resolucao{(a)
\[
N=\frac{\varepsilon\,\Delta t}{\Delta\Phi}
=\frac{5(\pot{4}{-3})}{\pot{2}{-4}}=100\ \text{espiras}.
\]
(b) Com carga de $\SI{10}{\kilo\ohm}$ e $\SI{5}{\volt}$, a corrente é
$\SI{0.5}{\milli\ampere}$ --- já abaixo do limite. A queda na resistência
interna deve ser pequena:
\[
\frac{R_{bobina}}{R_{bobina}+10\,\mathrm{k}}\le1\%
\;\Longrightarrow\;
R_{bobina}\le\SI{100}{\ohm}.
\]
Com $100$ espiras, isso é facilmente atingível --- fio de $\SI{0.2}{\milli\meter}$
e alguns metros de comprimento dão poucos ohms.\\
(c) \textbf{Resposta em frequência.} A bobina é um derivador: sua saída é
proporcional a $\mathrm{d}\Phi/\mathrm{d}t$, logo a sensibilidade
\textbf{cresce com a frequência}.
\begin{itemize}
\item \textbf{Vantagem:} excelente para detectar transitórios rápidos e
descargas.
\item \textbf{Limitação inferior:} ela \emph{não} detecta campo estático ---
para isso é preciso sensor Hall ou fluxgate.
\item \textbf{Limitação superior:} a indutância própria da bobina, associada à
capacitância entre espiras, forma um ressonador. Acima da frequência de
autorressonância a resposta deixa de ser proporcional a
$\mathrm{d}\Phi/\mathrm{d}t$.
\end{itemize}
\textbf{Para recuperar $\Phi$ (e não sua derivada):} acrescente um integrador na
saída --- é assim que funciona a bobina de Rogowski para medição de corrente.}
\end{questao}

\begin{questao}[4]{}
\textbf{(Freio por correntes parasitas.)} Modele quantitativamente um freio
magnético de disco.
\begin{itensq}
\item Estabeleça a dependência da força com os parâmetros.
\item Determine por que a força é proporcional à velocidade.
\item Analise o comportamento em velocidades muito altas.
\end{itensq}
\resposta{$F\propto\sigma B^{2}tv$; proporcionalidade a $v$ decorre de
$\varepsilon\propto v$ e $i\propto\varepsilon$}
\resolucao{(a) \textbf{Dependência.} A f.e.m. induzida localmente é
proporcional a $Bv$; a corrente parasita, a $\sigma Bv$ (com $\sigma$ a
condutividade); e a força, ao produto da corrente pelo campo:
\[
F\propto\sigma\,B^{2}\,t\,v\,A,
\]
com $t$ a espessura e $A$ a área sob o polo.\\
\textbf{Note a dependência quadrática em $B$} --- dobrar o campo quadruplica a
frenagem, o que torna ímãs de neodímio muito mais eficazes que ferrites.\\
(b) \textbf{Por que $F\propto v$.} Dois fatores lineares se combinam:
\[
\varepsilon\propto v
\;\Longrightarrow\;
i\propto v
\;\Longrightarrow\;
F=Bi\ell\propto v.
\]
É o mesmo mecanismo do arrasto viscoso --- e leva à mesma equação diferencial,
com decaimento exponencial da velocidade.\\
(c) \textbf{Velocidades muito altas.} A proporcionalidade \emph{deixa de valer}
por dois motivos:
\begin{itemize}
\item \textbf{Reação de armadura:} as correntes parasitas tornam-se intensas o
bastante para criar campo próprio comparável ao aplicado, distorcendo-o. A força
satura e chega a \emph{diminuir}.
\item \textbf{Efeito pelicular:} em alta frequência efetiva, as correntes se
concentram na superfície do disco, reduzindo a espessura útil $t$.
\end{itemize}
\textbf{Existe, portanto, uma velocidade de força máxima}, acima da qual é o freio
perde eficácia --- comportamento oposto ao do freio de atrito e que precisa ser
levado em conta no dimensionamento.}
\end{questao}

\begin{questao}[4]{}
\textbf{(Convenção dos pontos na indução mútua.)} Explique o significado físico
da convenção dos pontos e sua determinação experimental.
\begin{itensq}
\item Enuncie a convenção.
\item Determine o sinal de $M$ nas equações de malha.
\item Descreva o procedimento experimental.
\end{itensq}
\resposta{Pontos marcam terminais magneticamente correspondentes; o sinal de
$M$ depende de as correntes entrarem ou não pelos pontos}
\resolucao{(a) \textbf{Convenção.} Correntes que entram pelos terminais
\emph{pontuados} produzem fluxos que se \textbf{reforçam} no núcleo comum.\\
Isso é informação \emph{física} (sentido de enrolamento), que o desenho
esquemático não transmite --- daí a necessidade da marcação.\\
(b) \textbf{Sinal nas equações.} Para duas bobinas acopladas:
\[
v_1=L_1\frac{\mathrm{d}i_1}{\mathrm{d}t}
\pm M\frac{\mathrm{d}i_2}{\mathrm{d}t};
\qquad
v_2=L_2\frac{\mathrm{d}i_2}{\mathrm{d}t}
\pm M\frac{\mathrm{d}i_1}{\mathrm{d}t}.
\]
Usa-se $+M$ se \emph{ambas} as correntes entram pelos pontos (ou ambas saem), e
$-M$ caso contrário.\\
\textbf{Em série}, isso dá $L_{eq}=L_1+L_2\pm2M$ --- resultado já obtido no banco
de indutores.\\
(c) \textbf{Procedimento experimental.} Meça $L_{eq}$ com as bobinas ligadas em
série nas \emph{duas} formas possíveis:
\[
M=\frac{L_{aditivo}-L_{subtrativo}}{4}.
\]
A ligação de \emph{maior} indutância identifica o acoplamento aditivo --- e
portanto os terminais correspondentes.\\
\textbf{Método alternativo, mais direto:} aplique um pulso de tensão a uma
bobina e observe a polaridade do pulso induzido na outra com osciloscópio. Se
ambas as pontas positivas estiverem nos terminais correspondentes, os sinais têm
a mesma polaridade.\\
\textbf{Por que isso importa na prática:} inverter a marcação de um
transformador em um conversor inverte a fase da saída --- e pode transformar
realimentação negativa em positiva, destruindo o circuito.}
\end{questao}

\begin{questao}[4]{}
\textbf{(Verificação experimental da lei de Faraday.)} Projete um experimento
para verificar $\varepsilon\propto N$ e $\varepsilon\propto v$.
\begin{itensq}
\item Descreva o arranjo e o procedimento.
\item Estabeleça o critério de validação.
\item Aponte as principais fontes de erro.
\end{itensq}
\resposta{Ímã caindo por tubo com bobinas de $N$ variável; verificar
linearidade em $N$ e em $v$}
\resolucao{(a) \textbf{Arranjo.} Um ímã permanente cai por dentro de um tubo
isolante vertical, atravessando bobinas montadas ao longo dele. Um osciloscópio
registra a f.e.m. de cada bobina.
\begin{enumerate}
\item \textbf{Variar $N$:} bobinas de $50$, $100$, $200$ e $400$ espiras, todas
com a mesma área e posição relativa. Compare os picos de f.e.m.
\item \textbf{Variar $v$:} use \emph{a mesma} bobina em alturas diferentes. Como
$v=\sqrt{2gh}$, a velocidade é conhecida a partir da altura de queda.
\end{enumerate}
(b) \textbf{Critério de validação.}
\begin{itemize}
\item Gráfico de $\varepsilon_{pico}$ contra $N$: deve ser uma \textbf{reta pela
origem}. Ajuste linear com $R^{2}>0{,}99$ e intercepto compatível com zero.
\item Gráfico de $\varepsilon_{pico}$ contra $\sqrt{h}$: também reta pela
origem.
\item \textbf{Teste mais forte:} a \emph{área} sob cada pulso de f.e.m. deve ser
$N\Delta\Phi$ --- \emph{independente} da velocidade. Verificar isso confirma
simultaneamente a lei e a independência da carga induzida.
\end{itemize}
(c) \textbf{Fontes de erro.}
\begin{itemize}
\item \textbf{Frenagem do próprio ímã} pelas correntes induzidas: se as bobinas
estiverem em curto ou a carga for baixa, o ímã desacelera e $v$ deixa de ser
$\sqrt{2gh}$. \emph{Solução:} medir em circuito aberto (osciloscópio de alta
impedância).
\item \textbf{Tubo condutor:} um tubo de alumínio ou cobre freia
drasticamente o ímã. Use acrílico ou vidro.
\item \textbf{Posição de disparo do osciloscópio} e largura de banda
insuficiente, que arredondam o pico.
\item \textbf{Orientação do ímã} variando durante a queda --- use guia.
\end{itemize}}
\end{questao}

\begin{questao}[4]{}
\textbf{(Transformador em corrente contínua.)} Explique por que um
transformador não transfere potência em regime de corrente contínua.
\begin{itensq}
\item Analise o regime permanente.
\item Analise o transitório de ligação.
\item Determine o que ocorre com o enrolamento primário.
\end{itensq}
\resposta{Fluxo constante $\Rightarrow$ $\mathrm{d}\Phi/\mathrm{d}t=0
\Rightarrow$ nenhuma f.e.m. no secundário}
\resolucao{(a) \textbf{Regime permanente.} Com corrente contínua constante, o
fluxo no núcleo é constante:
\[
\frac{\mathrm{d}\Phi}{\mathrm{d}t}=0
\;\Longrightarrow\;
\varepsilon_2=0.
\]
\textbf{Nenhuma tensão no secundário}, por maior que seja a corrente no
primário. A indução exige \emph{variação}.\\
(b) \textbf{Transitório de ligação.} No instante em que a fonte CC é ligada, o
fluxo sobe de zero ao valor final --- e \emph{há} f.e.m. no secundário durante
esse intervalo. Aparece um pulso, cuja integral é $N_2\Delta\Phi$.\\
O mesmo ocorre, com polaridade invertida, ao desligar.\\
\textbf{É esse pulso que se explora na bobina de ignição} de motores: uma
corrente contínua é interrompida abruptamente, e o pulso resultante atinge
dezenas de quilovolts.\\
(c) \textbf{O que ocorre com o primário --- e é o ponto crítico.} Em regime CC, a
única coisa que limita a corrente é a \textbf{resistência ôhmica} do
enrolamento, que é propositalmente pequena (poucos ohms).\\
Um primário projetado para $\SI{220}{\volt}$ em $\SI{60}{\hertz}$ tem sua
corrente limitada pela \emph{reatância} $\omega L$, tipicamente centenas de
ohms. Em CC, $\omega=0$ e essa limitação \textbf{desaparece}:
\[
I_{CC}=\frac{220}{R_{enrolamento}}\gg I_{CA}.
\]
\textbf{Resultado:} corrente dezenas de vezes maior que a nominal, saturação
imediata do núcleo e queima do enrolamento em segundos.\\
\textbf{Conclusão:} não é apenas que o transformador "não funciona" em CC --- ele
se \emph{destrói}.}
\end{questao}

\begin{questao}[4]{}
\textbf{(Espira entrando em campo.)} Deduza a f.e.m. durante a entrada de uma
espira retangular em uma região de campo uniforme, com velocidade constante.
\begin{itensq}
\item Deduza pela variação de área e pela força sobre portadores.
\item Determine a força necessária para manter a velocidade.
\item Analise o balanço de energia.
\end{itensq}
\resposta{$\varepsilon=B\ell v$ pelos dois caminhos;
$F=\dfrac{B^{2}\ell^{2}v}{R}$}
\resolucao{(a) \textbf{Pela variação de área.} Com o lado dianteiro a uma
profundidade $x$ dentro do campo, $A_{imersa}=\ell x$ e
\[
\Phi=B\ell x
\;\Longrightarrow\;
\varepsilon=-\frac{\mathrm{d}\Phi}{\mathrm{d}t}=-B\ell\frac{\mathrm{d}x}{\mathrm{d}t}
=-B\ell v.
\]
\textbf{Pela força sobre portadores.} Os portadores do lado dianteiro movem-se
com velocidade $v$ e sofrem força $qvB$ ao longo do lado. O trabalho por unidade
de carga ao percorrer o comprimento $\ell$ é
\[
\varepsilon=\frac{W}{q}=vB\ell.
\]
\textbf{Os dois caminhos coincidem} --- e essa coincidência não é trivial: um
usa a lei de Faraday (variação de fluxo), o outro a força de Lorentz. Sua
equivalência é uma das simetrias mais notáveis do eletromagnetismo.\\
(b) \textbf{Força necessária.} A corrente $i=B\ell v/R$ percorre o lado
dianteiro, imerso no campo, sofrendo
\[
F=Bi\ell=\frac{B^{2}\ell^{2}v}{R},
\]
oposta ao movimento. É essa força que o agente externo deve equilibrar.\\
\textbf{Note:} os lados laterais também estão parcialmente no campo, mas suas
forças são opostas entre si e se cancelam.\\
(c) \textbf{Balanço.}
\[
P_{mec}=Fv=\frac{B^{2}\ell^{2}v^{2}}{R}
=\frac{(B\ell v)^{2}}{R}=\frac{\varepsilon^{2}}{R}=P_{el}.
\]
\textbf{Fecha exatamente.} E note o que acontece quando a espira fica
\emph{totalmente dentro}: os dois lados no campo geram f.e.m. opostas, a
corrente cessa, a força desaparece --- e a espira desliza livremente. A oposição
existe apenas enquanto o fluxo varia.}
\end{questao}

\begin{questao}[4]{}
\textbf{(Balanço volt-segundo.)} Demonstre que, em regime periódico, a tensão
média sobre um indutor é nula, e explore as consequências.
\begin{itensq}
\item Demonstre.
\item Aplique a um conversor com dois patamares de tensão.
\item Explique o que ocorre se o balanço não for respeitado.
\end{itensq}
\resposta{$\bar v_L=0$; daí $V_1t_1=V_2t_2$ em regime}
\resolucao{(a) \textbf{Demonstração.} Em regime periódico de período $T$, a
corrente retorna ao mesmo valor a cada ciclo: $i(T)=i(0)$. Como
$v_L=L\,\mathrm{d}i/\mathrm{d}t$,
\[
\int_0^{T}v_L\,\mathrm{d}t=L\int_{i(0)}^{i(T)}\mathrm{d}i=L\left[i(T)-i(0)\right]=0.
\]
Logo $\bar v_L=0$. Equivalentemente, o \textbf{produto volt-segundo} aplicado
em um sentido deve igualar o aplicado no outro.\\
(b) \textbf{Aplicação a conversor.} Se o indutor recebe $+V_1$ durante
$t_1$ e $-V_2$ durante $t_2$:
\[
V_1t_1=V_2t_2.
\]
\textbf{Exemplo (conversor abaixador):} com entrada $\SI{12}{\volt}$, saída
$\SI{5}{\volt}$ e razão cíclica $D$, o indutor vê $(12-5)$ durante $DT$ e
$-5$ durante $(1-D)T$:
\[
7D=5(1-D)
\;\Longrightarrow\;
D=\frac{5}{12}=0{,}417.
\]
\textbf{O balanço volt-segundo produziu a relação de conversão} $V_o=DV_i$ sem
nenhuma análise de circuito --- é o método padrão de projeto.\\
(c) \textbf{Se o balanço não for respeitado.} A corrente do indutor não retorna
ao valor inicial: ela cresce (ou decresce) a cada ciclo, acumulando
indefinidamente.\\
Em um transformador, isso significa \textbf{saturação progressiva} do núcleo até
o colapso da indutância e a destruição do interruptor --- o fenômeno conhecido
como \emph{staircase saturation}.\\
\textbf{É por isso que conversores em ponte exigem simetria rigorosa} entre os
semiciclos, e por que se acrescenta capacitor de bloqueio em série com o
primário: ele impede qualquer componente contínua de tensão, garantindo o
balanço por construção.}
\end{questao}

\begin{questao}[4]{}
\textbf{(Correntes parasitas em condutores.)} Analise por que condutores
maciços de grande seção apresentam perdas superiores às previstas em corrente
alternada.
\begin{itensq}
\item Explique o efeito pelicular.
\item Determine a profundidade de penetração e avalie para cobre a
      $\SI{60}{\hertz}$.
\item Indique as soluções empregadas.
\end{itensq}
\resposta{$\delta=\sqrt{\dfrac{2\rho}{\omega\mu}}\approx\SI{8.5}{\milli\meter}$
para cobre a $\SI{60}{\hertz}$}
\resolucao{(a) \textbf{Efeito pelicular.} A corrente alternada em um condutor
cria fluxo variável \emph{dentro dele}. Esse fluxo induz correntes parasitas
que, pela lei de Lenz, se opõem à corrente no centro e a reforçam na periferia.\\
Resultado: a corrente concentra-se junto à superfície, a seção efetiva diminui e
a \textbf{resistência aumenta}.\\
(b) \textbf{Profundidade de penetração.}
\[
\delta=\sqrt{\frac{2\rho}{\omega\mu}}.
\]
Para cobre ($\rho=\pot{1.7}{-8}$, $\mu\approx\mu_0$) a
$\SI{60}{\hertz}$:
\[
\delta=\sqrt{\frac{2(\pot{1.7}{-8})}{2\pi(60)(4\pi\times10^{-7})}}
\approx\SI{8.5}{\milli\meter}.
\]
\textbf{Interpretação:} condutores com raio muito menor que
$\SI{8.5}{\milli\meter}$ praticamente não sofrem o efeito em
$\SI{60}{\hertz}$. Já um barramento de $\SI{50}{\milli\meter}$ tem seu centro
essencialmente inútil.\\
\textbf{Dependência com a frequência:} $\delta\propto1/\sqrt{f}$. A
$\SI{1}{\mega\hertz}$, $\delta$ cai para cerca de
$\SI{66}{\micro\meter}$.\\
(c) \textbf{Soluções.}
\begin{itemize}
\item \textbf{Condutores tubulares} ou perfis chatos, que aproveitam melhor a
periferia --- comuns em barramentos de subestação.
\item \textbf{Fio Litz:} muitos fios finos isolados e \emph{transpostos}, de modo
que cada um ocupe alternadamente centro e periferia. Usado em enrolamentos de
alta frequência.
\item \textbf{Vários condutores em paralelo}, transpostos ao longo do percurso.
\end{itemize}
\textbf{Observação:} apenas colocar fios em paralelo \emph{sem} transposição não
resolve --- os externos conduziriam mais que os internos, pelo mesmo mecanismo.
É a transposição que equaliza.}
\end{questao}



\begin{questao}[4]{}
 Um feixe de prótons atravessa inicialmente uma região onde atuam simultaneamente um campo elétrico uniforme de intensidade
\[
E = 70\,\mathrm{kV/m}
\]
e um campo magnético uniforme de intensidade
\[
B = 50\,\mathrm{mT},
\]
perpendiculares entre si e também à direção de movimento das partículas. Os campos são ajustados de modo que apenas os prótons que atravessam essa região sem sofrer desvio alcancem uma segunda região, na qual existe apenas o campo magnético de mesma intensidade, conforme ilustrado na figura.

Desprezando os efeitos gravitacionais e considerando
\[
m_p = 1{,}67\times10^{-27}\,\mathrm{kg},
\qquad
q_p = 1{,}60\times10^{-19}\,\mathrm{C},
\]
determine:

\begin{itensq}
    \item a velocidade dos prótons que atravessam o seletor sem sofrer desvio;
    \item o raio \(r\) da trajetória circular descrita pelos prótons na região onde atua apenas o campo magnético;
    \item a distância \(2r\) entre o ponto de entrada na região magnética e o ponto em que os prótons atingem a placa detectora.
\end{itensq}



 \begin{tikzpicture}[
    >=Latex,
    line cap=round,
    line join=round,
    font=\sffamily
]

% ==================================================
% PARÂMETROS
% ==================================================
\def\r{2.35}
\def\xleft{0.0}
\def\xfilter{3.4}
\def\xslit{6.45}
\pgfmathsetmacro{\xfield}{\xslit-0.35}
\def\ybeam{0.0}
\def\ybottom{-4.7}

% Centro do arco semicircular
\coordinate (C) at (\xfield,\ybottom/2);
\coordinate (TopArc) at ($(C)+(90:\r)$);
\coordinate (BotArc) at ($(C)+(-90:\r)$);

% ==================================================
% PRÓTON INICIAL
% ==================================================
\draw[thick] (-1.05,0) circle (0.23);
\node at (-1.05,0) {\large $+$};

% Feixe horizontal
\draw[thick,dashed] (-0.75,0) -- (\xfield,0);

% ==================================================
% FENDA DE ENTRADA
% ==================================================
\draw[line width=3.6pt] (0.12,0.32) -- (0.12,0.82);
\draw[line width=3.6pt] (0.12,-0.32) -- (0.12,-0.82);

% ==================================================
% SELETOR DE VELOCIDADES
% ==================================================
% Placas
\draw[line width=2.2pt] (2.15,0.78) -- (4.65,0.78);
\draw[line width=2.2pt] (2.15,-0.78) -- (4.65,-0.78);

% Terminais das placas
\draw[line width=1.8pt] (3.40,0.78) -- (3.40,1.62);
\draw[line width=1.8pt] (3.40,-0.78) -- (3.40,-1.62);

% Polaridade
\node[font=\Large] at (2.75,1.15) {$-$};
\node[font=\Large] at (2.75,-1.15) {$+$};

% Campo elétrico (setas para cima)
\foreach \x in {2.45,2.95,3.45,3.95,4.45}{
    \draw[myE,line width=1.5pt,-{Latex[length=3mm,width=2mm]}]
        (\x,-0.62) -- (\x,0.62);
}

% Pontos de B no seletor
\foreach \x in {2.55,3.25,3.95,4.35}{
    \foreach \y in {-0.42,0,0.42}{
        \fill[myB] (\x,\y) circle (2.3pt);
    }
}

% Rótulos vetoriais
\node[myE,above] at (4.95,0.25) {$\vec E$};
\node[myB,right] at (4.55,0.0) {$\vec B$};

% Velocidade
\draw[thick,-{Latex[length=3mm,width=2mm]}] (1.05,0.28) -- (1.95,0.28);
\node[above] at (1.5,0.28) {$\vec v$};

% ==================================================
% FENDA DE SAÍDA
% ==================================================
\draw[line width=3.6pt] (\xslit-0.35,0.32) -- (\xslit-0.35,0.82);
\draw[line width=3.6pt] (\xslit-0.35,-0.32) -- (\xslit-0.35,-0.82);

% ==================================================
% REGIÃO COM CAMPO MAGNÉTICO
% ==================================================
\foreach \x in {6.7,7.4,8.1,8.8,9.5}{
    \foreach \y in {0.45,-0.15,-0.85,-1.55,-2.25,-2.95,-3.65,-4.35, -5}{
        \fill[myB] (\x,\y) circle (2.5pt);
    }
}
\foreach \x/\y in {
    7.05/-0.50,7.75/-1.20,8.45/-0.50,9.15/-1.20,
    7.05/-1.90,7.75/-2.60,8.45/-1.90,9.15/-2.60,
    7.05/-3.30,7.75/-4.00,8.45/-3.30,9.15/-4.00
}{
    \fill[myB] (\x,\y) circle (2.5pt);
}

\node[myB] at (9.75,-0.25) {$\vec B$};

% ==================================================
% TRAJETÓRIA SEMICIRCULAR
% ==================================================
\draw[thick,dashed,-{Latex[length=3mm,width=2mm]}]
    (TopArc) arc[start angle=90,end angle=-90,radius=\r];

% ==================================================
% MEDIDA 2r
% ==================================================
\draw[thick,{Latex[length=2.8mm]}-{Latex[length=2.8mm]}]
    (5.65,-0.02) -- (5.65,\ybottom+0.02);

\draw[thin] (5.42,0) -- (5.90,0);
\draw[thin] (5.42,\ybottom) -- (5.90,\ybottom);

\node[fill=white,inner sep=1.5pt,font=\Large] at (5.65,-2.35) {$2r$};

% ==================================================
% PLACA DETECTORA
% ==================================================
\draw[line width=4pt] (\xslit-0.35,\ybottom+0.55) -- (\xslit-0.35,\ybottom-0.55);
\draw[line width=1.7pt] (\xslit-0.35,\ybottom) -- (\xslit+0.15,\ybottom);

\node[align=center,font=\Large] at (3.95,-4.55) {Detector\\plate};

% ==================================================
% TEXTOS DOS CAMPOS
% ==================================================
\node[anchor=west,text=myE,font=\bfseries\Large] at (-0.55,-2.05)
    {$E = 70\ \mathrm{kV/m}$};

\node[anchor=west,text=myB,font=\bfseries\Large] at (-0.55,-2.85)
    {$B = 50\ \mathrm{mT}$};

\end{tikzpicture}

\resposta{$v=\dfrac{E}{B}=\SI{1.4e6}{\meter\per\second}$,
\quad
$r=\dfrac{m_pv}{q_pB}\approx\SI{0.292}{\meter}$,
\quad
$2r\approx\SI{0.585}{\meter}$}

\resolucao{%
(a) \textbf{Velocidade dos prótons no seletor.}
Para que o próton atravesse a região de campos cruzados sem sofrer desvio,
a força elétrica e a força magnética devem possuir o mesmo módulo e sentidos
opostos:
\[
F_E=F_B.
\]
Como
\[
F_E=qE
\qquad\text{e}\qquad
F_B=qvB,
\]
tem-se
\[
qE=qvB.
\]
Cancelando a carga $q$:
\[
v=\frac{E}{B}.
\]
Substituindo
\[
E=\SI{70}{\kilo\volt\per\meter}
\qquad\text{e}\qquad
B=\SI{50}{\milli\tesla}=\SI{0.050}{\tesla},
\]
resulta
\[
v=
\frac{70\times10^{3}}{0.050}
=
\boxed{\SI{1.4e6}{\meter\per\second}}.
\]

\textbf{Interpretação:} somente os prótons com essa velocidade atravessam o
seletor em linha reta. Prótons mais rápidos ou mais lentos sofrem uma força
resultante e são desviados.\\

(b) \textbf{Raio da trajetória na região magnética.}
Após deixar o seletor, o próton entra em uma região onde atua apenas o campo
magnético. Como a força magnética é sempre perpendicular à velocidade, ela não
altera o módulo da velocidade, mas muda continuamente sua direção, atuando como
força centrípeta:
\[
F_B=F_c.
\]
Portanto,
\[
q_pvB=\frac{m_pv^2}{r}.
\]
Isolando o raio:
\[
r=\frac{m_pv}{q_pB}.
\]
Para
\[
m_p=\SI{1.67e-27}{\kilogram},
\qquad
q_p=\SI{1.60e-19}{\coulomb},
\]
tem-se
\[
r=
\frac{
(\pot{1.67}{-27})(\pot{1.4}{6})
}{
(\pot{1.60}{-19})(0.050)
}
\approx
\boxed{\SI{0.292}{\meter}}.
\]
Assim,
\[
r\approx\SI{29.2}{\centi\meter}.
\]

(c) \textbf{Distância até a placa detectora.}
A trajetória mostrada corresponde a uma semicircunferência. Dessa forma, a
distância vertical entre o ponto de entrada na região magnética e o ponto de
impacto na placa detectora corresponde ao diâmetro da trajetória:
\[
d=2r.
\]
Logo,
\[
d=2(0.292)
\approx
\boxed{\SI{0.585}{\meter}}
\]
ou, equivalentemente,
\[
\boxed{d\approx\SI{58.5}{\centi\meter}}.
\]

\textbf{Observação:} o seletor de velocidades utiliza a compensação entre as
forças elétrica e magnética para selecionar partículas com uma velocidade
específica. Na segunda região, a presença apenas do campo magnético permite
determinar a razão entre momento linear e carga da partícula a partir da
curvatura de sua trajetória.
}
\end{questao}



% ============================================================
% QUESTÃO 5 — PARTÍCULA ACELERADA E ANALISADOR MAGNÉTICO
% ============================================================

\begin{questao}[4]{}

Um elétron parte praticamente do repouso e é acelerado por uma diferença de
potencial de
\[
\Delta V=\SI{1.8}{\kilo\volt}.
\]
Após atravessar uma pequena abertura, ele penetra perpendicularmente em uma
região onde existe um campo magnético uniforme de intensidade
\[
B=\SI{25}{\milli\tesla},
\]
perpendicular ao plano da figura.

Considere
\[
m_e=\SI{9.11e-31}{\kilogram},
\qquad
|q_e|=\SI{1.60e-19}{\coulomb},
\]
e despreze efeitos relativísticos.

Determine:

\begin{itensq}
    \item a velocidade do elétron ao deixar a região de aceleração;
    \item o raio da trajetória circular no campo magnético;
    \item o período do movimento circular;
    \item o tempo necessário para o elétron percorrer um quarto de circunferência.
\end{itensq}

\begin{center}
\begin{tikzpicture}[
    >=Latex,
    line cap=round,
    line join=round,
    scale=0.95,
    font=\sffamily
]

% --------------------------------------------------
% Região de aceleração
% --------------------------------------------------

\node at (-4.6,0) {\Large $e^-$};

\draw[thick,-{Latex[length=3mm]}]
    (-4.25,0) -- (-3.65,0)
    node[midway,above] {$\vec v$};

% Placas aceleradoras
\draw[line width=2.5pt] (-3.3,-1.1) -- (-3.3,1.1);
\draw[line width=2.5pt] (-1.3,-1.1) -- (-1.3,1.1);

\node[font=\Large] at (-3.65,0.65) {$-$};
\node[font=\Large] at (-0.95,0.65) {$+$};

% Campo elétrico
\foreach \y in {-0.7,-0.35,0,0.35,0.7}{
    \draw[myE,thick,-{Latex[length=2.5mm]}]
        (-1.55,\y) -- (-3.05,\y);
}

\node[myE] at (-2.3,1.45) {$\Delta V=\SI{1.8}{\kilo\volt}$};

% Fenda
\draw[line width=3pt] (-0.65,0.28) -- (-0.65,1.0);
\draw[line width=3pt] (-0.65,-0.28) -- (-0.65,-1.0);

% Feixe após aceleração
\draw[thick,dashed,-{Latex[length=3mm]}]
    (-0.6,0) -- (0.6,0);

% --------------------------------------------------
% Região magnética
% --------------------------------------------------

\foreach \x in {0.8,1.5,2.2,2.9,3.6}{
    \foreach \y in {-2.1,-1.4,-0.7,0,0.7,1.4,2.1}{
        \draw[myB,thick] (\x-0.08,\y-0.08) -- (\x+0.08,\y+0.08);
        \draw[myB,thick] (\x-0.08,\y+0.08) -- (\x+0.08,\y-0.08);
    }
}

\node[myB] at (3.65,2.55)
    {$\vec B$ entrando no plano};

% Trajetória
\draw[very thick,dashed,-{Latex[length=3mm]}]
    (0.6,0)
    arc[start angle=90,end angle=0,radius=2.0];

% Centro
\fill (2.6,-2.0) circle (1.5pt);

\draw[thin,dashed]
    (2.6,-2.0) -- (0.6,0);

\node at (1.55,-1.05) {$r$};

\end{tikzpicture}
\end{center}

\resposta{$v\approx\SI{2.52e7}{\meter\per\second}$,
\quad
$r\approx\SI{5.72}{\milli\meter}$,
\quad
$T\approx\SI{1.43}{\nano\second}$,
\quad
$t_{1/4}\approx\SI{0.357}{\nano\second}$}

\resolucao{%
(a) \textbf{Velocidade após a aceleração.}
A energia potencial elétrica adquirida pelo elétron transforma-se em energia
cinética:
\[
|q_e|\Delta V=\frac{1}{2}m_ev^2.
\]
Portanto,
\[
v=\sqrt{\frac{2|q_e|\Delta V}{m_e}}.
\]
Substituindo os valores:
\[
v=
\sqrt{
\frac{
2(1.60\times10^{-19})(1.8\times10^3)
}{
9.11\times10^{-31}
}}
\approx
\boxed{\SI{2.52e7}{\meter\per\second}}.
\]

(b) \textbf{Raio da trajetória.}
Como $\vec v\perp\vec B$, a força magnética atua como força centrípeta:
\[
|q_e|vB=\frac{m_ev^2}{r}.
\]
Assim,
\[
r=\frac{m_ev}{|q_e|B}.
\]
Logo,
\[
r=
\frac{(9.11\times10^{-31})(2.52\times10^7)}
{(1.60\times10^{-19})(25\times10^{-3})}
\approx
\boxed{\SI{5.72}{\milli\meter}}.
\]

(c) \textbf{Período do movimento.}
\[
T=\frac{2\pi m_e}{|q_e|B}.
\]
Portanto,
\[
T\approx
\boxed{\SI{1.43e-9}{\second}}
=
\boxed{\SI{1.43}{\nano\second}}.
\]

(d) Para um quarto de circunferência:
\[
t_{1/4}=\frac{T}{4}.
\]
Assim,
\[
t_{1/4}\approx
\boxed{\SI{3.57e-10}{\second}}
=
\boxed{\SI{0.357}{\nano\second}}.
\]

\textbf{Interpretação:} a diferença de potencial determina a energia cinética
da partícula, enquanto o campo magnético modifica apenas a direção da
velocidade, produzindo uma trajetória circular.
}

\end{questao}


% ============================================================
% QUESTÃO 6 — ESPECTRÔMETRO DE MASSAS COM DOIS ISÓTOPOS
% ============================================================

\begin{questao}[4]{}

Um espectrômetro de massas é utilizado para separar dois isótopos de neônio,
$^{20}\mathrm{Ne}^{+}$ e $^{22}\mathrm{Ne}^{+}$. Inicialmente, os íons
atravessam um seletor de velocidades constituído por campos elétrico e
magnético perpendiculares, de intensidades
\[
E=\SI{30}{\kilo\volt\per\meter},
\qquad
B_s=\SI{20}{\milli\tesla}.
\]

Os íons selecionados entram em seguida em uma região contendo apenas um campo
magnético uniforme de intensidade
\[
B_a=\SI{0.35}{\tesla}.
\]

Considere
\[
u=\SI{1.66e-27}{\kilogram},
\qquad
q=+e=\SI{1.60e-19}{\coulomb}.
\]

Determine:

\begin{itensq}
    \item a velocidade dos íons que atravessam o seletor sem desvio;
    \item os raios das trajetórias dos dois isótopos;
    \item a distância entre os pontos de impacto na placa detectora;
    \item qual isótopo atinge a região mais distante do ponto de entrada.
\end{itensq}

\begin{center}
\begin{tikzpicture}[
    >=Latex,
    line cap=round,
    line join=round,
    scale=0.88,
    font=\sffamily
]

% Fonte de íons
\draw[thick] (-4.7,0) circle (0.3);
\node at (-4.7,0) {$+$};

\draw[thick,-{Latex[length=3mm]}]
    (-4.35,0) -- (-3.65,0);

% Seletor
\draw[line width=2pt] (-3.2,0.8) -- (-1.0,0.8);
\draw[line width=2pt] (-3.2,-0.8) -- (-1.0,-0.8);

\node at (-2.85,1.12) {$-$};
\node at (-2.85,-1.12) {$+$};

\foreach \x in {-2.9,-2.4,-1.9,-1.4}{
    \draw[myE,thick,-{Latex[length=2.5mm]}]
        (\x,-0.6) -- (\x,0.6);
}

\foreach \x in {-2.7,-2.0,-1.3}{
    \foreach \y in {-0.4,0,0.4}{
        \fill[myB] (\x,\y) circle (2pt);
    }
}

\node at (-2.1,1.55) {Seletor de velocidades};

% Fenda
\draw[line width=3pt] (-0.45,0.3) -- (-0.45,1);
\draw[line width=3pt] (-0.45,-0.3) -- (-0.45,-1);

% Região analisadora
\foreach \x in {0.2,0.9,1.6,2.3,3.0,3.7}{
    \foreach \y in {-4.1,-3.4,-2.7,-2.0,-1.3,-0.6,0.1}{
        \fill[myB] (\x,\y) circle (2.3pt);
    }
}

\node[myB] at (3.5,0.7) {$\vec B_a$};

% Trajetória 20
\draw[very thick,dashed]
    (-0.35,0)
    arc[start angle=90,end angle=-90,radius=1.8];

\node at (1.25,-1.1) {$^{20}\mathrm{Ne}^{+}$};

% Trajetória 22
\draw[very thick,dash dot]
    (-0.35,0)
    arc[start angle=90,end angle=-90,radius=2.15];

\node at (2.65,-2.2) {$^{22}\mathrm{Ne}^{+}$};

% Detector
\draw[line width=4pt] (-0.45,-3.2) -- (-0.45,-4.8);

\node[rotate=90] at (-0.85,-4.0) {Detector};

% Marcas de impacto
\fill (-0.35,-3.6) circle (2.2pt);
\fill (-0.35,-4.3) circle (2.2pt);

\draw[<->,thick]
    (-1.25,-3.6) -- (-1.25,-4.3)
    node[midway,left] {$\Delta d$};

\end{tikzpicture}
\end{center}

\resposta{$v=\SI{1.50e6}{\meter\per\second}$,
\quad
$r_{20}\approx\SI{0.889}{\meter}$,
\quad
$r_{22}\approx\SI{0.977}{\meter}$,
\quad
$\Delta d\approx\SI{0.178}{\meter}$}

\resolucao{%
(a) \textbf{Velocidade selecionada.}
No seletor:
\[
qE=qvB_s.
\]
Logo,
\[
v=\frac{E}{B_s}.
\]
Assim,
\[
v=
\frac{30\times10^3}{20\times10^{-3}}
=
\boxed{\SI{1.50e6}{\meter\per\second}}.
\]

(b) \textbf{Raios das trajetórias.}
Para um íon de massa $m$:
\[
r=\frac{mv}{qB_a}.
\]

Para o $^{20}\mathrm{Ne}^{+}$:
\[
m_{20}=20u=20(1.66\times10^{-27}),
\]
resultando em
\[
r_{20}\approx
\boxed{\SI{0.889}{\meter}}.
\]

Para o $^{22}\mathrm{Ne}^{+}$:
\[
m_{22}=22u,
\]
e
\[
r_{22}\approx
\boxed{\SI{0.977}{\meter}}.
\]

(c) Como cada partícula descreve uma semicircunferência, a posição de impacto
é determinada pelo diâmetro:
\[
d=2r.
\]
A separação entre os impactos é
\[
\Delta d=2(r_{22}-r_{20}).
\]
Logo,
\[
\Delta d
=
2(0.977-0.889)
\approx
\boxed{\SI{0.178}{\meter}}.
\]

(d) Como
\[
r=\frac{mv}{qB},
\]
mantidos $v$, $q$ e $B$, o raio é diretamente proporcional à massa. Portanto,
o isótopo
\[
\boxed{^{22}\mathrm{Ne}^{+}}
\]
atinge a região mais distante.

\textbf{Interpretação:} o espectrômetro converte diferenças de massa em
diferenças geométricas mensuráveis na posição de impacto das partículas.
}

\end{questao}


% ============================================================
% QUESTÃO 7 — TRAJETÓRIA HELICOIDAL
% ============================================================

\begin{questao}[4]{}

Uma partícula alfa entra em uma região onde existe um campo magnético uniforme
de intensidade
\[
B=\SI{0.18}{\tesla}.
\]
A velocidade inicial da partícula possui módulo
\[
v=\SI{7.0e5}{\meter\per\second}
\]
e forma um ângulo de
\[
\theta=40^\circ
\]
com a direção do campo magnético.

Considere
\[
m_\alpha=\SI{6.64e-27}{\kilogram},
\qquad
q_\alpha=\SI{3.20e-19}{\coulomb}.
\]

Determine:

\begin{itensq}
    \item as componentes $v_{\parallel}$ e $v_{\perp}$ da velocidade;
    \item o raio da trajetória helicoidal;
    \item o período de uma volta;
    \item o passo $p$ da hélice;
    \item o que aconteceria com a trajetória se $\theta=90^\circ$.
\end{itensq}

\begin{center}
\begin{tikzpicture}[
    >=Latex,
    line cap=round,
    line join=round,
    scale=1.0,
    font=\sffamily
]

% Eixo do campo
\draw[myB,line width=1.3pt,-{Latex[length=3mm]}]
    (-4,0) -- (4.7,0)
    node[right] {$\vec B$};

% Linhas paralelas de campo
\foreach \y in {-1.4,-0.7,0.7,1.4}{
    \draw[myB!70,thin,-{Latex[length=2mm]}]
        (-3.8,\y) -- (4.3,\y);
}

% Hélice estilizada
\draw[very thick]
plot[
    domain=0:720,
    samples=150,
    smooth
]
({-3.3+0.010*\x},{0.60*sin(\x)});

% Partícula
\fill (-3.3,0) circle (2.8pt);
\node[above left] at (-3.3,0) {$\alpha$};

% Vetor velocidade
\draw[thick,-{Latex[length=3mm]}]
    (-3.3,0) -- (-2.35,0.8)
    node[above] {$\vec v$};

% Componente paralela
\draw[dashed,thick,-{Latex[length=2.5mm]}]
    (-3.3,-0.15) -- (-2.0,-0.15)
    node[midway,below] {$v_{\parallel}$};

% Componente perpendicular
\draw[dashed,thick,-{Latex[length=2.5mm]}]
    (-3.3,0) -- (-3.3,0.85)
    node[left] {$v_{\perp}$};

% Ângulo
\draw (-2.75,0)
    arc[start angle=0,end angle=40,radius=0.55];

\node at (-2.55,0.25) {$\theta$};

% Passo
\draw[<->,thick]
    (0.3,-1.15) -- (3.9,-1.15)
    node[midway,fill=white,inner sep=2pt] {$p$};

\node at (0,-2.0)
    {Trajetória helicoidal};

\end{tikzpicture}
\end{center}

\resposta{$v_{\parallel}\approx\SI{5.36e5}{\meter\per\second}$,
\quad
$v_{\perp}\approx\SI{4.50e5}{\meter\per\second}$,
\quad
$r\approx\SI{5.18}{\centi\meter}$,
\quad
$T\approx\SI{0.724}{\micro\second}$,
\quad
$p\approx\SI{0.388}{\meter}$}

\resolucao{%
(a) \textbf{Decomposição da velocidade.}
A componente paralela ao campo é
\[
v_{\parallel}=v\cos\theta,
\]
e a componente perpendicular:
\[
v_{\perp}=v\sin\theta.
\]

Logo,
\[
v_{\parallel}
=
(7.0\times10^5)\cos40^\circ
\approx
\boxed{\SI{5.36e5}{\meter\per\second}},
\]
e
\[
v_{\perp}
=
(7.0\times10^5)\sin40^\circ
\approx
\boxed{\SI{4.50e5}{\meter\per\second}}.
\]

(b) \textbf{Raio da hélice.}
Somente $v_{\perp}$ produz movimento circular:
\[
q_\alpha v_{\perp}B
=
\frac{m_\alpha v_{\perp}^2}{r}.
\]
Portanto,
\[
r=
\frac{m_\alpha v_{\perp}}{q_\alpha B}.
\]
Substituindo:
\[
r\approx
\boxed{\SI{5.18}{\centi\meter}}.
\]

(c) \textbf{Período.}
\[
T=
\frac{2\pi m_\alpha}{q_\alpha B}.
\]
Assim,
\[
T\approx
\boxed{\SI{7.24e-7}{\second}}
=
\boxed{\SI{0.724}{\micro\second}}.
\]

(d) \textbf{Passo da hélice.}
Durante uma volta, o deslocamento paralelo ao campo é
\[
p=v_{\parallel}T.
\]
Logo,
\[
p=
(5.36\times10^5)(7.24\times10^{-7})
\approx
\boxed{\SI{0.388}{\meter}}.
\]

(e) Para $\theta=90^\circ$:
\[
v_{\parallel}=0.
\]
Consequentemente, desaparece o avanço longitudinal da hélice e a trajetória
passa a ser
\[
\boxed{\text{circular}}.
\]

\textbf{Interpretação:} o movimento helicoidal resulta da superposição de um
movimento retilíneo uniforme paralelo ao campo com um movimento circular
uniforme perpendicular ao campo.
}

\end{questao}


% ============================================================
% QUESTÃO 8 — CÍCLOTRON
% ============================================================

\begin{questao}[4]{}

Um cíclotron é utilizado para acelerar prótons em uma região de campo magnético
uniforme de intensidade
\[
B=\SI{0.80}{\tesla}.
\]
Os prótons descrevem órbitas de raio progressivamente maior até atingirem o
raio máximo
\[
R=\SI{25}{\centi\meter}.
\]

Considere
\[
m_p=\SI{1.67e-27}{\kilogram},
\qquad
q_p=\SI{1.60e-19}{\coulomb},
\]
e despreze correções relativísticas.

Determine:

\begin{itensq}
    \item a frequência que deve ser aplicada entre os eletrodos do cíclotron;
    \item o período do movimento dos prótons;
    \item a velocidade do próton ao atingir o raio máximo;
    \item a energia cinética máxima do próton, em joules e em elétron-volts.
\end{itensq}

\begin{center}
\begin{tikzpicture}[
    >=Latex,
    line cap=round,
    line join=round,
    scale=0.95,
    font=\sffamily
]

% Região magnética
\foreach \x in {-3,-2.25,-1.5,-0.75,0.75,1.5,2.25,3}{
    \foreach \y in {-2.2,-1.5,-0.8,0,0.8,1.5,2.2}{
        \fill[myB!75] (\x,\y) circle (1.8pt);
    }
}

% Dees
\draw[line width=2pt]
    (0.18,2.55)
    arc[start angle=90,end angle=270,radius=2.55];

\draw[line width=2pt]
    (-0.18,-2.55)
    arc[start angle=-90,end angle=90,radius=2.55];

% Fenda central
\draw[line width=2pt] (-0.18,-2.55) -- (-0.18,2.55);
\draw[line width=2pt] (0.18,-2.55) -- (0.18,2.55);

\node at (-1.45,0) {\Large D};
\node at (1.45,0) {\Large D};

% Campo elétrico alternado na fenda
\draw[myE,very thick,<->]
    (-0.13,0.75) -- (0.13,0.75);

\node[myE,right] at (0.25,0.75) {$\vec E(t)$};

% Espiral
\draw[very thick,-{Latex[length=3mm]}]
plot[
    domain=0:900,
    samples=250,
    smooth
]
({0.0027*\x*cos(\x)},
 {0.0027*\x*sin(\x)});

% Centro
\fill (0,0) circle (2.5pt);

% Raio final
\draw[dashed,thick]
    (0,0) -- (2.43,0);

\node[below] at (1.25,0) {$R$};

% Campo
\node[myB] at (3.7,2.0)
    {$\vec B$};

% Saída
\draw[thick,-{Latex[length=3mm]}]
    (2.45,-0.35) -- (3.65,-0.7);

\node[right] at (3.65,-0.7)
    {feixe};

\end{tikzpicture}
\end{center}

\resposta{$f\approx\SI{12.2}{\mega\hertz}$,
\quad
$T\approx\SI{82.0}{\nano\second}$,
\quad
$v_{\max}\approx\SI{1.92e7}{\meter\per\second}$,
\quad
$K_{\max}\approx\SI{3.07e-13}{\joule}
\approx\SI{1.92}{\mega\electronvolt}$}

\resolucao{%
(a) \textbf{Frequência ciclotrônica.}
A força magnética fornece a força centrípeta:
\[
q_pvB=\frac{m_pv^2}{r}.
\]
Como
\[
v=\omega r,
\]
obtém-se
\[
\omega=\frac{q_pB}{m_p}.
\]
Portanto,
\[
f=\frac{\omega}{2\pi}
=
\frac{q_pB}{2\pi m_p}.
\]
Substituindo:
\[
f=
\frac{
(1.60\times10^{-19})(0.80)
}{
2\pi(1.67\times10^{-27})
}
\approx
\boxed{\SI{12.2}{\mega\hertz}}.
\]

(b) \textbf{Período.}
\[
T=\frac{1}{f}.
\]
Logo,
\[
T\approx
\boxed{\SI{82.0}{\nano\second}}.
\]

(c) \textbf{Velocidade máxima.}
Da condição de movimento circular:
\[
r=\frac{m_pv}{q_pB}.
\]
Para $r=R$:
\[
v_{\max}=
\frac{q_pBR}{m_p}.
\]
Assim,
\[
v_{\max}
=
\frac{
(1.60\times10^{-19})(0.80)(0.25)
}{
1.67\times10^{-27}
}
\approx
\boxed{\SI{1.92e7}{\meter\per\second}}.
\]

(d) \textbf{Energia cinética máxima.}
\[
K_{\max}
=
\frac{1}{2}m_pv_{\max}^2.
\]
Substituindo:
\[
K_{\max}
\approx
\boxed{\SI{3.07e-13}{\joule}}.
\]

Como
\[
\SI{1}{\electronvolt}
=
\SI{1.60e-19}{\joule},
\]
tem-se
\[
K_{\max}
=
\frac{3.07\times10^{-13}}
{1.60\times10^{-19}}
\approx
1.92\times10^6\,\mathrm{eV}.
\]
Portanto,
\[
\boxed{K_{\max}\approx\SI{1.92}{\mega\electronvolt}}.
\]

\textbf{Interpretação:} a frequência ciclotrônica permanece aproximadamente
constante enquanto os efeitos relativísticos forem desprezíveis. A cada
passagem pela região entre os eletrodos, o campo elétrico fornece energia à
partícula, aumentando sua velocidade e, consequentemente, o raio da órbita.
}

\end{questao}

% END BANCO COMPLEMENTAR

\gabarito[12.5 Lei de Lenz e indução eletromagnética]
